\chapter{Chiral modules and geometric resolutions}
\label{chap:chiral-modules}

\section{Factorization modules: definitions and categories}

A chiral algebra without its representations is like a group without
its actions: the abstract structure exists, but the richest
mathematics is inaccessible.  Conformal blocks, fusion rules, the
Kazhdan--Lusztig equivalence, and the Verlinde formula all live
at the module level.  The bar-cobar adjunction of
Chapters~\ref{chap:bar-cobar}--\ref{chap:higher-genus} must
therefore extend from algebras to modules.

The extension is not formal: the bar complex of a module
$\bar{B}(\cA, M)$ lives on \emph{pointed} configuration spaces
(one distinguished point carrying the module, the rest carrying
algebra data), and its Verdier dual is a comodule over
$\bar{B}(\cA)$.  The Koszul property of the pair $(\cA, \cA^!)$
descends to modules: if $\cA$ is Koszul, then every module $M$
admits a bar resolution whose $E_2$-collapse computes Ext groups.
The resulting module Koszul duality exchanges categories of
representations under the level involution
$k \leftrightarrow -k - 2h^\vee$ --- the representation-theoretic
face of the same Verdier duality that governs the algebra-level
theory.

We begin by establishing the precise categorical framework for modules
over chiral and factorization algebras.  The key idea---originating in the
``Swiss-cheese'' picture of Voronov~\cite{Voronov99} and developed
systematically by Costello--Gwilliam~\cite{CG17} and
Ayala--Francis~\cite{AF19}---is that a module is a factorization algebra
on a pointed space, where the marked point carries the module data and the
remaining points carry algebra data.

\begin{definition}[Factorization module]\label{def:factorization-module}
\index{factorization module|textbf}
Let $\mathcal{F}$ be a factorization algebra on a manifold $M$ with
values in a symmetric monoidal $\infty$-category $\mathcal{V}$, and let
$x \in M$ be a point.  A \emph{factorization module} over $\mathcal{F}$
supported at $x$ is a rule $\mathcal{M}$ that assigns:
\begin{enumerate}
\item To each open disk $D \ni x$: an object
   $\mathcal{M}(D) \in \mathcal{V}$.
\item To each finite collection of pairwise disjoint opens
   $U_1, \ldots, U_n \subset D$ with $U_i \cap \{x\} = \emptyset$
   for all~$i$ and $x \in D$: an \emph{action map}
   \[
   \alpha_{U_1,\ldots,U_n}^D:
   \mathcal{F}(U_1) \otimes \cdots \otimes \mathcal{F}(U_n)
   \otimes \mathcal{M}(D)
   \;\longrightarrow\; \mathcal{M}(D),
   \]
   compatible with the factorization structure of $\mathcal{F}$ on
   $M \setminus \{x\}$.
\end{enumerate}
These data satisfy associativity and unit axioms parallel to
(FA2)--(FA5) of Definition~\ref{def:factorization-algebra-AF},
with $\mathcal{M}$ replacing $\mathcal{F}$ at the marked point~$x$.
\end{definition}

\begin{remark}[The Swiss-cheese picture]\label{rem:swiss-cheese}
When $M = \mathbb{R}^n$ and $x = 0$, the collection of disks around the
origin with disjoint smaller disks away from the origin defines the
\emph{Swiss-cheese operad} $\mathrm{SC}_n$ of
Voronov~\cite{Voronov99}.  Algebras over $\mathrm{SC}_n$ are pairs
$(A, M)$ where $A$ is an $E_n$-algebra and $M$ is an $A$-module in
an appropriate homotopical sense.  Our
Definition~\ref{def:factorization-module} globalizes this picture
to arbitrary manifolds.
\end{remark}

\begin{definition}[Chiral module]\label{def:chiral-module}
\index{chiral module|textbf}
Let $\mathcal{A}$ be a chiral algebra (= $\Einf$-chiral algebra,
Definition~\ref{def:einf-chiral}) on a smooth curve $X$.  A
\emph{chiral $\mathcal{A}$-module} supported at $x \in X$ is a
$\mathcal{D}_X$-module $\mathcal{M}$ together with a chiral action map
\[
\alpha: j_* j^*(\mathcal{A} \boxtimes \mathcal{M})
\;\longrightarrow\; \Delta_!\, \mathcal{M}
\]
(where $j: X^2 \setminus \Delta \hookrightarrow X^2$ is the complement
of the diagonal) satisfying:
\begin{enumerate}[label=(\roman*)]
\item \emph{Associativity}: the two compositions
   $\mathcal{A} \boxtimes \mathcal{A} \boxtimes \mathcal{M}
   \rightrightarrows \Delta^{(3)}_!\,\mathcal{M}$ agree;
\item \emph{Unit}: the composition
   $\omega_X \boxtimes \mathcal{M} \xrightarrow{\eta \boxtimes
   \mathrm{id}} \mathcal{A} \boxtimes \mathcal{M}
   \xrightarrow{\alpha} \Delta_!\,\mathcal{M}$
   equals the canonical map.
\end{enumerate}
This is the $\mathcal{D}$-module counterpart of a module over a vertex
algebra~\cite[\S3.6]{BD04}.
\end{definition}

\begin{definition}[$\Eone$-chiral module]\label{def:e1-chiral-module}
\index{$\mathbb{E}_1$-module|textbf}
Let $\mathcal{A}$ be an $\Eone$-chiral algebra
(Definition~\ref{def:e1-chiral}) on $X$.  An \emph{$\Eone$-chiral
$\mathcal{A}$-module} is a $\mathcal{D}_X$-module $\mathcal{M}$
equipped with a chiral action
\[
\alpha: j_* j^*(\mathcal{A} \boxtimes \mathcal{M})
\;\longrightarrow\; \Delta_!\, \mathcal{M}
\]
satisfying associativity and unitality as in
Definition~\ref{def:chiral-module}, but indexed by \emph{ordered}
configurations: the action of $n$ copies of~$\mathcal{A}$ on
$\mathcal{M}$ is parameterized by the ordered configuration space
$\widetilde{C}_{n,1}(X) := \{(z_1, \ldots, z_n; w) \in X^{n+1}
\mid z_i \neq z_j,\; z_i \neq w\}$
with no $\Sigma_n$-equivariance imposed.

In vertex algebra language, $\mathcal{M}$ is a module over a nonlocal
vertex algebra: it carries a state-field map
$Y_{\mathcal{M}}: \mathcal{A} \to \mathrm{End}(\mathcal{M})[[z, z^{-1}]]$
satisfying weak associativity (the module Borcherds identity) but not
the commutativity condition
$Y_{\mathcal{M}}(a,z) Y_{\mathcal{M}}(b,w) \sim
Y_{\mathcal{M}}(b,w) Y_{\mathcal{M}}(a,z)$.
\end{definition}

\begin{definition}[Module categories]\label{def:module-categories}
\index{module category!equivalence}
For a chiral algebra $\mathcal{A}$ on $X$, we define the following
categories of modules:
\begin{enumerate}[label=(\roman*)]
\item $\mathrm{Mod}_{\mathcal{A}}$: the dg category of chiral
   $\mathcal{A}$-modules (Definition~\ref{def:chiral-module}),
   with morphisms the $\mathcal{D}_X$-module maps compatible with
   the chiral action.  This is a stable $\infty$-category when
   $\mathcal{V} = \mathrm{Ch}(\mathbb{C})$.
\item $\mathrm{Mod}^{\Eone}_{\mathcal{A}}$: the dg category of
   $\Eone$-chiral $\mathcal{A}$-modules
   (Definition~\ref{def:e1-chiral-module}), when $\mathcal{A}$ is
   $\Eone$-chiral.  There is a forgetful functor
   $\mathrm{Mod}_{\mathcal{A}} \to
   \mathrm{Mod}^{\Eone}_{\mathcal{A}}$ when $\mathcal{A}$ is
   $\Einf$-chiral.
\item $D(\mathrm{Mod}_{\mathcal{A}})$: the derived category,
   obtained by inverting quasi-isomorphisms.  We assume throughout
   that $\mathrm{Mod}_{\mathcal{A}}$ is \emph{compactly generated}:
   i.e., it is generated under filtered colimits by the compact
   objects $\mathrm{Mod}_{\mathcal{A}}^c$.  This holds whenever
   $\mathcal{A}$ is finitely generated as a chiral algebra
   (the free modules
   $\mathrm{Free}_{\mathcal{A}}(V)$ for finite-dimensional $V$
   are compact generators).
\end{enumerate}
\end{definition}

\begin{remark}[Tensor structure on module categories]
\label{rem:tensor-structure-modules}
\index{vertex tensor category}
For a rational vertex algebra~$V$ ($C_2$-cofinite with finitely many
simple modules; Theorem~\ref{thm:zhu-correspondence}(iii)), the module
category $\mathrm{Mod}_V$ carries a braided tensor structure: the
\emph{fusion product} $\mathcal{M}_1 \boxtimes_V \mathcal{M}_2$,
constructed via analytic continuation of intertwining operators
(Huang--Lepowsky--Zhang~\cite{HLZ}).  The braiding is determined by the monodromy
of conformal blocks on~$\mathbb{P}^1$.

For the affine algebra $\widehat{\mathfrak{g}}_k$ at positive integer
level, this tensor structure on
$\mathcal{O}_k^{\mathrm{int}}(\widehat{\mathfrak{g}})$ is identified
with the quantum group tensor product on
$\mathrm{Rep}(U_q(\mathfrak{g}))$ via the Kazhdan--Lusztig equivalence
(Theorem~\ref{thm:km-quantum-groups}).

The construction of fusion products uses the factorization structure
on configuration spaces: the fusion of modules at points $z_1, z_2$ is
computed by integrating over the configuration space $C_2(X)$ and
taking the limit $z_2 \to z_1$.  This is one reason why module
categories for chiral algebras are more structured than for general
associative algebras.
\end{remark}

\begin{proposition}[Fusion product of Heisenberg Fock modules; \ClaimStatusProvedHere]
\label{prop:fock-fusion-product}
\index{Fock module!fusion product}
\index{fusion product!Heisenberg}
The fusion product of Fock modules over the Heisenberg algebra
$\mathcal{H}_\kappa$ is:
\[
F_\lambda \boxtimes_{\mathcal{H}_\kappa} F_\mu
\;\cong\; F_{\lambda + \mu}.
\]
The braiding isomorphism
$\sigma: F_\lambda \boxtimes F_\mu \xrightarrow{\sim}
F_\mu \boxtimes F_\lambda$
is multiplication by $e^{2\pi i \lambda\mu/\kappa}$
\textup{(}the monodromy of the intertwining operator around
$z = 0$\textup{)}.
\end{proposition}

\begin{proof}
The intertwining operator $\mathcal{Y}\colon F_\lambda \otimes
F_\mu \to F_{\lambda+\mu}((z))$ is given by
\[
\mathcal{Y}(|\lambda\rangle, z)
= z^{\lambda\mu/\kappa}\,
\exp\!\Bigl(\lambda \sum_{n > 0} \frac{J_{-n}}{n}\, z^n\Bigr)\,
\exp\!\Bigl(-\lambda \sum_{n > 0} \frac{J_n}{n}\, z^{-n}\Bigr),
\]
where $|\lambda\rangle$ is the highest weight vector of $F_\lambda$
(Definition~\ref{def:fock-module}).  The fusion product is the image of
$\mathcal{Y}$ under analytic continuation (in the sense
of~\cite{HLZ}), which is isomorphic to $F_{\lambda+\mu}$: the
charge is additive ($J_0(|\lambda\rangle \boxtimes |\mu\rangle)
= (\lambda + \mu)(|\lambda\rangle \boxtimes |\mu\rangle)$) and
the Fock space structure is generated by $J_{-n}$ for $n \geq 1$.

The braiding comes from the monodromy of the factor
$z^{\lambda\mu/\kappa}$: analytically continuing $z \mapsto
e^{2\pi i} z$ gives the phase $e^{2\pi i \lambda\mu/\kappa}$.
Since $\mathcal{H}_\kappa$ is abelian, this determines the braiding
completely.

In the bar-cobar framework, the fusion product is realized on the
configuration space $C_2(X)$: two Fock modules at points $z_1, z_2$
are fused by integrating the intertwining operator over a small circle
$|z_1 - z_2| = \varepsilon$ and taking $\varepsilon \to 0$.  The
charge additivity follows from the residue calculus on the boundary
stratum $D_{12} = \{z_1 = z_2\} \subset \overline{C}_2(X)$.
\end{proof}

\begin{remark}[Fusion under Koszul duality]
\label{rem:fusion-koszul}
Under the module Koszul duality of
Corollary~\ref{cor:heisenberg-module-equivalence},
$F_\lambda^! \boxtimes_{\mathcal{H}^!} F_\mu^! \cong
F_{\lambda+\mu}^!$: the fusion product is preserved.  The braiding
phase $e^{2\pi i\lambda\mu/\kappa}$ transforms to
$e^{-2\pi i\lambda\mu/\kappa}$ (the sign change reflects the
curvature reversal $\kappa \mapsto -\kappa$ on the Koszul dual side).
This is consistent with the general principle that Koszul duality
exchanges a braided tensor category with its \emph{reverse} braiding.
\end{remark}

\begin{remark}[Koszul duality for module categories]
\label{rem:koszul-module-categories}
For a Koszul dual pair $(\mathcal{A}, \mathcal{A}^!)$ of chiral algebras,
the bar-cobar adjunction induces an equivalence
\[
\mathrm{Mod}^{\mathrm{compl}}_{\mathcal{A}}
\;\simeq\;
\mathrm{CoMod}^{\mathrm{conil}}_{\mathcal{A}^!}
\]
between \emph{completed} $\mathcal{A}$-modules and \emph{conilpotent}
$\mathcal{A}^!$-comodules (see also
Appendix~\ref{app:existence-criteria}).  When $\mathcal{A}^!$ is
finite-type graded, the coalgebra--algebra adjunction provides a further
equivalence $\mathrm{CoMod}^{\mathrm{conil}}_{\mathcal{A}^!} \simeq
\mathrm{Mod}^{\mathrm{compl}}_{\mathcal{A}^!}$, so one recovers
$\mathrm{Mod}^{\mathrm{compl}}_{\mathcal{A}} \simeq
\mathrm{Mod}^{\mathrm{compl}}_{\mathcal{A}^!}$.
Without the finite-type hypothesis, the module--module equivalence fails
in general: the intermediate step through comodules is essential.

The bar resolution of a module
$\mathcal{M} \in \mathrm{Mod}_{\mathcal{A}}$ produces a comodule
over $\bar{B}(\mathcal{A}) \simeq \mathcal{A}^{!,c}$ (the Koszul dual
coalgebra), and the cobar construction inverts this on conilpotent
objects.  This equivalence is made explicit in the character computations
of \S\ref{sec:complete-calculations} below.
\end{remark}

\begin{theorem}[Monoidal module Koszul duality; \ClaimStatusProvedHere]
\label{thm:monoidal-module-koszul}
\index{Koszul duality!monoidal}
\index{fusion product!Koszul duality}
Let $(\mathcal{A}, \mathcal{A}^!)$ be a Koszul dual pair of chiral
algebras with $\mathcal{A}$ rational ($C_2$-cofinite with finitely
many simple modules).  Then the module Koszul duality equivalence
\[
\Phi\colon
\mathrm{Mod}^{\mathrm{compl}}_{\mathcal{A}}
\;\xrightarrow{\;\sim\;}
\mathrm{CoMod}^{\mathrm{conil}}_{\mathcal{A}^!}
\]
of Remark~\textup{\ref{rem:koszul-module-categories}} is
\emph{lax monoidal} with respect to the fusion products on both sides:
there is a natural transformation
\begin{equation}\label{eq:monoidal-bar}
\mu_{M,N}\colon
\Phi(M \boxtimes_{\mathcal{A}} N)
\;\longrightarrow\;
\Phi(M) \boxtimes_{\mathcal{A}^{!,c}} \Phi(N)
\end{equation}
and $\mu_{M,N}$ is a quasi-isomorphism when $M$ and $N$ are
Koszul modules \textup{(}i.e., the bar spectral sequence for each
degenerates at~$E_2$\textup{)}.  When $\mathcal{A}^!$ is
finite-type graded, the composite equivalence
$\mathrm{Mod}^{\mathrm{compl}}_{\mathcal{A}} \simeq
\mathrm{Mod}^{\mathrm{compl}}_{\mathcal{A}^!}$
is monoidal with reversed braiding: $\sigma^! = \sigma^{-1}$.
\end{theorem}

\begin{proof}
\emph{Step~1: Construction of $\mu_{M,N}$.}
The bar construction on the fusion product
$M \boxtimes_{\mathcal{A}} N$ is computed on the configuration space
$C_n(X)$ with two marked points.  By the factorization property of
the bar complex (Theorem~\ref{thm:geometric-bar-definition}), the
summand of $\bar{B}(M \boxtimes N)$ where $p$ points cluster
near the first marked point and $q$ points near the second yields
$\bar{B}_p(M) \otimes \bar{B}_q(N)$ on
$\overline{C}_p(X) \times \overline{C}_q(X)$.  The Alexander--Whitney
map assembles these summands into the cotensor product
$\bar{B}(M) \square_{\bar{B}(\mathcal{A})} \bar{B}(N)$, defining
$\mu_{M,N}$.

\emph{Step~2: Quasi-isomorphism for Koszul modules.}
When $M$ is a Koszul module, the bar spectral sequence
$E_2^{p,q} = \mathrm{Ext}^p_{\mathcal{A}}(\mathbf{1}, M)$
degenerates.  The K\"unneth theorem on
$\overline{C}_p(X) \times \overline{C}_q(X)$ then gives
$H^*(\bar{B}(M \boxtimes N))
\cong H^*(\bar{B}(M)) \otimes_{H^*(\bar{B}(\mathcal{A}))}
H^*(\bar{B}(N))$
(the required Tor-independence holds because, for Koszul modules,
$H^*(\bar{B}(M)) = M^!$ is concentrated on the diagonal and hence
flat over $\mathcal{A}^!$ by the linearity of the bar resolution),
so $\mu_{M,N}$ induces an isomorphism on cohomology.

\emph{Step~3: Braiding reversal.}
The argument proceeds through four ingredients:

\noindent (a) \emph{KZ connection from the bar propagator.}
The KZB connection on the bundle of conformal blocks is realized
as the bar complex propagator $\eta_{12}$ on
$\overline{C}_2(\mathbb{P}^1)$
(Proposition~\ref{prop:kzb-bar-complex}): the flat connection is
$\nabla = d - \sum_{i<j} \Omega_{ij}/(k+h^\vee) \cdot \eta_{ij}$,
where $\Omega_{ij}$ is the Casimir insertion and $\eta_{ij}$ is
the propagator $1$-form on the configuration space.

\noindent (b) \emph{Monodromy and $R$-matrix.}
The monodromy of the KZ connection around the divisor
$D_{12} = \{z_1 = z_2\} \subset \overline{C}_2(\mathbb{P}^1)$
is the $R$-matrix $R \in \mathrm{End}(M \otimes N)$.  By the
Kohno--Drinfeld theorem~\cite{Kohno87, Drinfeld90}, this monodromy
representation is equivalent to the universal $R$-matrix of the
quantum group $U_q(\mathfrak{g})$ at $q = e^{\pi i/(k+h^\vee)}$.

\noindent (c) \emph{Verdier involution reverses cycle orientation.}
The Koszul duality functor $\Phi$ applies the Verdier involution
$\sigma$ (Theorem~\ref{thm:bar-cobar-verdier}) to the bar complex.
On $\overline{C}_2(\mathbb{P}^1)$, $\sigma$ acts as Poincar\'{e}
duality, which reverses the orientation of the cycle
$\gamma_{12}$ encircling $D_{12}$.

\noindent (d) \emph{Orientation reversal inverts $R$.}
The monodromy around $\gamma_{12}^{-1}$ (the reversed cycle) is
$R^{-1}$.  Therefore $\Phi$ transforms
$\sigma \mapsto \sigma^{-1}$, i.e., the braiding on the dual
category is the inverse braiding.
For the Heisenberg algebra, this recovers the phase
reversal $e^{2\pi i\lambda\mu/\kappa} \mapsto
e^{-2\pi i\lambda\mu/\kappa}$ of
Remark~\ref{rem:fusion-koszul}.
\end{proof}

\begin{remark}[Towards a derived Drinfeld--Kohno theorem]\label{rem:towards-derived-dk}
\index{Drinfeld--Kohno!derived!preview}
The mechanism above---Verdier duality reverses collision cycle
orientation, hence inverts the $R$-matrix monodromy---admits a
homotopy-coherent enhancement.  Ordered configuration spaces on
the line carry the braid group action (Arnol'd), and bar-cobar
duality implements Verdier duality on these configuration spaces.
A \emph{derived Drinfeld--Kohno theorem} would assert that
this extends to an equivalence of $\mathsf{E}_1$-factorization
categories whose braid monodromy realizes the Drinfeld--Kohno
correspondence up to derived completion: the bar-cobar-enhanced
$\mathsf{E}_1$-factorization data of the Yangian $Y(\mathfrak{g})$
is equivalent to that of $U_q(\mathfrak{g})^{\mathrm{op}}$ with
reversed $R$-matrix.  See
Conjecture~\textup{\ref{conj:derived-drinfeld-kohno}} for the
precise formulation.
\end{remark}

\begin{remark}[Quantum groups and braiding reversal]
\label{rem:quantum-group-braiding}
For $\mathcal{A} = \widehat{\mathfrak{g}}_k$ at positive integer level,
the monoidal module Koszul duality combined with the
Kazhdan--Lusztig equivalence
(Theorem~\ref{thm:kazhdan-lusztig-equivalence}) gives a commutative
diagram of braided monoidal categories:
\[
\begin{tikzcd}
\mathcal{O}_k^{\mathrm{int}}(\widehat{\mathfrak{g}})
  \ar[r, "\Phi_{\mathrm{KL}}", "\sim"']
  \ar[d, "\Phi"']
& \mathrm{Rep}(U_q(\mathfrak{g}))
  \ar[d, "q \mapsto q^{-1}"]
\\
\mathcal{O}_{k'}(\widehat{\mathfrak{g}})
  \ar[r, "\Phi_{\mathrm{KL}}", "\sim"']
& \mathrm{Rep}(U_{q^{-1}}(\mathfrak{g}))
\end{tikzcd}
\]
where $q = e^{\pi i/(k + h^\vee)}$ and $k' = -k - 2h^\vee$.
Since $k' < 0$, the bottom-left entry is the full category~$\mathcal{O}$
(there are no integrable modules at negative level); the KL equivalence
at level~$k'$ uses the Kazhdan--Lusztig--Finkelberg category
(see \cite{KL93}).  The
vertical arrow $q \mapsto q^{-1}$ reverses the braiding
($R_q \mapsto R_{q^{-1}} = R_q^{-1}$) while preserving the tensor
product (since the underlying fusion rules depend only on
$\mathfrak{g}$, not on the sign of~$q$).
\end{remark}

\begin{proposition}[Ext--Tor exchange via module Koszul duality; \ClaimStatusProvedHere]
\label{prop:ext-tor-exchange}
\index{Ext--Tor exchange}
Let $(\mathcal{A}, \mathcal{A}^!)$ be a Koszul chiral pair with
$\mathcal{A}$ uncurved ($m_0 = 0$, i.e., genus~$0$ or critical
level) and with finite-type grading (each conformal weight space
is finite-dimensional).  For $\mathcal{A}$-modules $M, N$ with Koszul dual
comodules $M^! = \bar{B}_{\mathcal{A}}(M)$,
$N^! = \bar{B}_{\mathcal{A}}(N)$:
\begin{equation}\label{eq:ext-tor-exchange}
\mathrm{Ext}^n_{\mathcal{A}}(M, N)
\;\cong\;
\mathrm{Tor}_n^{\mathcal{A}^!}(M^!, N^!)^{\vee}
\end{equation}
where $\vee$ denotes the graded linear dual.
\end{proposition}

\begin{proof}
By the module Koszul duality equivalence
(Remark~\ref{rem:koszul-module-categories}),
$\Phi: \mathrm{Mod}^{\mathrm{compl}}_{\mathcal{A}}
\xrightarrow{\sim}
\mathrm{CoMod}^{\mathrm{conil}}_{\mathcal{A}^{!,c}}$.
Under this equivalence, the internal Hom and cotensor product
are interchanged:
\begin{align*}
\underline{\mathrm{Hom}}_{\mathcal{A}}(M, N) &\simeq
\underline{\mathrm{Hom}}_{\mathcal{A}^{!,c}}
(\Phi(M), \Phi(N)), \\
M \boxtimes_{\mathcal{A}} N &\simeq
\Phi^{-1}(\Phi(M) \square_{\mathcal{A}^{!,c}} \Phi(N)).
\end{align*}
Taking cohomology of the first line gives the Ext exchange.
For the Tor identification, the cotensor product
$\Phi(M) \square_{\mathcal{A}^{!,c}} \Phi(N)$
of conilpotent comodules dualizes (by finite-type grading)
to the tensor product of modules over $\mathcal{A}^!$:
\[
H_n(\Phi(M) \square_{\mathcal{A}^{!,c}} \Phi(N))^{\vee}
\;\cong\;
\mathrm{Tor}_n^{\mathcal{A}^!}(M^!, N^!).
\]
Combining these identifications
yields~\eqref{eq:ext-tor-exchange}.
\end{proof}

\begin{remark}[Deformation--extension duality]\label{rem:ext-tor-interpretation}
Equation~\eqref{eq:ext-tor-exchange} has a concrete
interpretation: extensions of $\mathcal{A}$-modules
(classified by $\mathrm{Ext}^1_{\mathcal{A}}$, which encode
first-order deformations of~$M$) are dual to torsion
products of $\mathcal{A}^!$-modules
(which encode co-extensions of~$M^!$).  For
$\mathcal{A} = \widehat{\mathfrak{g}}_k$ and $M$ a Verma module,
$\mathrm{Ext}^1(M(\lambda), M(\mu))$ classifies intertwining
operators, while
$\mathrm{Tor}_1^{\mathcal{A}^!}(M(\lambda)^!, M(\mu)^!)$
classifies dual singular vectors.
In the curved case ($m_0 \neq 0$), the Ext and Tor groups
must be replaced by their curved analogues in the sense of
Positselski~\cite{Positselski11}; see
Remark~\ref{rem:koszul-module-categories}.
\end{remark}

\subsection{Conformal blocks and the bar complex}
\label{sec:conformal-blocks}

\begin{definition}[Space of conformal blocks]
\label{def:conformal-blocks}
\index{conformal blocks|textbf}
Let $\mathcal{A}$ be a chiral algebra on a smooth projective
curve $X$ of genus~$g$, and let
$\mathcal{M}_1, \ldots, \mathcal{M}_r$ be $\mathcal{A}$-modules
inserted at distinct points $x_1, \ldots, x_r \in X$.  The
\emph{space of conformal blocks} is the space of coinvariants:
\begin{equation}\label{eq:conformal-blocks}
\mathbb{V}(\mathcal{A};\, \mathcal{M}_1, \ldots, \mathcal{M}_r;\, X)
= H^0\bigl(
  X \setminus \{x_1, \ldots, x_r\},\;
  \mathcal{A} \bigr) \Big\backslash
  \mathcal{M}_1 \otimes \cdots \otimes \mathcal{M}_r.
\end{equation}
Here the left action is by the residues of the chiral bracket
along the insertions $x_i$.  When $r = 0$ (no module insertions),
$\mathbb{V}(\mathcal{A}; X) = H^0(X, \mathcal{A}) \backslash
\mathbf{1}$ is the \emph{vacuum conformal block}.
\end{definition}

\begin{proposition}[Conformal blocks via the bar complex; \ClaimStatusProvedHere]
\label{prop:conformal-blocks-bar}
\index{conformal blocks!bar complex computation}
For a Koszul chiral algebra $\mathcal{A}$ on a curve $X$ of
genus~$g$ with module insertions
$\mathcal{M}_1, \ldots, \mathcal{M}_r$, the space of conformal
blocks is computed by the bar complex on $X$:
\begin{equation}\label{eq:conformal-blocks-bar}
\mathbb{V}(\mathcal{A};\, \mathcal{M}_1, \ldots, \mathcal{M}_r;\, X)
\;\cong\;
H^0\bigl(
  \mathrm{Tot}\,
  \bar{B}(\mathcal{M}_1, \ldots, \mathcal{M}_r;\, X)
\bigr)
\end{equation}
where $\bar{B}(\mathcal{M}_1, \ldots, \mathcal{M}_r; X)$ is the
bar complex on $\overline{C}_{n+r}(X)$ with module data at the
first $r$ marked points and algebra data at the remaining~$n$
points.
\end{proposition}

\begin{proof}
The bar complex on $\overline{C}_{n+r}(X)$ computes the derived
coinvariants of $\mathcal{A}$ acting on
$\mathcal{M}_1 \otimes \cdots \otimes \mathcal{M}_r$.  The
$n$-th term of the bar complex is:
\[
\bar{B}_n(\vec{\mathcal{M}}; X) = \Gamma\bigl(
\overline{C}_{n+r}(X),\;
j_*j^*\bigl(\mathcal{M}_1 \boxtimes \cdots \boxtimes
\mathcal{M}_r \boxtimes \mathcal{A}^{\boxtimes n}\bigr)
\otimes \Omega^n_{\log}
\bigr)
\]
where $j\colon C_{n+r}(X) \hookrightarrow \overline{C}_{n+r}(X)$
and $\Omega^n_{\log}$ are the logarithmic forms with poles along
the boundary divisors.  The differential is
$d = d_{\mathrm{config}} + d_{\mathrm{internal}}$ as in
Definition~\ref{def:geometric-bar}, with $d_{\mathrm{config}}$
encoding the chiral bracket residues of $\mathcal{A}$ acting on
each~$\mathcal{M}_i$.

By the Koszul property, the spectral sequence associated to the
bar filtration degenerates at~$E_2$, and the $E_1$ page computes
$H^*(X \setminus \{x_i\}, \mathcal{A}) \backslash
\mathcal{M}_1 \otimes \cdots \otimes \mathcal{M}_r$, which is
precisely the conformal block.  The $E_2$ degeneration ensures that
$H^0(\mathrm{Tot}\, \bar{B})$ is exactly the conformal block
space, without higher corrections.
\end{proof}

\begin{remark}[Conformal blocks and Koszul duality]
\label{rem:conformal-blocks-koszul}
\index{conformal blocks!Koszul duality}
The conformal block computation via bar complexes has a natural
Koszul dual:
\[
\mathbb{V}(\mathcal{A};\, \vec{\mathcal{M}};\, X)
\;\cong\;
\mathbb{V}(\mathcal{A}^!;\, \vec{\mathcal{M}}^!;\, X)
\]
for Koszul modules $\mathcal{M}_i$ with duals $\mathcal{M}_i^!$,
by the bar-cobar quasi-isomorphism applied termwise.  This
identification gives a geometric proof that the dimension
$\dim \mathbb{V}$ is invariant under the Feigin--Frenkel level shift
$k \mapsto -k - 2h^\vee$.  This is a consequence of chiral Koszul
duality, distinct from the classical level-rank duality
$\widehat{\mathfrak{sl}}_N$ at level $k$ $\leftrightarrow$
$\widehat{\mathfrak{sl}}_k$ at level $N$
(\cite{Beauville}, \cite{NakanishiTsuchiya}).

For $\widehat{\mathfrak{sl}}_2$ at level~$k$ on
$\mathbb{P}^1$ with $r$ insertions of weight $\lambda_i$,
the Verlinde formula gives
$\dim \mathbb{V} = \#\{(\lambda_1, \ldots, \lambda_r) :
\sum \lambda_i \in 2\mathbb{Z},\, \text{fusion rules satisfied}\}$.
The bar complex computes this via the Euler characteristic of
$\overline{C}_{n+r}(\mathbb{P}^1)$.
\end{remark}

\begin{definition}[Bundle of conformal blocks]
\label{def:bundle-conformal-blocks}
\index{conformal blocks!bundle|textbf}
\index{factorization!conformal blocks}
Let $\mathcal{A}$ be a rational chiral algebra and
$\Lambda_1, \ldots, \Lambda_r$ a collection of highest weights.
As the pointed curve $(X; x_1, \ldots, x_r)$ varies over the moduli
space $\mathcal{M}_{g,r}$, the spaces of conformal blocks
(Definition~\ref{def:conformal-blocks}) assemble into a vector bundle
\begin{equation}\label{eq:conformal-block-bundle}
\mathbb{V}(\mathcal{A};\, \Lambda_1, \ldots, \Lambda_r)
\;\to\; \mathcal{M}_{g,r}
\end{equation}
of finite rank, called the \emph{bundle of conformal blocks} (or
\emph{bundle of vacua}).
\begin{enumerate}[label=\textup{(\roman*)}]
\item The bundle carries a projective flat connection: the
  \emph{Knizhnik--Zamolodchikov--Bernard (KZB) connection},
  whose horizontal sections satisfy the KZB equations.
\item The rank of $\mathbb{V}$ is independent of the genus-$g$
  curve $X$ and depends only on the algebra $\mathcal{A}$, the
  weights $\Lambda_i$, and the genus~$g$.
\item The factorization property of $\mathcal{A}$ implies a
  \emph{factorization} of the bundle under degeneration: when $X$
  degenerates to a nodal curve $X_0 = X_1 \cup_p X_2$, the fiber
  satisfies
  \[
  \mathbb{V}(\mathcal{A};\, \vec{\Lambda};\, X_0)
  \;\cong\;
  \bigoplus_{\mu \in \mathrm{Irr}(\mathcal{A})}
  \mathbb{V}(\mathcal{A};\, \vec{\Lambda}_1, \mu;\, X_1)
  \otimes
  \mathbb{V}(\mathcal{A};\, \mu^*, \vec{\Lambda}_2;\, X_2).
  \]
\end{enumerate}
\end{definition}

\begin{proposition}[KZB connection from the bar complex; \ClaimStatusProvedHere]
\label{prop:kzb-bar-complex}
\index{KZB connection!bar complex}
The projective flat connection on
$\mathbb{V}(\mathcal{A}; \vec{\Lambda})$ arises from the
bar complex as follows.  The bar complex
$\bar{B}(\vec{\mathcal{M}}; X)$ of
Proposition~\textup{\ref{prop:conformal-blocks-bar}} depends on the
complex structure of~$X$ through the propagator
$\eta_{ij} \in \Omega^1(\overline{C}_2(X))$.  The variation
of $\eta_{ij}$ with respect to moduli $t_\alpha \in T\mathcal{M}_{g,r}$
yields:
\[
\nabla_{\partial/\partial t_\alpha}
= \partial_{t_\alpha}
+ \sum_{i<j}
\mathrm{Res}_{D_{ij}} \frac{\partial \eta_{ij}}{\partial t_\alpha},
\]
where the residue operator acts on bar complex terms on
$\overline{C}_{n+r}(X)$.  Flatness
$[\nabla_\alpha, \nabla_\beta] = 0$ follows from $d^2 = 0$ on
the bar complex (Theorem~\ref{thm:bar-nilpotency-complete}):
the flatness obstruction equals the commutator of bar
differentials, which vanishes.

For $\mathcal{A} = \widehat{\mathfrak{g}}_k$ on $\mathbb{P}^1$,
this recovers the KZ connection
$\nabla_i = \partial_{z_i} + \frac{1}{k + h^\vee}
\sum_{j \neq i} \frac{\Omega_{ij}}{z_i - z_j}$, where
$\Omega_{ij} = \sum_a T^a_i \otimes T^a_j$ is the Casimir
element.
\end{proposition}

\begin{proof}
The bar complex $\bar{B}(\vec{\mathcal{M}}; X)$ is a complex of
sheaves on $\mathcal{M}_{g,r}$ (via the universal curve
$\pi\colon \mathcal{C}_{g,r} \to \mathcal{M}_{g,r}$), and its
cohomology sheaf $R^0\pi_*(\bar{B})$ is the bundle
$\mathbb{V}$.  The Gauss--Manin connection on
$R^0\pi_*(\bar{B})$ is the KZB connection.

The propagator $\eta_{ij}$ on $\overline{C}_2(X)$ is the unique
logarithmic $1$-form with simple pole along the diagonal $D_{ij}$
and residue $+1$ (the Cauchy kernel on the curve).  On
$\mathbb{P}^1$: $\eta_{ij} = dz_i/(z_i - z_j)$.  Its variation
$\partial_{t_\alpha} \eta_{ij}$ contributes a term to the
connection that acts on $\mathcal{M}_i \otimes \mathcal{M}_j$
via the chiral bracket --- this is the Casimir insertion
$\Omega_{ij}/(k + h^\vee)$ (the factor $(k + h^\vee)^{-1}$
is the inverse of the quadratic pole coefficient in the OPE).

Flatness: $[\nabla_\alpha, \nabla_\beta]$ involves
$\partial_\alpha \partial_\beta \eta_{ij}
- \partial_\beta \partial_\alpha \eta_{ij} = 0$ (symmetry
of second derivatives) plus commutators of residue operators.
The commutator of residue operators at $D_{ij}$ and $D_{kl}$
vanishes when $\{i,j\} \cap \{k,l\} = \emptyset$ (disjoint
divisors).  When $\{i,j\} \cap \{k,l\} \neq \emptyset$, the
commutator is a residue at the codimension-$2$ stratum
$D_{ij} \cap D_{kl}$, which is controlled by $d^2 = 0$ on
the bar complex.
\end{proof}

\begin{theorem}[Verlinde formula via the bar complex; \ClaimStatusProvedElsewhere]
\label{thm:verlinde-bar}
\index{Verlinde formula!bar complex|textbf}
Let $\mathcal{A} = \widehat{\mathfrak{g}}_k$ at positive integer
level~$k$, with integrable highest weights
$\Lambda_1, \ldots, \Lambda_r$.  The dimension of the space of
conformal blocks on a genus-$g$ curve is:
\begin{equation}\label{eq:verlinde}
\dim \mathbb{V}(\widehat{\mathfrak{g}}_k;\,
\Lambda_1, \ldots, \Lambda_r;\, \Sigma_g)
= \sum_{\mu \in P_+^k}
\prod_{i=1}^r \frac{S_{\Lambda_i,\mu}}{S_{0,\mu}}
\cdot (S_{0,\mu})^{2-2g}
\end{equation}
where $S_{\Lambda,\mu}$ is the modular $S$-matrix of
$\widehat{\mathfrak{g}}_k$\textup{:}
\[
S_{\Lambda,\mu}
= C \sum_{w \in W} \varepsilon(w)\,
e^{-2\pi i \langle w(\Lambda + \rho),\, \mu + \rho \rangle/(k + h^\vee)},
\qquad
C = \frac{i^{|\Delta^+|}}{|P/(k+h^\vee)Q|^{1/2}}.
\]
\end{theorem}

\begin{proof}
The proof assembles three ingredients.

\emph{Step~1: Factorization reduction.}
By the factorization property of conformal blocks
(Definition~\ref{def:bundle-conformal-blocks}(iii)), the
genus-$g$ conformal block decomposes under degeneration to a
maximally degenerate curve (a tree of $\mathbb{P}^1$'s with
$g$ self-intersections):
\begin{equation}\label{eq:verlinde-factorization}
\dim \mathbb{V}(g, \vec{\Lambda})
= \sum_{\mu_1, \ldots, \mu_g \in \mathrm{Irr}}
\prod_{\text{vertices}} N_{\Lambda_a \Lambda_b}^{\Lambda_c}
\end{equation}
where $N_{\Lambda \mu}^{\nu}$ are the fusion coefficients
(= dimensions of genus-$0$ three-point conformal blocks).

\emph{Step~2: Diagonalization via the $S$-matrix.}
The fusion coefficients are diagonalized by the modular
$S$-matrix: the Verlinde formula for fusion coefficients gives
\[
N_{\Lambda\mu}^{\nu}
= \sum_{\sigma \in P_+^k}
\frac{S_{\Lambda,\sigma}\, S_{\mu,\sigma}\,
\overline{S_{\nu,\sigma}}}{S_{0,\sigma}}.
\]
Substituting into~\eqref{eq:verlinde-factorization} and summing
over the internal weights $\mu_1, \ldots, \mu_g$ using the
orthogonality $\sum_\mu S_{\mu,\sigma} \overline{S_{\mu,\tau}}
= \delta_{\sigma,\tau}$, the sum telescopes to the claimed
formula~\eqref{eq:verlinde}.

\emph{Step~3: Bar complex interpretation.}
In the bar-cobar framework, each ingredient has a geometric
counterpart:
\begin{itemize}
\item The fusion coefficients $N_{\Lambda\mu}^{\nu}$ are computed
  by the bar complex on $\overline{C}_3(\mathbb{P}^1)$
  (Proposition~\ref{prop:conformal-blocks-bar} with $g = 0$,
  $r = 3$).
\item The modular $S$-matrix is the monodromy of the KZB connection
  (Proposition~\ref{prop:kzb-bar-complex}) around the $A$-cycle
  of a genus-$1$ curve.
\item The factorization property is the bar complex's
  factorization along boundary strata of $\overline{C}_{n+r}(X)$
  (Theorem~\ref{thm:geometric-bar-definition}).
\end{itemize}
The Koszul property ensures that the bar spectral sequence
degenerates, so the rank of $\mathbb{V}$ is computed exactly by
the $E_2$ page.
\end{proof}

\begin{remark}[Verlinde formula and Koszul duality]
\label{rem:verlinde-koszul}
\index{Verlinde formula!Koszul duality}
The Verlinde formula~\eqref{eq:verlinde} applies at positive
integer level $k$, where $\widehat{\mathfrak{g}}_k$ has finitely
many integrable highest weights.  The Koszul dual level
$k' = -k - 2h^\vee < 0$ lies outside this regime: at $k' < 0$
there are no integrable highest weights, and the Verlinde formula
does not directly apply.

Nonetheless, the complementarity theorem
(Theorem~\ref{thm:quantum-complementarity-main}) implies that
the genus-$g$ quantum corrections satisfy
$Q_g(\widehat{\mathfrak{g}}_k) +
Q_g(\widehat{\mathfrak{g}}_{k'}^!)
= H^*(\mathcal{M}_g, Z(\widehat{\mathfrak{g}}_k))$,
which constrains the conformal block dimensions at
level~$k'$ through the center of the algebra.  At admissible
levels (where both $k$ and $k'$ are admissible and the
Verlinde formula extends to the admissible category), the
complementarity gives an explicit relationship between the
conformal block ranks of the two sides.  For positive integer
level, the bar complex computes the Verlinde dimensions
exactly (Theorem~\ref{thm:verlinde-bar}), and the duality
acts on the bar complex via the Verdier pairing
(Theorem~\ref{thm:verdier-duality-operations}).
\end{remark}

\subsection{Conformal block ranks under Koszul duality}
\label{sec:conformal-blocks-kd}

\begin{proposition}[Conformal block duality; \ClaimStatusProvedHere]
\label{prop:conformal-block-duality}
\index{conformal blocks!Koszul duality}
Let $(\mathcal{A}, \mathcal{A}^!)$ be a Koszul pair of chiral
algebras, both $C_2$-cofinite, with finitely many simple modules.
Denote by
$\Phi\colon \mathrm{Irr}(\mathcal{A})
\xrightarrow{\sim} \mathrm{Irr}(\mathcal{A}^!)$ the bijection
on simple modules induced by bar-cobar duality
\textup{(}the heart of the Koszul t-structure, see
Proposition~\textup{\ref{prop:koszul-t-structures})}.
Then the conformal block ranks satisfy:
\begin{enumerate}[label=\textup{(\roman*)}]
\item \textup{(}Genus~$0$\textup{)} The fusion coefficients are
  preserved: $N_{\Lambda\mu}^{\nu}(\mathcal{A})
  = N_{\Phi(\Lambda)\Phi(\mu)}^{\Phi(\nu)}(\mathcal{A}^!)$.
\item \textup{(}Genus~$g$\textup{)}  If $\mathcal{A}$ is rational,
  then $\dim \mathbb{V}(\mathcal{A};\, \vec{\Lambda};\, \Sigma_g)
  = \dim \mathbb{V}(\mathcal{A}^!;\, \Phi(\vec{\Lambda});\,
  \Sigma_g)$.
\item The KZB connection on the bundle of conformal blocks
  for $\mathcal{A}$ is \emph{dual} to the KZB connection for
  $\mathcal{A}^!$: the monodromy representations are related
  by the bar-cobar functor.
\end{enumerate}
\end{proposition}

\begin{proof}
For (i): the monoidal module Koszul duality
(Theorem~\ref{thm:monoidal-module-koszul}) provides a braided
tensor equivalence $\Phi\colon \mathcal{C}_{\mathcal{A}} \to
\mathcal{C}_{\mathcal{A}^!}^{\mathrm{rev}}$
on the module categories (with reversed braiding).  Since $\Phi$
is a tensor functor, it preserves the fusion product:
$\Phi(\mathcal{L}(\Lambda) \boxtimes \mathcal{L}(\mu))
\cong \Phi(\mathcal{L}(\Lambda)) \boxtimes
\Phi(\mathcal{L}(\mu))$.  In particular,
$N_{\Lambda\mu}^{\nu}
= \dim \mathrm{Hom}_{\mathcal{A}}(\mathcal{L}(\nu),
\mathcal{L}(\Lambda) \boxtimes \mathcal{L}(\mu))
= \dim \mathrm{Hom}_{\mathcal{A}^!}(\Phi(\mathcal{L}(\nu)),
\Phi(\mathcal{L}(\Lambda)) \boxtimes \Phi(\mathcal{L}(\mu)))
= N_{\Phi(\Lambda)\Phi(\mu)}^{\Phi(\nu)}$.

For (ii): since $\Phi$ is a tensor equivalence preserving fusion
coefficients, it preserves the fusion matrices
$(N_\lambda)_{\mu\nu} = N_{\lambda\mu}^\nu$.  The quantum
dimensions $d_\lambda = S_{0\lambda}/S_{00}$ are the
Frobenius--Perron eigenvalues of these matrices, hence
$d_{\Phi(\lambda)} = d_\lambda$.  The Verlinde formula
expresses genus-$g$ conformal block dimensions as
$\dim \mathbb{V}(\Sigma_g) = \sum_\lambda d_\lambda^{2-2g}$,
which depends only on the quantum dimensions.  Since $\Phi$ is a
bijection on simple objects preserving quantum dimensions, the
genus-$g$ ranks agree.  The $C_2$-cofiniteness of $\mathcal{A}^!$
(a hypothesis) ensures that conformal blocks are finite-dimensional
(Zhu~\cite{Zhu96}), and the factorization property holds for all
$C_2$-cofinite algebras (Huang--Lepowsky--Zhang~\cite{HLZ}).

For (iii): the KZB connection arises from the variation of
the propagator $\eta_{ij}$ on $\overline{C}_2(X)$
(Proposition~\ref{prop:kzb-bar-complex}).  The propagator
depends only on the curve~$X$, not on the algebra
$\mathcal{A}$: it is the Cauchy kernel $dz/(z-w)$ on
$\mathbb{P}^1$ and the Szeg\H{o} kernel on higher-genus
curves.  The algebra enters through the Casimir insertion
$\Omega_{ij}/(k + h^\vee)$.  Under Koszul duality, the Casimir
$\Omega$ is exchanged for the dual Casimir $\Omega^!$ with
level shift $k \mapsto k'$, and the monodromy representations
are intertwined by the tensor functor $\Phi$.
\end{proof}

\begin{remark}[Conformal block duality and complementarity]
\label{rem:conformal-block-complementarity}
Proposition~\ref{prop:conformal-block-duality} is consistent
with the complementarity theorem
(Theorem~\ref{thm:quantum-complementarity-main}): the genus-$g$
quantum corrections $Q_g(\mathcal{A})$ and
$Q_g(\mathcal{A}^!)$ are exchanged by Verdier duality, and the
conformal block ranks are invariants of this exchange.
In particular, when the $S$-matrices satisfy
$S^{\mathcal{A}}_{\Lambda,\mu}
= \overline{S^{\mathcal{A}^!}_{\Phi(\Lambda),\Phi(\mu)}}$
(as occurs for Feigin--Frenkel dual affine algebras at admissible
levels), the Verlinde formula yields identical ranks directly.

The hypothesis ``both $C_2$-cofinite'' is essential: Koszul duality
can map a rational algebra to a non-rational one (e.g.,
$V_1(\mathfrak{sl}_2)$ is rational with $X_{V_1} = \{0\}$, but
$V_{-3}(\mathfrak{sl}_2)$ has $X_{V_{-3}} = \mathfrak{sl}_2^*$).
Part~(ii) of the proposition does not apply in such cases.
\end{remark}

\subsection{Derived categories and t-structures}
\label{sec:derived-t-structures}

\begin{proposition}[Koszul duality and t-structures; \ClaimStatusProvedHere]
\label{prop:koszul-t-structures}
\index{t-structure!Koszul duality}
\index{derived category!module Koszul duality}
Let $(\mathcal{A}, \mathcal{A}^!)$ be a Koszul pair of chiral
algebras, and let $D^b(\mathcal{O}_k)$,
$D^b(\mathcal{O}_{k'})$ be the bounded derived categories of
their respective categories~$\mathcal{O}$.  The module Koszul
duality functor $\Phi$ of
Remark~\textup{\ref{rem:koszul-module-categories}} induces an
equivalence of triangulated categories
\begin{equation}\label{eq:derived-koszul}
\Phi\colon D^b(\mathcal{O}_k)
\;\xrightarrow{\;\sim\;}\;
D^b(\mathcal{O}_{k'})
\end{equation}
that exchanges the standard t-structure $(\mathcal{D}^{\leq 0},
\mathcal{D}^{\geq 0})$ on $D^b(\mathcal{O}_k)$ with the
\emph{Koszul t-structure} on $D^b(\mathcal{O}_{k'})$:
\begin{align}
\Phi(\mathcal{D}^{\leq 0}_{\mathrm{std}})
&= \mathcal{D}^{\leq 0}_{\mathrm{Kosz}}, \\
\Phi(\mathcal{L}(\lambda))
&\in \mathcal{D}^{\leq 0}_{\mathrm{Kosz}}
  \cap \mathcal{D}^{\geq 0}_{\mathrm{Kosz}}
\quad \text{for all } \lambda.
\end{align}
The heart of the Koszul t-structure on $D^b(\mathcal{O}_{k'})$
is the Serre quotient
$\mathcal{O}_{k'}^{\mathrm{Kosz}} \simeq \mathcal{O}_k$.
\end{proposition}

\begin{proof}
The module-comodule equivalence of
Remark~\ref{rem:koszul-module-categories} restricts to category
$\mathcal{O}$ because the bar resolution of a highest-weight module
has terms in category $\mathcal{O}$ (the bar differential preserves
the highest-weight condition and finite-dimensionality of weight
spaces, since OPE residues do not increase conformal weight).
The bar resolution of a simple module
$\mathcal{L}(\lambda) \in \mathcal{O}_k$ thus produces a complex in
$D^b(\mathcal{O}_{k'})$ whose terms are modules over $\mathcal{A}^!$.
For a Koszul algebra, the bar resolution is \emph{linear}: the
$n$-th term $\bar{B}_n(\mathcal{L}(\lambda))$ is concentrated in
internal degree $n + h(\lambda)$, where $h(\lambda)$ is the
conformal weight.

This linearity defines the Koszul t-structure: an object
$M^\bullet \in D^b(\mathcal{O}_{k'})$ is in
$\mathcal{D}^{\leq 0}_{\mathrm{Kosz}}$ if and only if
$H^n(M^\bullet)$ is concentrated in internal degrees $\geq n$
for all~$n$.  The Koszul property ensures that $\Phi(\mathcal{L}(\lambda))$
lies in the heart of this t-structure.

The identification of the heart with $\mathcal{O}_k$ follows from
the invertibility of the bar-cobar adjunction: $\Omega \circ B \simeq
\mathrm{id}$ on Koszul modules, so the bar construction provides a
bijection between simples of $\mathcal{O}_k$ and simples of the
Koszul heart of $\mathcal{O}_{k'}$, preserving composition series.
\end{proof}

\begin{remark}[Geometric interpretation of t-structure exchange]
\label{rem:t-structure-geometry}
In the configuration space picture, the Koszul t-structure has a
transparent geometric meaning: the bar complex
$\bar{B}_\bullet(\mathcal{L}(\lambda))$ on
$\overline{C}_{n+1}(X)$ has a filtration by the number of
boundary strata.  The standard t-structure corresponds to the
filtration by bar degree (number of points in configuration
space), while the Koszul t-structure corresponds to the
filtration by \emph{internal degree} (the total conformal weight).

For the Kazhdan--Lusztig equivalence
(Theorem~\ref{thm:kazhdan-lusztig-equivalence}), the Koszul
t-structure on $D^b(\mathrm{Rep}(U_{q^{-1}}(\mathfrak{g})))$
corresponds to the t-structure associated with the canonical
basis of $U_{q^{-1}}$ --- the basis elements are precisely the
images of the simple modules under bar-cobar duality.
\end{remark}

\section{Highest weight modules}\label{sec:highest-weight-modules}

The chiral algebras of primary interest---affine Kac--Moody, Virasoro,
$\mathcal{W}$-algebras---carry a conformal grading, and their module
theory is governed by highest weight representations.  We establish the
definitions and structural results needed for the resolution and
character computations below.

\begin{definition}[Weight grading and highest weight modules]
\label{def:highest-weight-module}
\index{highest weight module|textbf}
Let $\mathcal{A}$ be a chiral algebra equipped with a conformal vector
(a Virasoro subalgebra generated by $T(z) = \sum_{n} L_n z^{-n-2}$).
A \emph{weight grading} on an $\mathcal{A}$-module $\mathcal{M}$ is a
decomposition $\mathcal{M} = \bigoplus_{h} \mathcal{M}_h$ into
generalized eigenspaces of $L_0$.

A module $\mathcal{M}$ is \emph{highest weight} of weight~$\lambda$
(encoding the $L_0$-eigenvalue~$h$ together with any additional Cartan
data) if $\mathcal{M} = \mathcal{A} \cdot v_\lambda$ is generated by a
vector $v_\lambda$ satisfying:
\begin{enumerate}[label=(\roman*)]
\item $\mathcal{A}_{>0} \cdot v_\lambda = 0$ (annihilation by
   positive modes);
\item $L_0 \cdot v_\lambda = h \cdot v_\lambda$ (weight condition);
\item for affine algebras, $\mathfrak{h} \cdot v_\lambda = \lambda|_\mathfrak{h} \cdot v_\lambda$
   (Cartan eigenvalue).
\end{enumerate}
\end{definition}

\begin{definition}[Verma module]\label{def:verma-module}
\index{Verma module|textbf}
The \emph{Verma module} $\mathcal{M}(\lambda)$ is the universal highest
weight module of weight~$\lambda$: it is induced from the one-dimensional
representation of the non-negative subalgebra
$\mathcal{A}_{\geq 0} = \mathcal{A}_0 \oplus \mathcal{A}_{>0}$ on which
$v_\lambda$ has weight~$\lambda$ and $\mathcal{A}_{>0}$ acts by zero:
\[
\mathcal{M}(\lambda)
= \mathrm{Ind}_{\mathcal{A}_{\geq 0}}^{\mathcal{A}}\,\mathbb{C}_\lambda
= U(\mathcal{A}_{<0}) \otimes_\mathbb{C} \mathbb{C}_\lambda.
\]
Every highest weight module of weight~$\lambda$ is a quotient
of~$\mathcal{M}(\lambda)$.  The unique simple quotient is
$\mathcal{L}(\lambda) = \mathcal{M}(\lambda) / \mathcal{J}(\lambda)$,
where $\mathcal{J}(\lambda)$ is the maximal proper submodule.
\end{definition}

\begin{definition}[Shapovalov form]\label{def:shapovalov-form}
\index{Shapovalov form|textbf}
The \emph{Shapovalov form} on $\mathcal{M}(\lambda)$ is the unique
symmetric bilinear form
$\langle\cdot,\cdot\rangle_\lambda \colon
\mathcal{M}(\lambda) \times \mathcal{M}(\lambda) \to \mathbb{C}$
satisfying $\langle v_\lambda, v_\lambda \rangle_\lambda = 1$ and
the contravariance property
\[
\langle a_n \cdot u,\, w \rangle_\lambda
= \langle u,\, a_{-n} \cdot w \rangle_\lambda
\]
for all modes $a_n$ of generating fields.  Equivalently, $\langle u, w
\rangle_\lambda = \varepsilon(u \cdot \sigma(w))$, where $\sigma$ is the
anti-involution sending $a_n \mapsto a_{-n}$ and $\varepsilon$ is the
counit $\mathcal{M}(\lambda) \to \mathbb{C}_\lambda$.

The form is non-degenerate on the weight-$(h{+}n)$ subspace
$\mathcal{M}(\lambda)_{h+n}$ if and only if $\mathcal{M}(\lambda)$ has
no singular vectors of weight~$h{+}n$.  The \emph{Kac determinant}
$\det\langle\cdot,\cdot\rangle_{\lambda,n}$ on this subspace admits a
product formula (\cite{Kac}, Theorem~9.1; \cite{KR87}, Lecture~4) whose
vanishing locus determines the reducibility of~$\mathcal{M}(\lambda)$.
\end{definition}

\begin{proposition}[Generic irreducibility; \ClaimStatusProvedElsewhere]
\label{prop:generic-irreducibility}
\index{Kac determinant}
For the affine algebras $\widehat{\mathfrak{g}}_k$, the Virasoro algebra
$\mathrm{Vir}_c$, and the $\mathcal{W}$-algebras $\mathcal{W}^k(\mathfrak{g})$:
the Verma module $\mathcal{M}(\lambda)$ is irreducible for generic values
of the level\/central charge and the highest weight.  Precisely, the Kac
determinant is a non-trivial polynomial in the parameters, so its
vanishing defines a proper subvariety of the parameter space.

The reducibility locus includes:
\begin{enumerate}[label=(\roman*)]
\item \emph{Minimal model values:} $c_{p,q} = 1 - 6(p-q)^2/pq$
   for $\mathrm{Vir}_c$ (Feigin--Fuchs~\cite{FF84});
\item \emph{Admissible levels:} $k = -h^\vee + p/q$ for
   $\widehat{\mathfrak{g}}_k$ (Kac--Wakimoto~\cite{KW88});
\item \emph{Critical level:} $k = -h^\vee$ for
   $\widehat{\mathfrak{g}}_k$ (the Sugawara construction is
   undefined; see Remark~\ref{rem:feigin-frenkel-center}).
\end{enumerate}
\end{proposition}

\begin{remark}[Connection to Koszul duality]\label{rem:verma-koszul}
Generic irreducibility of Verma modules is the representation-theoretic
counterpart of the Koszul property.  At generic level, the Ext groups
\[
\mathrm{Ext}^q_{\mathcal{O}}(\mathcal{M}(\lambda),
\mathcal{M}(\mu)) = 0, \qquad q > 0,\; \lambda \neq \mu,
\]
vanish between distinct Verma modules; this is the
condition ensuring that the bar complex spectral sequence degenerates
at~$E_2$ (cf.\ the proof of Theorem~\ref{thm:virasoro-periodicity}).

At special levels---critical, admissible, minimal model---Verma modules
become reducible, the Ext groups become non-trivial, and the bar complex
acquires higher-order obstructions.  These are the quantum corrections
studied throughout Part~1 of this monograph.
\end{remark}

\subsection{Category \texorpdfstring{$\mathcal{O}$}{O} and structural results}

\begin{definition}[Category $\mathcal{O}$ for affine Lie algebras]
\label{def:category-O}
\index{Category O@Category $\mathcal{O}$|textbf}
Let $\widehat{\mathfrak{g}} = \mathfrak{g}[t,t^{-1}] \oplus
\mathbb{C}K \oplus \mathbb{C}d$ be an untwisted affine Lie algebra at
level~$k$ (i.e., $K$ acts as $k \cdot \mathrm{id}$).
\emph{Category~$\mathcal{O}_k(\widehat{\mathfrak{g}})$} is the full
subcategory of $\widehat{\mathfrak{g}}$-modules~$M$ satisfying:
\begin{enumerate}[label=\textup{(\roman*)}]
\item $M$ is finitely generated over $\widehat{\mathfrak{g}}$;
\item $\mathfrak{h} \oplus \mathbb{C}d$ acts semisimply on $M$
  (i.e., $M$ decomposes into weight spaces
  $M = \bigoplus_{\lambda} M_\lambda$
  with $\dim M_\lambda < \infty$);
\item $\mathfrak{n}_+[t] \oplus \mathfrak{n}_+ \oplus
  \mathfrak{h}[t] \cdot t$ acts locally nilpotently on $M$.
\end{enumerate}
The objects of $\mathcal{O}_k$ include all Verma modules
$\mathcal{M}(\lambda)$ and their simple quotients
$\mathcal{L}(\lambda)$ (Definition~\ref{def:verma-module}).
\end{definition}

\begin{definition}[Integrable modules]
\label{def:integrable-modules}
\index{integrable module|textbf}
The \emph{integrable} subcategory
$\mathcal{O}_k^{\mathrm{int}}(\widehat{\mathfrak{g}})
\subset \mathcal{O}_k(\widehat{\mathfrak{g}})$
consists of modules on which the Chevalley generators
$e_i, f_i$ ($i = 0, \ldots, \mathrm{rk}\,\mathfrak{g}$) act locally
nilpotently.  For $k$ a non-negative integer,
$\mathcal{O}_k^{\mathrm{int}}$ is a semisimple abelian category with
simple objects the integrable highest-weight modules
$L(\Lambda)$, $\Lambda$ dominant of level~$k$; there are finitely many
such (their number equals the cardinality of $P^k_+ / W$, where
$P^k_+$ is the set of level-$k$ dominant weights).  This is the
starting point for the Verlinde formula
(see Part~II for explicit computations).
\end{definition}

\begin{definition}[Weyl modules]\label{def:weyl-module}
\index{Weyl module|textbf}
\index{standard module|see{Weyl module}}
For a dominant integral weight $\Lambda$ of level~$k$, the
\emph{Weyl module} (or \emph{standard module}) is the maximal
integrable quotient of the Verma module:
\[
V(\Lambda) = \mathcal{M}(\Lambda) \Big/ \sum_{i=0}^r
\langle f_i^{\langle\Lambda, \alpha_i^\vee\rangle + 1}
\cdot v_\Lambda \rangle.
\]
At integrable level, $V(\Lambda)$ coincides with the simple module
$L(\Lambda)$ (the Weyl module is irreducible).  At non-integral or
non-dominant levels, $V(\Lambda)$ and $L(\Lambda)$ may differ.

In the bar-cobar framework, the Weyl module is the natural
``free'' object: the bar complex
$\bar{B}^{\mathrm{ch}}(\widehat{\mathfrak{g}}_k, V(\Lambda))$
computes $\mathrm{Ext}^*_{\mathcal{O}_k}(V(\Lambda), -)$, and the
bar spectral sequence degenerates at~$E_2$ precisely when
$V(\Lambda)$ is Koszul.
\end{definition}

\begin{theorem}[Kazhdan--Lusztig equivalence; \ClaimStatusProvedElsewhere]
\label{thm:kazhdan-lusztig-equivalence}
\index{Kazhdan--Lusztig equivalence|textbf}
\index{quantum group!Kazhdan--Lusztig equivalence}
Let $k$ be a non-negative integer and set
$q = e^{\pi i/(k + h^\vee)}$.  There is an equivalence of
braided tensor categories
\begin{equation}\label{eq:kl-equivalence}
\Phi_{\mathrm{KL}} \colon
\mathcal{O}_k^{\mathrm{int}}(\widehat{\mathfrak{g}})
\;\xrightarrow{\;\sim\;}
\mathcal{C}(U_q(\mathfrak{g})),
\end{equation}
where $\mathcal{C}(U_q(\mathfrak{g}))$ denotes the
\emph{semisimplified tilting category} of $U_q(\mathfrak{g})$
at the root of unity $q$.  The equivalence sends the integrable
module $L(\Lambda)$ to the simple tilting module
$T_q(\Lambda)$ and identifies the fusion product
$\boxtimes_{\widehat{\mathfrak{g}}_k}$ with the fusion tensor
product on~$\mathcal{C}$.

The equivalence sends the braiding
\textup{(}monodromy of conformal blocks on $\mathbb{P}^1$\textup{)}
to the $R$-matrix braiding on $\mathcal{C}(U_q(\mathfrak{g}))$.
\end{theorem}

\begin{proof}[References]
For $\mathfrak{g} = \mathfrak{sl}_N$, this is proved by
Kazhdan--Lusztig~\cite{KL93} via a deformation argument comparing
the monodromy representation of the KZ connection with the
quasi-$R$-matrix of $U_q(\mathfrak{g})$.
Finkelberg~\cite{Finkelberg96} extended the result to arbitrary
semisimple~$\mathfrak{g}$.  The braided tensor structure on
$\mathcal{O}_k^{\mathrm{int}}$ is constructed by
Huang--Lepowsky--Zhang~\cite{HLZ} using the theory of intertwining
operators and analytic continuation.
\end{proof}

\begin{remark}[Categorical precision]\label{rem:kl-category-precision}
The target of the KL equivalence is \emph{not} the full
representation category $\mathrm{Rep}(U_q(\mathfrak{g}))$, which
at a root of unity $q$ is non-semisimple and has infinitely many
indecomposable objects.  Rather, the correct target is the
semisimplified quotient of the tilting category: one takes
the full subcategory of tilting modules and quotients by the
tensor ideal of negligible morphisms.  The resulting modular
tensor category $\mathcal{C}(U_q(\mathfrak{g}))$ is semisimple
with $|P_+^k|$ simple objects, matching the integrable modules
at level~$k$.  For $\mathfrak{sl}_2$ at level~$k$, this gives
$k + 1$ simples on both sides.
\end{remark}

\begin{remark}[Koszul duality and $q \leftrightarrow q^{-1}$]
\label{rem:kl-koszul}
The level-shifting duality $\widehat{\mathfrak{g}}_k \leftrightarrow
\widehat{\mathfrak{g}}_{-k-2h^\vee}$
(Theorem~\ref{thm:universal-kac-moody-koszul}) transforms the KL
parameter as $q \mapsto q^{-1}$
(Theorem~\ref{thm:km-quantum-groups}).
Since $U_q(\mathfrak{g})$ and $U_{q^{-1}}(\mathfrak{g})$ are
related by the Drinfeld--Jimbo bar involution
(which reverses the $R$-matrix braiding), the KL equivalence
identifies chiral Koszul duality with the standard quantum group
duality:
\[
\begin{tikzcd}
\mathcal{O}_k^{\mathrm{int}} \arrow[r, "\Phi_{\mathrm{KL}}"]
   \arrow[d, "\mathrm{Koszul}"'] &
\mathcal{C}(U_q(\mathfrak{g}))
   \arrow[d, "q \mapsto q^{-1}"] \\
\mathcal{O}_{-k-2h^\vee}^{\mathrm{int}}
   \arrow[r, "\Phi_{\mathrm{KL}}"'] &
\mathcal{C}(U_{q^{-1}}(\mathfrak{g}))
\end{tikzcd}
\]
The diagram commutes on simple objects (both sides send
$L(\Lambda)$ to $T_{q^{-1}}(\Lambda)$); commutativity on
tensor products follows from the braiding reversal on both sides.
\end{remark}

\begin{definition}[Contragredient module]
\label{def:contragredient-module}
\index{contragredient module|textbf}
\index{restricted dual}
For $M \in \mathcal{O}_k(\widehat{\mathfrak{g}})$, the
\emph{contragredient} (restricted dual) module is
\[
M^\vee = \bigoplus_\lambda M_\lambda^*,
\]
the graded dual (direct sum of duals of weight spaces), equipped
with the $\widehat{\mathfrak{g}}$-action
$(a \cdot f)(m) = f(\sigma(a) \cdot m)$
where $\sigma$ is the Chevalley anti-involution
$\sigma(e_i) = f_i$, $\sigma(f_i) = e_i$,
$\sigma(h) = h$, $\sigma(d) = d$, $\sigma(K) = K$.

The contragredient satisfies:
\begin{enumerate}[label=\textup{(\roman*)}]
\item $\mathcal{M}(\lambda)^\vee \cong \mathcal{M}(\lambda)$ (Verma
  modules are self-contragredient);
\item $\mathcal{L}(\lambda)^\vee \cong \mathcal{L}(\lambda)$ (simple
  modules are self-contragredient);
\item the Shapovalov form (Definition~\ref{def:shapovalov-form})
  realizes the natural pairing
  $M \otimes M^\vee \to \mathbb{C}$.
\end{enumerate}
\end{definition}

\begin{remark}[Contragredient modules and Koszul duality]
\label{rem:contragredient-koszul}
For a Koszul pair $(\widehat{\mathfrak{g}}_k, \widehat{\mathfrak{g}}_{k'})$
with $k' = -k - 2h^\vee$ (Corollary~\ref{cor:level-shifting-part1}),
the contragredient functor $(-)^\vee$ on
$\mathcal{O}_k$ should \emph{not} be confused with the Koszul
duality functor $\mathrm{Ind}\colon
\mathrm{Mod}^{\mathrm{compl}}_{\widehat{\mathfrak{g}}_k} \to
\mathrm{CoMod}^{\mathrm{conil}}_{\widehat{\mathfrak{g}}_{k'}}$
(Theorem~\ref{thm:e1-module-koszul-duality}).
The contragredient operates \emph{within} the category at a fixed level;
Koszul duality exchanges categories at \emph{different} levels.
Both preserve the set of highest weight labels (the dominant weights
are the same for $\mathfrak{g}$ regardless of level), but they act on
different structures.
\end{remark}

\begin{theorem}[BGG reciprocity for affine algebras; \ClaimStatusProvedElsewhere]
\label{thm:bgg-reciprocity}
\index{BGG reciprocity|textbf}
In category $\mathcal{O}_k(\widehat{\mathfrak{g}})$ at generic level,
each simple module $\mathcal{L}(\lambda)$ has a projective cover
$P(\lambda)$ satisfying the \emph{BGG reciprocity}:
\begin{equation}\label{eq:bgg-reciprocity}
[P(\lambda) : \mathcal{M}(\mu)]
= [\mathcal{M}(\mu) : \mathcal{L}(\lambda)]
\end{equation}
where $[P : \mathcal{M}]$ denotes the Verma module multiplicity in
a Verma flag of the projective, and
$[\mathcal{M} : \mathcal{L}]$ denotes the composition factor
multiplicity.

Equivalently, the Grothendieck group $K_0(\mathcal{O}_k)$ has two
natural bases: the Verma modules $\{[\mathcal{M}(\lambda)]\}$ and the
simple modules $\{[\mathcal{L}(\lambda)]\}$.  The transition matrix
between them is upper-triangular (with respect to the partial order
on weights), and BGG reciprocity states that the inverse transition
matrix is the \emph{transpose} of the transition matrix.

This is proved for finite-dimensional semisimple Lie algebras by
Bernstein--Gel'fand--Gel'fand~\cite{BGG76} and extended to
symmetrizable Kac--Moody algebras by Rocha-Caridi--Wallach and
Kashiwara--Tanisaki~\cite{KT95}.
\end{theorem}

\begin{remark}[BGG reciprocity and the bar-cobar resolution]
\label{rem:bgg-bar-cobar}
\index{BGG resolution!Koszul interpretation}
In the language of this monograph, BGG reciprocity has a transparent
homological interpretation.  The projective cover $P(\lambda)$
is the bar resolution of $\mathcal{L}(\lambda)$ by Verma modules
(Theorem~\ref{thm:weyl-kac-geometric}, Step~1).  The multiplicity
$[P(\lambda) : \mathcal{M}(\mu)]$ counts the number of times
$\mathcal{M}(\mu)$ appears in the bar complex.  The composition
multiplicity $[\mathcal{M}(\mu) : \mathcal{L}(\lambda)]$ measures
extensions between simples.

For a \emph{Koszul} algebra, the Koszul resolution is minimal, so the
bar complex multiplicities are determined by the Ext groups:
$[P(\lambda) : \mathcal{M}(\mu)] =
\dim \mathrm{Ext}^*_{\mathcal{O}}(\mathcal{L}(\lambda),
\mathcal{L}(\mu))$.
BGG reciprocity then becomes the statement that the bar complex is
self-dual under the Shapovalov form---a consequence of the
bar-cobar duality (Theorem~\ref{thm:bar-cobar-isomorphism-main})
combined with the self-contragredience of Verma and simple modules.
\end{remark}

\begin{remark}[BGG reciprocity from bar spectral sequence]\label{rem:bgg-bar-spectral}
\index{BGG resolution!spectral sequence interpretation}
The bar spectral sequence provides a refinement of BGG reciprocity.  The $E_2$ page $E_2^{p,q} = H^q(\bar{B}^p(\mathcal{A}))$ of the bar filtration encodes the module multiplicities in each bar degree: $E_2^{p,0}$ is the space of weight-$p$ bar generators, and $E_2^{p,q}$ for $q > 0$ measures syzygies.  At admissible levels, the rationality of $L_k(\mathfrak{g})$ (Theorem~\ref{thm:admissible-rep-theory}) implies finite-dimensionality of all $E_r^{p,q}$ (Corollary~\ref{cor:bar-admissible-finiteness}).  A deeper question is whether the BGG multiplicities---specifically the Kazhdan--Lusztig polynomials $P_{w,w'}(q) = \sum_i q^i \dim \mathrm{Ext}^i_{\mathcal{O}}(\mathcal{M}(w \cdot \lambda), \mathcal{L}(w' \cdot \lambda))$---can be read off from the $E_2$ page of the bar spectral sequence.  This would require a wall-crossing analysis for the spectral sequence as the level crosses admissible values, which remains open.
\end{remark}

\begin{definition}[Tilting modules]\label{def:tilting-module}
\index{tilting module|textbf}
A module $T \in \mathcal{O}_k(\widehat{\mathfrak{g}})$ is
\emph{tilting} if both $T$ and $T^\vee$ (the contragredient)
admit a Verma flag:
\[
0 = T_0 \subset T_1 \subset \cdots \subset T_r = T
\quad\text{with}\quad
T_i/T_{i-1} \cong \mathcal{M}(\mu_i)
\text{ for each } i.
\]
For each dominant weight~$\lambda$, there exists a unique
indecomposable tilting module $T(\lambda)$ with highest
weight~$\lambda$.  At integrable level~$k$, the indecomposable
tilting modules $\{T(\Lambda) : \Lambda \in P_+^k\}$ form a
complete set of projective objects in the fusion category
$\mathcal{O}_k^{\mathrm{int}}(\widehat{\mathfrak{g}})$
(Andersen~\cite{Andersen92}, Soergel~\cite{Soergel97}).
\end{definition}

\begin{proposition}[Tilting modules and the bar complex; \ClaimStatusProvedHere]
\label{prop:tilting-bar}
\index{tilting module!bar complex}
\index{bar complex!tilting modules}
For an indecomposable tilting module $T(\lambda)$ over
$\widehat{\mathfrak{g}}_k$ at generic level:
\begin{enumerate}[label=\textup{(\roman*)}]
\item The bar complex $\bar{B}^{\mathrm{ch}}(\widehat{\mathfrak{g}}_k,
  T(\lambda))$ is \emph{self-dual} under the Verdier involution:
  $\bar{B}(T(\lambda))^\vee \cong \bar{B}(T(\lambda))$.
\item The Verma flag multiplicities are:
  $[T(\lambda) : \mathcal{M}(\mu)]
  = \dim \mathrm{Hom}_{\mathcal{O}}(P(\mu), T(\lambda))
  = [\mathcal{M}(\mu) : \mathcal{L}(\lambda)]$
  where the last equality is BGG reciprocity
  \textup{(}Theorem~\textup{\ref{thm:bgg-reciprocity})}.
\item Under the Kazhdan--Lusztig equivalence
  $\Phi_{\mathrm{KL}}\colon
  \mathcal{O}_k^{\mathrm{int}} \xrightarrow{\sim}
  \mathrm{Rep}(U_q(\mathfrak{g}))$,
  the tilting module $T(\Lambda)$ maps to the indecomposable
  tilting module $T_q(\Lambda) \in \mathrm{Rep}(U_q(\mathfrak{g}))$.
\item Under the module Koszul duality $\Phi$ of
  Theorem~\textup{\ref{thm:monoidal-module-koszul}},
  tilting modules are sent to tilting modules:
  $\Phi(T_k(\Lambda)) = T_{k'}(\Lambda)$.
\end{enumerate}
\end{proposition}

\begin{proof}
(i) follows from the self-contragredience
$T(\lambda)^\vee \cong T(\lambda)$ (the defining property of tilting
modules), combined with the compatibility of the bar construction
with contragredience.

(ii) is a direct consequence of BGG reciprocity: a projective
$P(\mu)$ has a Verma flag with multiplicities
$[P(\mu) : \mathcal{M}(\nu)] = [\mathcal{M}(\nu) : \mathcal{L}(\mu)]$,
and $T(\lambda)$, being both Verma-filtered and co-Verma-filtered, has
multiplicities equal to the composition multiplicities by the
double-flag property.

(iii) follows from the uniqueness of indecomposable tilting modules
with a given highest weight: the KL functor preserves highest
weights and the Verma flag structure (it sends
$\mathcal{M}(\Lambda)$ to the quantum Verma module
$M_q(\Lambda)$), so it must send tilting to tilting.

(iv) The Koszul duality functor preserves the Verma flag structure
(since Verma modules are Koszul, the bar complex of a Verma module is
a comodule with the correct filtration).  The self-duality of tilting
modules (both Verma-filtered and co-Verma-filtered) maps to the
corresponding property at the dual level.  By uniqueness of
indecomposable tilting modules, $\Phi(T_k(\Lambda)) = T_{k'}(\Lambda)$.
\end{proof}

\begin{proposition}[Verma module bar complex; \ClaimStatusProvedHere]
\label{prop:verma-bar-complex}
\index{Verma module!bar complex}
Let $\mathcal{M}(\lambda)$ be the Verma module over
$\widehat{\mathfrak{g}}_k$ at generic level~$k$.  The geometric bar
complex \textup{(}Definition~\textup{\ref{def:bar-complex-modules})}
takes the form:
\[
\bar{B}_n^{\mathrm{ch}}(\widehat{\mathfrak{g}}_k, \mathcal{M}(\lambda))
= \Gamma\bigl(\overline{C}_{n+2}(X),\;
j_*j^*(\widehat{\mathfrak{g}}_k \boxtimes
\overline{\widehat{\mathfrak{g}}}_k^{\boxtimes n} \boxtimes
\mathcal{M}(\lambda)) \otimes \Omega^n_{\log}\bigr).
\]
This complex is acyclic at generic level
\textup{(}Proposition~\textup{\ref{prop:generic-irreducibility}}:
Verma modules are irreducible, so the bar complex has no higher
cohomology\textup{)}.  At special levels where the Verma module is
reducible, $H^1(\bar{B}_\bullet(\widehat{\mathfrak{g}}_k,
\mathcal{M}(\lambda)))$ is generated by the singular vectors of
$\mathcal{M}(\lambda)$: each singular vector
$v_\mu \in \mathcal{M}(\lambda)_\mu$ satisfying
$\widehat{\mathfrak{g}}_{>0} \cdot v_\mu = 0$ contributes a class
in $H^1$ at conformal weight~$|\mu|$.
\end{proposition}

\begin{proof}
At generic level, $\mathcal{M}(\lambda)$ is irreducible
(Proposition~\ref{prop:generic-irreducibility}), so it is a free
module over $U(\widehat{\mathfrak{g}}_{<0})$.  The bar complex is
then acyclic by the standard contracting homotopy
(\S\ref{sec:chiral-module-resolution} below).

At special level, the maximal submodule
$\mathcal{J}(\lambda) \subset \mathcal{M}(\lambda)$ is non-zero.
Each inclusion $\mathcal{M}(\mu) \hookrightarrow \mathcal{M}(\lambda)$
of a Verma submodule (generated by a singular vector of weight~$\mu$)
produces a class in $H^1$ of the bar complex: the inclusion
$\iota: v_\mu \hookrightarrow \mathcal{M}(\lambda)$ satisfies
$d_{\mathrm{bar}}(\iota) = 0$ (since $v_\mu$ is annihilated by
$\widehat{\mathfrak{g}}_{>0}$) and is non-trivial in cohomology
(since $v_\mu$ is not in the image of $d_{\mathrm{bar}}$ from
$\bar{B}_2$: this would require $v_\mu$ to factor through the
chiral product, which contradicts primitivity of singular vectors).

The vanishing locus of the Shapovalov form
(Definition~\ref{def:shapovalov-form}) determines which weight
spaces have singular vectors, and the Kac determinant formula
$\det \langle -, - \rangle_{\lambda, n}$ gives the multiplicity.
In the bar-cobar framework, the Kac determinant is the determinant
of the pairing between $\bar{B}_1$ and $\bar{B}^1$ at conformal
weight~$n$: its vanishing detects the failure of bar-complex
acyclicity, which is the bar-cobar translation of Verma module
reducibility.
\end{proof}

\begin{definition}[Kazhdan--Lusztig polynomials]
\label{def:kl-polynomials}
\index{Kazhdan--Lusztig polynomial|textbf}
The \emph{Kazhdan--Lusztig polynomials}
$P_{y,w}(q) \in \mathbb{Z}[q]$, indexed by pairs of elements
$y, w$ in the (affine) Weyl group with $y \leq w$ in the Bruhat
order, are defined by the recursive formula:
\[
P_{y,w}(q) = q^{(\ell(w) - \ell(y) - 1)/2}\,
P_{y,sw}(q^{-1})
+ q^{(\ell(w) - \ell(y))/2}\, P_{sy,sw}(q^{-1})
- \sum_{y \leq z < sw} \mu(z, sw)\, q^{(\ell(w) - \ell(z))/2}\, P_{y,z}(q)
\]
with initial conditions $P_{w,w} = 1$ and $P_{y,w} = 0$ for
$y \not\leq w$, where $s$ is a simple reflection with $sw < w$
and $\mu(z,w) = \deg P_{z,w}$ is the leading coefficient
(see Kazhdan--Lusztig~\cite{KL79}).

The representation-theoretic significance: for category
$\mathcal{O}$ of a symmetrizable Kac--Moody algebra,
the composition multiplicities satisfy
\begin{equation}\label{eq:kl-multiplicity}
[\mathcal{M}(w \cdot \lambda) : \mathcal{L}(y \cdot \lambda)]
= P_{y,w}(1)
\end{equation}
as conjectured by Kazhdan--Lusztig~\cite{KL79} and proved by
Beilinson--Bernstein~\cite{BB81} and Brylinski--Kashiwara~\cite{BK81}
(for finite type) and Kashiwara--Tanisaki~\cite{KT95} (for affine type).
\end{definition}

\begin{remark}[Kazhdan--Lusztig polynomials and the bar spectral sequence]
\label{rem:kl-bar}
\index{Kazhdan--Lusztig polynomial!bar complex interpretation}
The KL multiplicity formula~\eqref{eq:kl-multiplicity} has a direct
interpretation in the bar complex framework.  The Koszul resolution
of $\mathcal{L}(\lambda)$ (Theorem~\ref{thm:koszul-resolution-module})
is the BGG resolution, whose $n$-th term is
$\bigoplus_{\ell(w) = n} \mathcal{M}(w \cdot \lambda)$.
The characters of the simple modules are then determined by inverting
the Verma character matrix, and the inversion coefficients are the
values $P_{y,w}(1)$ of the KL polynomials at $q = 1$.

For Koszul algebras, the full $q$-polynomial $P_{y,w}(q)$ (not just
its value at $1$) is recovered from the $E_2$~page of the bar
spectral sequence: $P_{y,w}(q) = \sum_i \dim
\mathrm{Ext}^i_{\mathcal{O}}(\mathcal{M}(w \cdot \lambda),
\mathcal{L}(y \cdot \lambda))\, q^{i/2}$.
This is the Koszul--BGG correspondence between graded Ext groups and
KL polynomials (Beilinson--Ginzburg--Soergel~\cite{BGS96}).
\end{remark}

\begin{example}[KL polynomials for $\widehat{\mathfrak{sl}}_2$]
\label{ex:kl-sl2}
\index{Kazhdan--Lusztig polynomial!sl2@$\mathfrak{sl}_2$ computation}
The affine Weyl group of $\widehat{\mathfrak{sl}}_2$ is
$\mathcal{W} = \mathbb{Z} \rtimes \mathbb{Z}/2$, generated by
reflections $s_0$ (affine) and $s_1$ (finite).  Elements of minimal
length in each coset $\mathcal{W}/\langle s_1 \rangle$ are
parameterized by $n \in \mathbb{Z}_{\geq 0}$ with $\ell(w_n) = n$,
where $w_n = s_0 s_1 s_0 \cdots$ ($n$~factors).  The Bruhat order is
the linear order $w_0 < w_1 < w_2 < \cdots$.

For this linear Bruhat order, the KL polynomials are:
\[
P_{w_m, w_n}(q) = \begin{cases}
1 & \text{if } m = n, \\
0 & \text{if } m > n, \\
1 & \text{if } n - m = 1 \text{ (adjacent)}, \\
\displaystyle\sum_{j=0}^{\lfloor(n-m-1)/2\rfloor} q^j
& \text{in general.}
\end{cases}
\]
In particular, the polynomials are truncated geometric series of
degree $\lfloor(n-m-1)/2\rfloor$, and the evaluation
$P_{w_m, w_n}(1) = \lceil(n-m)/2\rceil$ counts the total
composition factor multiplicity.

\emph{Bar spectral sequence interpretation.}
By Remark~\ref{rem:kl-bar}, $P_{w_m, w_n}(q) = \sum_i \dim
\mathrm{Ext}^i_{\mathcal{O}}(\mathcal{M}(w_n \cdot \lambda),
\mathcal{L}(w_m \cdot \lambda))\, q^{i/2}$.  The formula
$P_{w_m, w_n}(q) = \sum_{j=0}^{\lfloor(n-m-1)/2\rfloor} q^j$ says
that the Ext groups between Verma and simple modules grow linearly
with the Weyl group distance, with non-trivial contributions in even
cohomological degrees.  On the $E_2$~page of the bar spectral
sequence, these Ext groups appear as non-trivial entries that
obstruct degeneration at $E_2$ when the level is not generic.
\end{example}

\subsection{Zhu's algebra and module classification}

\begin{definition}[Zhu's algebra]\label{def:zhu-algebra}
\index{Zhu's algebra|textbf}
Let $V$ be a vertex algebra with conformal vector.
Define $O(V) \subset V$ to be the subspace spanned by elements
\[
a \circ b = \mathrm{Res}_{z=0}\, Y(a,z) \frac{(1+z)^{\mathrm{wt}(a)}}{z^2}\, b,
\]
and define a bilinear product $*$ on~$V$ by
\[
a * b = \mathrm{Res}_{z=0}\, Y(a,z) \frac{(1+z)^{\mathrm{wt}(a)}}{z}\, b.
\]
\emph{Zhu's algebra} is the quotient
$A(V) = V / O(V)$,
equipped with the product induced by~$*$.  This is a
unital associative algebra (Zhu~\cite{Zhu}).
\end{definition}

\begin{theorem}[Zhu's correspondence; \ClaimStatusProvedElsewhere]
\label{thm:zhu-correspondence}
\index{Zhu's algebra!correspondence}
Let $V$ be a vertex algebra of countable dimension satisfying Zhu's
$C_2$-cofiniteness condition.  Then:
\begin{enumerate}[label=(\roman*)]
\item The functor $\Omega\colon \mathcal{M} \mapsto \mathcal{M}_0$
   (taking the lowest conformal weight space) is a bijection
   \[
   \{\text{simple positive-energy } V\text{-modules}\}
   \;\xrightarrow{\;\sim\;}
   \{\text{simple } A(V)\text{-modules}\}.
   \]
\item Under this bijection, a simple $V$-module $\mathcal{L}(\lambda)$
   corresponds to the simple $A(V)$-module with the same lowest
   weight space.
\item If $A(V)$ is finite-dimensional, then $V$ has finitely many
   simple positive-energy modules (rationality in the sense
   of Zhu~\cite{Zhu}).
\end{enumerate}
\end{theorem}

\begin{example}[Zhu's algebra for standard examples]\label{ex:zhu-algebras}
\index{Zhu's algebra!examples}
\leavevmode
\begin{enumerate}[label=(\roman*)]
\item $A(\widehat{\mathfrak{g}}_k) \cong U(\mathfrak{g})$ for all
   levels~$k$ (Frenkel--Zhu~\cite{FZ92}).  Simple $\widehat{\mathfrak{g}}_k$-modules in
   category~$\mathcal{O}$
   (Definition~\ref{def:category-O}) correspond to simple
   $U(\mathfrak{g})$-modules (i.e., finite-dimensional
   $\mathfrak{g}$-representations for the integrable subcategory).
\item $A(\mathrm{Vir}_c) \cong \mathbb{C}[x]$ \cite{FZ92} (polynomial ring in one
   variable, with $x = [L_{-2}\mathbf{1}]$).  Simple modules are
   parameterized by $h \in \mathbb{C}$ (the $L_0$-eigenvalue on the
   highest weight vector).
\item $A(\mathcal{W}^k(\mathfrak{sl}_2)) = A(\mathrm{Vir}_{c(k)}) \cong
   \mathbb{C}[x]$, since the Drinfeld--Sokolov reduction at level~$k$
   produces a Virasoro algebra with $c = 1 - 6(k+1)^2/(k+2)$.
\item For the Heisenberg algebra $\mathcal{H}_\kappa$,
   $A(\mathcal{H}_\kappa) \cong \mathbb{C}[\alpha]$, and simple
   modules (Fock modules $F_\lambda$) are parameterized by
   $\lambda \in \mathbb{C}$.
\end{enumerate}
\end{example}

\begin{proposition}[Zhu algebra under level-shifting Koszul duality; \ClaimStatusProvedHere]
\label{prop:zhu-koszul-compatibility}
\index{Zhu's algebra!Koszul compatibility}
For the following Koszul pairs, the Zhu algebras are canonically
isomorphic:
\begin{enumerate}[label=\textup{(\roman*)}]
\item \emph{Affine Kac--Moody:}
  $A(\widehat{\mathfrak{g}}_k) \cong U(\mathfrak{g})
  \cong A(\widehat{\mathfrak{g}}_{-k-2h^\vee})$
  for all generic levels~$k$;
\item \emph{Heisenberg:}
  $A(\mathcal{H}_\kappa) \cong \mathbb{C}[\alpha]
  \cong A(\mathrm{Sym}^{\mathrm{ch}}(V^*))$
  for all $\kappa \neq 0$.
\end{enumerate}
In both cases, the Koszul pairing at conformal weight~$1$ descends to
the identity on the Zhu algebra.  Consequently, the simple module
labels---dominant weights of $\mathfrak{g}$ in case~\textup{(i)},
charges $\lambda \in \mathbb{C}$ in case~\textup{(ii)}---are
preserved by module Koszul duality
\textup{(}Theorem~\textup{\ref{thm:e1-module-koszul-duality})}.
\end{proposition}

\begin{proof}
In both cases, the Zhu algebra is determined by the zero-mode
commutator $[\bar{a}, \bar{b}] = \overline{a_{(0)}b}$, which
encodes the Lie bracket on the weight-$1$ generators.

\emph{(i) Affine:}
By Frenkel--Zhu~\cite{FZ92},
$A(\widehat{\mathfrak{g}}_k) \cong U(\mathfrak{g})$ via
$[J^a_{-1}\mathbf{1}] \mapsto \bar{J}^a$, with commutator
$[\bar{J}^a, \bar{J}^b] = f_{ab}^c\, \bar{J}^c$.  The level~$k$
enters the OPE only through the double pole
$J^a_{(1)}J^b = k(a,b)$, which maps to a scalar
$k(a,b) \in \mathbb{C} \subset A(\widehat{\mathfrak{g}}_k)$.
Since the Koszul dual $\widehat{\mathfrak{g}}_{-k-2h^\vee}$
has the same generators and the same structure constants~$f_{ab}^c$
(only the level changes), both Zhu algebras equal~$U(\mathfrak{g})$.

\emph{(ii) Heisenberg:}
$A(\mathcal{H}_\kappa) = \mathbb{C}[J_0]$ because the Heisenberg
OPE $J(z)J(w) \sim \kappa/(z-w)^2$ has no simple pole:
$J_{(0)}J = 0$, so $[\bar{J}, \bar{J}] = 0$ and $A(\mathcal{H}_\kappa)$
is commutative.  Similarly,
$A(\mathrm{Sym}^{\mathrm{ch}}(V^*)) = \mathbb{C}[\alpha_0]$
(the commutative chiral algebra has trivial OPE).  The Koszul pairing
$\langle J, \alpha \rangle = 1$ identifies $J_0$ with $\alpha_0$.
\end{proof}

\begin{remark}[Scope and consequences]\label{rem:zhu-koszul-scope}
\leavevmode
\begin{enumerate}[label=(\roman*)]
\item \emph{Module labels preserved.}
  The isomorphism $A(\mathcal{A}) \cong A(\mathcal{A}^!)$ means that
  simple modules on both sides of the Koszul duality are indexed by
  the same set (simple $A(\mathcal{A})$-modules).  For
  $\widehat{\mathfrak{g}}_k$, both $\mathcal{O}_k$ and
  $\mathcal{O}_{-k-2h^\vee}$ have simples labeled by dominant weights
  of~$\mathfrak{g}$, though the integrability conditions differ
  (level-$k$ dominant vs.\ level-$(-k{-}2h^\vee)$-dominant).
\item \emph{$\mathcal{W}$-algebras.}
  The argument above requires that all generators have conformal
  weight~$1$.  For $\mathcal{W}$-algebras
  $\mathcal{W}^k(\mathfrak{g})$ (generators at weights
  $d_1, \ldots, d_r$ determined by the exponents of~$\mathfrak{g}$),
  a different argument is needed: the Zhu algebra is identified
  with~$Z(U(\mathfrak{g}))$ via DS reduction (Theorem~\ref{thm:w-algebra-zhu-koszul}),
  which is again level-independent.
\end{enumerate}
\end{remark}

\begin{corollary}[Virasoro Zhu algebra is Koszul-invariant; \ClaimStatusProvedHere]
\label{cor:virasoro-zhu-koszul}
\index{Zhu's algebra!Virasoro}
Despite the conformal weight~$2$ generator, the Virasoro Koszul pair
$(\mathrm{Vir}_c, \mathrm{Vir}_{26-c})$ satisfies
\[
A(\mathrm{Vir}_c) \;\cong\; \mathbb{C}[x]
\;\cong\; A(\mathrm{Vir}_{26-c})
\]
for all $c \in \mathbb{C}$, where $x = [L_{-2}\mathbf{1}]$.
\end{corollary}

\begin{proof}
By Frenkel--Zhu~\cite{FZ92},
$A(\mathrm{Vir}_c) \cong \mathbb{C}[x]$ with
$x = [L_{-2}\mathbf{1}]$ and the product
$[L_{-2}\mathbf{1}] * [L_{-2}\mathbf{1}]
= [L_{-2}^2\mathbf{1}] + [L_{-3}\mathbf{1}]$.
The central charge~$c$ enters the vertex algebra through the
relation $L_{(3)}L = c/2$, which maps to a scalar in
$A(\mathrm{Vir}_c)$ and does not alter the algebra structure.
Therefore $A(\mathrm{Vir}_c) \cong \mathbb{C}[x]$
independently of~$c$, and the isomorphism
$A(\mathrm{Vir}_c) \cong A(\mathrm{Vir}_{26-c})$ follows.
\end{proof}

\begin{theorem}[$\mathcal{W}$-algebra Zhu algebras are Koszul-invariant; \ClaimStatusProvedHere]
\label{thm:w-algebra-zhu-koszul}
\index{Zhu's algebra!W-algebra@$\mathcal{W}$-algebra}
\index{Zhu's algebra!Koszul invariance}
For any simple Lie algebra~$\mathfrak{g}$ of rank~$r$ and generic level~$k$,
the Zhu algebra of the principal $\mathcal{W}$-algebra is
\begin{equation}\label{eq:w-zhu-center}
A(\mathcal{W}^k(\mathfrak{g}))
  \;\cong\; Z(U(\mathfrak{g}))
  \;\cong\; \mathbb{C}[p_1, \ldots, p_r],
\end{equation}
where $Z(U(\mathfrak{g}))$ is the center of the universal enveloping
algebra and $p_1, \ldots, p_r$ are fundamental Casimir elements of
degrees $d_1 + 1, \ldots, d_r + 1$ \textup{(}$d_i$ the exponents
of~$\mathfrak{g}$\textup{)}.  In particular, this is level-independent:
\[
A(\mathcal{W}^k(\mathfrak{g}))
  \;\cong\; A(\mathcal{W}^{k'}(\mathfrak{g}))
  \qquad \text{for } k' = -k - 2h^\vee.
\]
\end{theorem}

\begin{proof}
The result follows from the composition of four published theorems.

\emph{Step~1 (Frenkel--Zhu).}
By~\cite{FZ92}, the Zhu algebra of the universal affine vertex algebra at
generic level is $A(V_k(\mathfrak{g})) \cong U(\mathfrak{g})$.

\emph{Step~2 (Feigin--Frenkel, Kac--Roan--Wakimoto).}
The principal $\mathcal{W}$-algebra is obtained by quantum Drinfeld--Sokolov
reduction: $\mathcal{W}^k(\mathfrak{g}) = H^0_{\mathrm{DS}}(V_k(\mathfrak{g}))$.

\emph{Step~3 (Arakawa).}
By~\cite{Ara07}, the Zhu functor commutes with quantum DS reduction:
$A(H^0_{\mathrm{DS}}(V)) = H^0_{\mathrm{DS}}(A(V))$.
Applying this to $V = V_k(\mathfrak{g})$ gives
$A(\mathcal{W}^k(\mathfrak{g})) = H^0_{\mathrm{DS}}(U(\mathfrak{g}))$.

\emph{Step~4 (Kostant, Harish-Chandra).}
Kostant's theorem on Whittaker vectors~\cite{Kostant78} identifies
the DS reduction of the enveloping algebra with respect to the
principal nilpotent character~$\chi$ as
$H^0_{\mathrm{DS}}(U(\mathfrak{g}), \chi) \cong Z(U(\mathfrak{g}))$.
By the Harish-Chandra isomorphism~\cite{HarishChandra},
$Z(U(\mathfrak{g})) \cong \mathbb{C}[p_1, \ldots, p_r]$.

The composition gives~\eqref{eq:w-zhu-center}.  Since
$Z(U(\mathfrak{g}))$ depends only on~$\mathfrak{g}$ (not on~$k$), the
isomorphism $A(\mathcal{W}^k(\mathfrak{g})) \cong A(\mathcal{W}^{k'}(\mathfrak{g}))$
at dual levels follows immediately.
\end{proof}

\begin{remark}[Explicit cases and admissible levels]\label{rem:zhu-w-algebras-explicit}
\leavevmode
\begin{enumerate}[label=(\roman*)]
\item \emph{Virasoro.}
  Theorem~\ref{thm:w-algebra-zhu-koszul} for
  $\mathfrak{g} = \mathfrak{sl}_2$ ($r = 1$, exponent~$1$, one
  Casimir of degree~$2$) recovers
  Corollary~\ref{cor:virasoro-zhu-koszul}:
  $A(\mathrm{Vir}_c) \cong \mathbb{C}[C_2] \cong \mathbb{C}[x]$.
\item \emph{$\mathcal{W}_3$.}
  For $\mathfrak{g} = \mathfrak{sl}_3$ ($r = 2$, exponents $1, 2$),
  the Zhu algebra is
  $A(\mathcal{W}^k(\mathfrak{sl}_3)) \cong \mathbb{C}[C_2, C_3]$,
  where $C_2, C_3$ are the quadratic and cubic Casimirs.
  The generators $x = [T]$ and $y = [W]$ map to (level-dependent
  normalizations of) $C_2$ and $C_3$ respectively, but the
  algebra structure is the free polynomial ring~$\mathbb{C}[x, y]$
  regardless of~$k$.
\item \emph{Admissible levels.}
  At admissible levels $k = -h^\vee + p/q$ (with
  $(p,q) = 1$), the simple quotient $L_k(\mathcal{W}^k(\mathfrak{g}))$
  has a \emph{finite-dimensional} Zhu algebra, which is a proper
  quotient of $Z(U(\mathfrak{g}))$ determined by the set of
  admissible highest weights.  These finite-dimensional Zhu algebras
  are generally \emph{not} isomorphic at dual admissible levels (the
  number of simple modules may differ), so Koszul invariance
  of the Zhu algebra holds at generic levels but not in general
  at special levels.
\end{enumerate}
\end{remark}

\subsection{Associated variety and rationality}

\begin{definition}[Li's filtration and the associated variety]
\label{def:associated-variety}
\index{associated variety|textbf}
\index{C2-cofiniteness@$C_2$-cofiniteness}
\index{Li filtration}
Let $V$ be a vertex algebra with conformal vector.  Define the
\emph{Li filtration} $F_p V = \mathrm{span}\{a_{(-n-2)} b \mid
a \in V,\, b \in F_{p-1} V,\, n \geq 0\}$ with $F_0 V = V$.
The associated graded $\mathrm{gr}^F V = \bigoplus_{p \geq 0}
F_p V / F_{p+1} V$ inherits a commutative Poisson algebra structure
(Zhu~\cite{Zhu} and Li): the product and Lie bracket are induced by
\begin{align}
\bar{a} \cdot \bar{b} &= \overline{a_{(-1)} b}, &
\{\bar{a}, \bar{b}\} &= \overline{a_{(0)} b}.
\end{align}
The \emph{associated variety} of $V$ is the Poisson variety
\[
X_V = \mathrm{Specm}\,(\mathrm{gr}^F V / \sqrt{0}) \subset
\mathfrak{g}^* \quad
\text{(when $V$ is generated by weight-$1$ fields)}.
\]
More generally, $X_V$ is a subvariety of the jet scheme
$J_\infty(\mathrm{Specm}\, A(V))$, where $A(V)$ is the Zhu algebra.
\end{definition}

\begin{theorem}[Arakawa's rationality criterion; \ClaimStatusProvedElsewhere]
\label{thm:arakawa-rationality}
\index{Arakawa!rationality criterion}
\index{associated variety!rationality criterion}
Let $V$ be a vertex algebra satisfying Zhu's $C_2$-cofiniteness
condition \textup{(}$\dim V / C_2(V) < \infty$, where
$C_2(V) = \mathrm{span}\{a_{(-2)} b\}$\textup{)}.  Then:
\begin{enumerate}[label=\textup{(\roman*)}]
\item $X_V = \{0\}$ if and only if $V$ is $C_2$-cofinite.
\item If $V$ is $C_2$-cofinite, then $V$ has finitely many simple
  modules \textup{(}Zhu~\cite{Zhu96}\textup{)}.
  \emph{Rationality}---that every grading-restricted module is
  completely reducible---is a strictly stronger condition; the
  triplet algebras $\mathcal{W}(p)$ are $C_2$-cofinite but
  \emph{not} rational
  \textup{(}see Example~\ref{ex:triplet-logarithmic}\textup{)}.
\item For the affine vertex algebra
  $V_k(\mathfrak{g})$ at admissible level~$k$, the associated variety
  is
  \[
  X_{V_k(\mathfrak{g})} = \overline{\mathbb{O}}
  \]
  where $\overline{\mathbb{O}}$ is the closure of a nilpotent
  orbit in $\mathfrak{g}^*$ determined by~$k$
  \textup{(}Arakawa~\cite{Arakawa17}\textup{)}.
  In particular, $V_k(\mathfrak{g})$ is rational if and only if
  $X_{V_k(\mathfrak{g})} = \{0\}$, which occurs precisely at the
  non-degenerate admissible levels.
\end{enumerate}
\end{theorem}

\begin{remark}[Associated variety and Koszul duality]
\label{rem:associated-variety-koszul}
\index{associated variety!Koszul duality}
The associated variety provides a geometric invariant that is
\emph{not} preserved by level-shifting Koszul duality.
For $\widehat{\mathfrak{sl}}_2$ at level $k = 1$ (integrable,
rational): $X_{V_1(\mathfrak{sl}_2)} = \{0\}$.  The Koszul dual
at level $k' = -3$ (critical minus 1): $X_{V_{-3}(\mathfrak{sl}_2)} =
\mathfrak{sl}_2^*$ (the full nilpotent cone, since $-3$ is below the
critical level).

More generally, Koszul duality interchanges levels
$k \leftrightarrow -k - 2h^\vee$, which transforms the admissible
level condition $k = -h^\vee + p/q$ into $k' = -h^\vee - p/q$.
This can map rational algebras ($X_V = \{0\}$) to
non-rational ones ($X_V \neq \{0\}$) and vice versa.

In the bar-cobar framework, the associated variety controls
convergence: when $X_V = \{0\}$, the bar spectral sequence
(Theorem~\ref{thm:deformation-acyclicity}) degenerates at a
finite page, while for $X_V \neq \{0\}$ the spectral sequence
may have non-trivial higher differentials.  This is the
geometric counterpart of the passage from Koszul (generic level)
to curved $A_\infty$ (special level) in the bar complex.
\end{remark}

\begin{proposition}[Orbit duality for affine vertex algebras; \ClaimStatusProvedHere]
\label{prop:orbit-duality}
\index{Barbasch--Vogan duality!orbit duality}
\index{associated variety!orbit duality}
\index{nilpotent orbit!Koszul duality}
Let $\mathfrak{g}$ be a simple Lie algebra of type~$A$
\textup{(}i.e., $\mathfrak{g} = \mathfrak{sl}_N$\textup{)}, and let
$V_k(\mathfrak{g})$ be the simple affine vertex algebra at admissible
level~$k$.  Then the associated varieties of $V_k(\mathfrak{g})$ and
its Koszul dual $V_{k'}(\mathfrak{g})$ ($k' = -k - 2h^\vee$)
are related by the \emph{Barbasch--Vogan duality} $d\colon
\mathcal{N}(\mathfrak{g}) \to \mathcal{N}(\mathfrak{g})$:
\begin{equation}\label{eq:orbit-duality}
X_{V_{k'}(\mathfrak{g})} = \overline{d(X_{V_k(\mathfrak{g})})}
\end{equation}
where $d(\overline{\mathbb{O}}_\lambda) =
\overline{\mathbb{O}}_{\lambda^t}$ is the orbit corresponding
to the transpose partition \textup{(}for type~$A$, this equals
the Barbasch--Vogan duality of~\cite{BV85}\textup{)}.
\end{proposition}

\begin{proof}
At admissible level $k = -h^\vee + p/q$ with $\gcd(p,q) = 1$ and
$p \geq h^\vee$, Arakawa~\cite{Arakawa17} shows that
$X_{V_k(\mathfrak{g})} = \overline{\mathbb{O}}_\lambda$ where
$\lambda$ is determined by the denominator~$q$ of the admissible
level.  For $\mathfrak{g} = \mathfrak{sl}_N$: at non-degenerate
admissible level, $X_{V_k} = \{0\}$ (trivial orbit, partition
$(1^N)$); at degenerate admissible level,
$X_{V_k} = \overline{\mathbb{O}}_\lambda$ for a non-trivial
partition~$\lambda$.

The Koszul dual level $k' = -k - 2h^\vee = -h^\vee - p/q$
lies outside the admissible regime
(Definition~\ref{def:admissible-level}), but the vertex algebra
$V_{k'}(\mathfrak{g})$ is well-defined.
The associated variety $X_{V_{k'}}$ is determined by the
Feigin--Frenkel involution: the shift $p/q \mapsto -p/q$
transposes the partition classifying the nilpotent orbit
(by the interplay between the dual Coxeter number and the
admissible weight classification;
cf.~Kac--Wakimoto~\cite{KW88}), giving
$X_{V_{k'}} = \overline{\mathbb{O}}_{\lambda^t}$.  For
type~$A$, $\lambda^t$ is the transpose partition, and
$\mathbb{O}_{\lambda^t} = d(\mathbb{O}_\lambda)$ by definition
of Barbasch--Vogan duality.
\end{proof}

\begin{example}[$\mathfrak{sl}_3$ orbit duality]
\label{ex:orbit-duality-sl3}
\index{orbit duality!sl3@$\mathfrak{sl}_3$}
For $\mathfrak{g} = \mathfrak{sl}_3$, there are three nilpotent
orbits, classified by partitions of~$3$:
\begin{center}
\renewcommand{\arraystretch}{1.3}
\begin{tabular}{lccc}
\toprule
\textbf{Partition} & \textbf{Orbit} & \textbf{Dimension}
& $d(\mathbb{O})$ \\
\midrule
$(1,1,1)$ & $\{0\}$ (trivial) & $0$ & $\mathbb{O}_{(3)}$ \\
$(2,1)$ & $\mathbb{O}_{\min}$ (minimal $=$ subregular)
& $4$ & $\mathbb{O}_{(2,1)}$ \\
$(3)$ & $\mathbb{O}_{\mathrm{prin}}$ (principal)
& $6$ & $\{0\}$ \\
\bottomrule
\end{tabular}
\end{center}
The minimal orbit is \emph{self-dual} under BV duality
(since $(2,1)^t = (2,1)$), while the trivial and principal
orbits are exchanged.

At non-degenerate admissible level $k = -3 + p/q$ ($p \geq 3$,
$\gcd(p,q) = 1$):
$V_k(\mathfrak{sl}_3)$ is rational with
$X_{V_k} = \{0\}$, while the Koszul dual at
$k' = -3 - p/q$ has
$X_{V_{k'}} = \overline{\mathbb{O}}_{\mathrm{prin}} = \mathcal{N}$.
This illustrates the general phenomenon that Koszul duality maps
rational vertex algebras to maximally non-rational ones.

At the ``boundary'' admissible level
$k = -3 + 3/2 = -3/2$, the associated variety is
$X_{V_{-3/2}} = \overline{\mathbb{O}}_{\min}$, and the Koszul
dual at $k' = -3 - 3/2 = -9/2$ also has
$X_{V_{-9/2}} = \overline{\mathbb{O}}_{\min}$: both lie on the
self-dual minimal orbit.
\end{example}

\subsection{Logarithmic modules and non-semisimple tensor categories}
\label{sec:logarithmic-modules}

When $V$ is not rational, the module category is
\emph{non-semisimple}: indecomposable modules need not be
simple, and $L_0$ may act non-semisimply (with Jordan blocks).
This occurs in two distinct situations:
\begin{enumerate}[label=\textup{(\alph*)}]
\item $X_V \neq \{0\}$ (not $C_2$-cofinite): infinitely many
  simples, and the module category has no finiteness control;
\item $X_V = \{0\}$ ($C_2$-cofinite) but not rational:
  finitely many simples, but the module category is a
  \emph{finite non-semisimple} tensor category---the
  logarithmic case.
\end{enumerate}
Such modules are called \emph{logarithmic} because the
correlation functions involve $\log(z - w)$ terms arising from
the Jordan block structure.

\begin{definition}[Logarithmic module]
\label{def:logarithmic-module}
\index{logarithmic module|textbf}
\index{Jordan block!L0@$L_0$}
A module $\mathcal{M}$ over a vertex algebra $V$ is
\emph{logarithmic} if the $L_0$-action on $\mathcal{M}$ is not
semisimple: there exist generalized eigenvectors
$v \in \mathcal{M}$ with $(L_0 - h)^k v = 0$ but
$(L_0 - h)^{k-1} v \neq 0$ for some $k \geq 2$.

The \emph{rank} of the logarithmic module (at conformal
weight~$h$) is the maximal Jordan block size of $L_0$ on
$\mathcal{M}_h$.  A rank-$1$ logarithmic module is an ordinary
(non-logarithmic) module.
\end{definition}

\begin{proposition}[Logarithmic modules and bar complex extensions;
\ClaimStatusProvedHere]
\label{prop:logarithmic-bar}
\index{logarithmic module!bar complex}
\index{bar complex!logarithmic modules}
Let $V$ be a vertex algebra that is not rational, and let
$\mathcal{L}_1, \mathcal{L}_2$ be simple $V$-modules.  A
non-trivial extension
\begin{equation}\label{eq:log-extension}
0 \to \mathcal{L}_1 \to \mathcal{E} \to \mathcal{L}_2 \to 0
\end{equation}
in the module category corresponds to a class in
$\mathrm{Ext}^1_V(\mathcal{L}_2, \mathcal{L}_1)$, which is
computed by the bar complex:
\begin{equation}\label{eq:ext-bar}
\mathrm{Ext}^n_V(\mathcal{L}_2, \mathcal{L}_1)
\;\cong\; H^n(\mathrm{Hom}_V(\bar{B}(\mathcal{L}_2),
\mathcal{L}_1)).
\end{equation}
\end{proposition}

\begin{proof}
The bar construction for modules provides a resolution
$\bar{B}(V, \mathcal{L}_2)$ of $\mathcal{L}_2$ by
$V$-modules built from the bar complex of~$V$; applying
$\mathrm{Hom}_V(-, \mathcal{L}_1)$ and taking cohomology
computes $\mathrm{Ext}^*_V(\mathcal{L}_2, \mathcal{L}_1)$.
\end{proof}

\begin{remark}[Jordan blocks and spectral sequence differentials]
\label{rem:jordan-bar-spectral}
\index{Jordan block!spectral sequence}
Heuristically, the Jordan block structure of logarithmic modules
is reflected in the bar spectral sequence as follows.  Filter an
indecomposable extension $\mathcal{E}$ by conformal weight; on
the $E_1$ page, $\mathrm{gr}\,\mathcal{E} \cong \mathcal{L}_1
\oplus \mathcal{L}_2$ as graded vector spaces.  The extension
class $[\mathcal{E}] \in \mathrm{Ext}^1$ is represented by a
bar $1$-cocycle appearing as a differential connecting the two
summands.  More generally, a length-$r$ chain of extensions
(iterated extension of $r$ simples) corresponds to non-trivial
differentials on the $E_r$ page, analogous to the relationship
between Massey products and higher spectral sequence differentials.
\end{remark}

\begin{example}[Triplet $\mathcal{W}(p)$ at $c = c_{2,1} = -2$]
\label{ex:triplet-logarithmic}
\index{triplet algebra|textbf}
\index{logarithmic module!triplet}
The triplet vertex algebra $\mathcal{W}(2)$ at $c = -2$
(the $(2,1)$~model) is $C_2$-cofinite but \emph{not} rational.
By Adamovi\'c--Milas~\cite{AdamovicMilas2008}, $\mathcal{W}(p)$ has exactly $2p$ simple
modules; for $p = 2$ these are:
\begin{itemize}
\item $\Lambda(1)$ ($h = 0$, vacuum) and
$\Lambda(2)$ ($h = -\tfrac{1}{8}$): non-projective simples;
\item $\Pi(1)$ ($h = -\tfrac{1}{8}$) and
$\Pi(2)$ ($h = \tfrac{3}{8}$): projective-simple (simple and
projective).
\end{itemize}
The module category is non-semisimple.  The indecomposable
projective covers of the non-projective simples are:
\begin{align}
P(1) &: \quad \Lambda(1) \hookrightarrow P(1) \twoheadrightarrow
\Lambda(1), \quad \text{with socle filtration }
\Lambda(1) \mid \Lambda(2) \oplus \Lambda(2) \mid
\Lambda(1), \\
P(2) &: \quad \Lambda(2) \hookrightarrow P(2) \twoheadrightarrow
\Lambda(2), \quad \text{with socle filtration }
\Lambda(2) \mid \Lambda(1) \mid \Lambda(2).
\end{align}
(The modules $\Pi(1)$ and $\Pi(2)$ are their own projective covers.)
The extension groups among non-projective simples are
$\dim \mathrm{Ext}^1(\Lambda(1), \Lambda(2)) = 2$
and $\dim \mathrm{Ext}^1(\Lambda(2), \Lambda(1)) = 1$.

Note that $\mathcal{W}(2)$ is $C_2$-cofinite, so
$X_{\mathcal{W}(2)} = \{0\}$; the non-semisimplicity arises
despite the trivial associated variety (case~(b) of the
classification above).

In the bar complex framework: $\bar{B}(\mathcal{W}(2))$ has
non-trivial $H^1$ detecting these extensions
(Proposition~\ref{prop:logarithmic-bar}).  As described in
Remark~\ref{rem:jordan-bar-spectral}, the rank-$2$ Jordan block
in $P(1)$ corresponds to a non-trivial $d_2$ differential
connecting the $\Lambda(1)$ and $\Lambda(2)$ sectors.

The fusion ring of $\mathcal{C}_{\mathcal{W}(2)}$ is
non-semisimple, with
$\Lambda(2) \boxtimes \Lambda(2) \cong P(1)$
(not $\Lambda(1)$: the fusion product is indecomposable).
Fusion with a projective-simple $\Pi(s)$ always gives a
projective module, reflecting the tensor ideal property of
projectives in the non-semisimple category.
Under Koszul duality, the non-semisimple fusion rules on the
$\mathcal{W}(2)$ side correspond to a non-trivial
$A_\infty$~structure on the Koszul dual, as predicted by the
curved $A_\infty$ framework of
\S\ref{sec:deviations-triviality}.
\end{example}

\begin{proposition}[$\mathrm{Ext}$ groups for $\mathcal{W}(2)$ via bar
resolution; \ClaimStatusProvedHere]
\label{prop:w2-ext-bar}
\index{triplet algebra!Ext computation}
\index{Ext groups!triplet}
The bar resolution of the non-projective simples $\Lambda(s)$
\textup{($s = 1,2$)} by $\mathcal{W}(2)$-modules yields
minimal projective resolutions:
\begin{align}
\label{eq:w2-resolution-1}
\cdots \to P(2)^{\oplus 2} &\to P(1) \to \Lambda(1) \to 0, \\
\label{eq:w2-resolution-2}
\cdots \to P(1) &\to P(2) \to \Lambda(2) \to 0.
\end{align}
The resulting $\mathrm{Ext}$ groups among non-projective simples are:
\begin{equation}\label{eq:w2-ext}
\dim \mathrm{Ext}^1_{\mathcal{W}(2)}(\Lambda(1), \Lambda(2)) = 2,
\qquad
\dim \mathrm{Ext}^1_{\mathcal{W}(2)}(\Lambda(2), \Lambda(1)) = 1.
\end{equation}
\end{proposition}

\begin{proof}
By Proposition~\ref{prop:logarithmic-bar}, the bar complex
$\bar{B}(\mathcal{W}(2), \Lambda(s))$ provides a resolution of
$\Lambda(s)$ by $\mathcal{W}(2)$-modules, and
equation~\eqref{eq:ext-bar} computes the Ext groups.  Since
$\mathcal{W}(2)$ is $C_2$-cofinite, the bar resolution can be
refined to a minimal projective resolution (the bar spectral
sequence degenerates at a finite page by
Proposition~\ref{prop:bar-admissible}(i)).

\emph{Step~1: Resolution of $\Lambda(1)$.}
The projective cover of $\Lambda(1)$ is $P(1)$, which maps
onto $\Lambda(1)$ with kernel
$\Omega^1(\Lambda(1)) = \mathrm{rad}(P(1))$.  By the socle
filtration of $P(1)$ stated above,
$\mathrm{rad}(P(1))/\mathrm{rad}^2(P(1)) \cong
\Lambda(2)^{\oplus 2}$: the radical quotient has $\Lambda(2)$
appearing with multiplicity~$2$.  The projective cover of
$\Omega^1(\Lambda(1))$ is therefore $P(2)^{\oplus 2}$, giving
the first two terms of~\eqref{eq:w2-resolution-1}.

\emph{Step~2: Resolution of $\Lambda(2)$.}
Similarly, $P(2) \twoheadrightarrow \Lambda(2)$ with kernel
$\Omega^1(\Lambda(2)) = \mathrm{rad}(P(2))$.  By the socle
filtration of $P(2)$,
$\mathrm{rad}(P(2))/\mathrm{rad}^2(P(2)) \cong \Lambda(1)$:
the radical quotient has $\Lambda(1)$ appearing once.  The
projective cover is $P(1)$, giving
\eqref{eq:w2-resolution-2}.

\emph{Step~3: Ext computation.}
Applying $\mathrm{Hom}_{\mathcal{W}(2)}(-, \Lambda(2))$ to the
resolution~\eqref{eq:w2-resolution-1}: since
$\mathrm{head}(P(s)) = \Lambda(s)$, we have
$\mathrm{Hom}(P(s), \Lambda(t)) = \delta_{st}\,\mathbb{C}$.
The complex becomes
\[
0 \to \underbrace{\mathrm{Hom}(P(1), \Lambda(2))}_{= 0}
\to \underbrace{\mathrm{Hom}(P(2)^{\oplus 2}, \Lambda(2))}_{= \mathbb{C}^2}
\to \cdots
\]
so $\mathrm{Ext}^1(\Lambda(1), \Lambda(2))
= \ker(d^1)/\mathrm{im}(d^0) = \mathbb{C}^2$.

Similarly, applying $\mathrm{Hom}(-, \Lambda(1))$ to
\eqref{eq:w2-resolution-2}: the complex becomes
\[
0 \to \underbrace{\mathrm{Hom}(P(2), \Lambda(1))}_{= 0}
\to \underbrace{\mathrm{Hom}(P(1), \Lambda(1))}_{= \mathbb{C}}
\to \cdots
\]
giving $\mathrm{Ext}^1(\Lambda(2), \Lambda(1)) = \mathbb{C}$.
The dimensions agree with the Adamovi\'c--Milas
classification~\cite{AdamovicMilas2008}.
\end{proof}

\begin{remark}[Ext quiver of $\mathcal{W}(2)$]
\label{rem:w2-ext-quiver}
\index{Ext quiver!triplet}
The Ext quiver of $\mathcal{W}(2)$ (restricted to non-projective
simples) has vertices $\{\Lambda(1), \Lambda(2)\}$ and arrows:
two arrows $\Lambda(1) \rightrightarrows \Lambda(2)$ and one
arrow $\Lambda(2) \to \Lambda(1)$.  The path algebra of this
quiver, modulo the relations imposed by the $A_\infty$ structure,
recovers the endomorphism algebra $\mathrm{End}(P(1) \oplus P(2))$
controlling the block.  This demonstrates concretely how the bar
complex (through its projective resolution) encodes the
representation-theoretic data of a non-semisimple module category.
\end{remark}

\begin{remark}[Non-semisimple tensor categories and bar-cobar duality]
\label{rem:non-semisimple-tensor}
\index{tensor category!non-semisimple}
\index{Verlinde formula!non-semisimple}
When $V$ is $C_2$-cofinite but not rational, the module category
$\mathcal{C}_V$ is a finite non-semisimple tensor category in
the sense of Etingof--Gelaki--Nikshych--Ostrik~\cite{EGNO}.
The fusion product $\boxtimes$ on $\mathcal{C}_V$ is exact in
each variable, and the Grothendieck ring $K_0(\mathcal{C}_V)$
is a $\mathbb{Z}_+$-ring.

The bar complex detects the non-semisimplicity as follows:
\begin{enumerate}[label=\textup{(\roman*)}]
\item \emph{Semisimple case} ($X_V = \{0\}$, rational):
  the bar spectral sequence degenerates at $E_2$, modules are
  projective, and the Verlinde formula computes
  $\dim \mathbb{V}$ exactly.
\item \emph{Non-semisimple case} ($X_V = \{0\}$,
  $C_2$-cofinite but not rational): the bar spectral sequence
  has non-trivial higher differentials encoding the extensions
  between simple modules.  The ``modified Verlinde formula'' of
  Creutzig--Ridout~\cite{CreutzigRidout2013} replaces the $S$-matrix
  by the \emph{pseudo-trace functions} (characters of projective
  modules), and the bar complex computes these via the full
  projective resolution.
\item \emph{Wild case} ($X_V \neq \{0\}$): the bar spectral
  sequence need not converge, and the module category may fail
  to be a finite tensor category.  The bar-cobar framework
  reduces to the curved $A_\infty$ regime
  (\S\ref{sec:deviations-triviality}).
\end{enumerate}
\end{remark}

\section{Resolutions and character formulas}\label{sec:complete-calculations}

\subsection{The fundamental principle of homological triviality}

For a bounded cochain complex $V^\bullet$ of vector spaces, the Euler
characteristic satisfies
$\chi = \sum_i (-1)^i \dim V^i = \sum_i (-1)^i \dim H^i$.
If $V^\bullet$ is a resolution of~$M$ (i.e., $H^0 = M$ and $H^i = 0$
for $i > 0$), this gives
$\dim M = \sum_{i \geq 0} (-1)^i \dim V^i$.
Passing to graded characters
$\operatorname{ch}(V) = \sum_n \dim V_n\, q^n$:
\[
\operatorname{ch}(M) = \sum_{i \geq 0} (-1)^i \operatorname{ch}(V^i).
\]
When the resolution terms $V^i$ have special structure (e.g., are free
or induced modules), this alternating sum collapses to a closed form.

\subsection{Chiral algebras}

In the chiral setting, vector spaces are replaced by
$\mathcal{D}_X$-modules, tensor products respect locality
(singularities only on diagonals), multiplication is encoded by OPEs,
and configuration spaces provide the geometric arena.

\section{Deriving the chiral module resolution}\label{sec:chiral-module-resolution}

\subsection{Free chiral modules}

\begin{lemma}[Structure of free chiral modules; \ClaimStatusProvedHere]
Let $\mathcal{A}$ be a chiral algebra on X and V a $\mathcal{D}_X$-module. The free chiral $\mathcal{A}$-module generated by V is:
\[
\text{Free}_{\mathcal{A}}(V) = \bigoplus_{n \geq 0} \Gamma(C_n(X), j_*j^*(\mathcal{A}^{\boxtimes n} \boxtimes V))
\]
\end{lemma}

\begin{proof}
The $n$-th summand accounts for $n$~copies of $\mathcal{A}$ acting
on~$V$: on the configuration space $C_n(X)$ of $n$~distinct points,
the $n$~copies of $\mathcal{A}$ are placed without coincidences, and
the extension $j_*j^*$ allows poles along diagonals to encode OPE
singularities.  Taking global sections produces the space of allowed
fields.  Universality follows from the factorization property: any
$\mathcal{A}$-module map $V \to \mathcal{M}$ extends uniquely by
iterating the chiral action on the $\mathcal{A}^{\boxtimes n}$
factors.
\end{proof}

\subsection{The bar resolution for chiral modules}

\begin{definition}[Bar complex for chiral modules]
\label{def:bar-complex-modules}
For a chiral algebra $\mathcal{A}$ with augmentation $\varepsilon: \mathcal{A} \to \omega_X$ and module $\mathcal{M}$, define:
\[
\overline{B}_n^{\text{ch}}(\mathcal{A}, \mathcal{M}) = \mathcal{A} \otimes \overline{\mathcal{A}}^{\otimes n} \otimes \mathcal{M}
\]
where $\overline{\mathcal{A}} = \ker(\varepsilon)$ and the differential is:
\begin{align}
d(a_0 \otimes [a_1|\cdots|a_n] \otimes m) &= \mu(a_0 \otimes a_1) \otimes [a_2|\cdots|a_n] \otimes m \\
&+ \sum_{i=1}^{n-1} (-1)^i a_0 \otimes [a_1|\cdots|a_i \cdot a_{i+1}|\cdots|a_n] \otimes m \\
&+ (-1)^n a_0 \otimes [a_1|\cdots|a_{n-1}] \otimes \mu_{\mathcal{M}}(a_n \otimes m)
\end{align}
\end{definition}

\begin{theorem}[Bar resolution is acyclic; \ClaimStatusProvedHere]
The bar complex is a resolution: $H^0(\overline{B}^{\text{ch}}) = \mathcal{M}$ and $H^i(\overline{B}^{\text{ch}}) = 0$ for $i > 0$.
\end{theorem}

\begin{proof}[First proof: direct]
Define a contracting homotopy $s: \overline{B}_n \to \overline{B}_{n+1}$ by:
\[
s(a_0 \otimes [a_1|\cdots|a_n] \otimes m) = 1 \otimes [a_0|a_1|\cdots|a_n] \otimes m
\]
where we use $a_0 = \varepsilon(a_0) \cdot 1 + \overline{a_0}$ with $\overline{a_0} \in \overline{\mathcal{A}}$.

Computing:
\begin{align}
(ds + sd)(a_0 \otimes [a_1|\cdots|a_n] \otimes m) &= \varepsilon(a_0) \cdot 1 \otimes [a_1|\cdots|a_n] \otimes m \\
&+ \text{terms with } \overline{a_0}
\end{align}

For normalized chains (where $a_i \in \overline{\mathcal{A}}$), we get $ds + sd = \text{id}$, proving acyclicity.
\end{proof}

\begin{proof}[Second proof: spectral sequence]
Filter the bar complex by the number of bars:
\[
F_p = \bigoplus_{n \leq p} \overline{B}_n
\]

The associated graded is:
\[
\text{gr}_p = \mathcal{A} \otimes (s\overline{\mathcal{A}})^{\otimes p} \otimes \mathcal{M}
\]

The $E_1$ page computes cohomology of the associated graded.  Since $\mathcal{A}$ is augmented and \emph{connected} (meaning $\overline{\mathcal{A}} = \ker(\varepsilon)$ is concentrated in positive conformal weights), the complex $T(s\overline{\mathcal{A}})$ (the tensor coalgebra on the suspension of $\overline{\mathcal{A}}$) is acyclic by the standard augmentation filtration argument: the filtration by conformal weight is bounded below and exhaustive, so the spectral sequence converges. Therefore $E_2^{p,q} = 0$ for $p > 0$, and the spectral sequence degenerates, proving acyclicity.

(The connectedness hypothesis is essential: for non-connected algebras such as $\mathcal{A} = \mathbb{C}[x]$ with $|x| = 0$, the bar complex need not be acyclic.)
\end{proof}

\subsection{Geometric realization on configuration spaces}

\begin{theorem}[Geometric bar complex; \ClaimStatusProvedHere]
\label{thm:geometric-bar-module}
The bar complex has a geometric realization:
\[
\overline{B}_n^{\text{geom}}(\mathcal{A}, \mathcal{M}) = \Gamma(\overline{C}_{n+2}(X), j_*j^*(\mathcal{A} \boxtimes \overline{\mathcal{A}}^{\boxtimes n} \boxtimes \mathcal{M}) \otimes \Omega^n_{\log})
\]
\end{theorem}

\begin{proof}
The construction parallels the algebra bar complex (Definition~\ref{def:geometric-bar}).  An element $a_0 \otimes [a_1|\cdots|a_n] \otimes m$ corresponds to a section of $\mathcal{A} \boxtimes \overline{\mathcal{A}}^{\boxtimes n} \boxtimes \mathcal{M}$ on the open configuration space $C_{n+2}(X)$, where the $n+2$ points are the output at $z_0$, intermediate insertions at $z_1, \ldots, z_n$, and the module input at $z_{n+1}$.  The $j_*$-extension to $\overline{C}_{n+2}(X)$ provides logarithmic poles along the collision divisors.

The differential $d = d_{\mathrm{int}} + d_{\mathrm{res}} + d_{\mathrm{dR}}$ is defined as in the algebra case, with the modification that the module point $z_{n+1}$ can collide with algebra points but uses the module action $\rho\colon \mathcal{A} \boxtimes \mathcal{M} \to \Delta_*\mathcal{M}$ rather than the algebra product.  The verification that $d^2 = 0$ follows the nine-term argument of Theorem~\ref{thm:bar-nilpotency-complete}, with the $\mathcal{A}$-module axiom (associativity of the action $\rho$) replacing the algebra associativity in Case~(ii) when one of the colliding points is the module point.
\end{proof}

\section{Computing characters via resolutions}\label{sec:computing-characters}

\subsection{The fundamental character formula}

\begin{theorem}[Character via acyclic resolution; \ClaimStatusProvedHere]
If $\mathcal{P}_\bullet \to \mathcal{M}$ is an acyclic resolution, then:
\[
\text{ch}(\mathcal{M}) = \sum_{n=0}^\infty (-1)^n \text{ch}(\mathcal{P}_n)
\]
\end{theorem}

\begin{proof}[First proof: Euler characteristic]
For each weight space, the complex $\mathcal{P}_\bullet^{(\lambda)}$ of weight $\lambda$ components has Euler characteristic:
\[
\chi(\mathcal{P}_\bullet^{(\lambda)}) = \sum_n (-1)^n \dim \mathcal{P}_n^{(\lambda)} = \dim \mathcal{M}^{(\lambda)}
\]
since the complex is acyclic. Summing over weights with $q^{\lambda}$ gives the character formula.
\end{proof}

\begin{proof}[Second proof: long exact sequences]
Write $Z_n = \ker(d_n)$, $B_n = \text{im}(d_{n+1})$. The short exact sequences:
\[
0 \to Z_n \to \mathcal{P}_n \to B_{n-1} \to 0
\]
give $\text{ch}(\mathcal{P}_n) = \text{ch}(Z_n) + \text{ch}(B_{n-1})$.

Since $H^n = Z_n/B_n = 0$ for $n > 0$, we have $Z_n = B_n$. Telescoping:
\[
\sum_{n=0}^N (-1)^n \text{ch}(\mathcal{P}_n) = \text{ch}(Z_0) - (-1)^N \text{ch}(B_N)
\]
As $N \to \infty$, $B_N \to 0$ (assuming appropriate convergence), giving $\text{ch}(\mathcal{M}) = \text{ch}(Z_0)$.
\end{proof}

\begin{proof}[Third proof: Hodge theory]
Equip $\mathcal{P}_\bullet$ with an inner product. The Hodge Laplacian $\Delta = dd^* + d^*d$ has:
\[
\ker \Delta \cong H^*(\mathcal{P}_\bullet)
\]
By the McKean--Singer theorem, the supertrace $\mathrm{STr}(e^{-t\Delta}) = \sum_n (-1)^n \mathrm{Tr}(e^{-t\Delta}|_{\mathcal{P}_n})$ is \emph{independent of $t$}.  Evaluating at $t \to 0^+$ gives $\sum_n (-1)^n \mathrm{ch}(\mathcal{P}_n)$; evaluating at $t \to \infty$ projects onto $\ker \Delta \cong H^*(\mathcal{P}_\bullet)$, giving $\mathrm{ch}(\mathcal{M})$ (since $H^n = 0$ for $n > 0$ by acyclicity, only $H^0 = \mathcal{M}$ contributes).
\end{proof}

\subsection{Koszul duality for modules}

\begin{theorem}[Koszul pairs simplify resolutions; \ClaimStatusProvedHere]
\label{thm:koszul-resolution-module}
If $(\mathcal{A}, \mathcal{A}^!)$ are Koszul dual chiral algebras, then for any $\mathcal{A}$-module $\mathcal{M}$:
\[
\mathcal{P}_n(\mathcal{M}) = \mathcal{A} \otimes (\mathcal{A}^!)_n \otimes \mathcal{M}
\]
provides a minimal resolution.
\end{theorem}

\begin{proof}
Koszul duality means $\text{Ext}^i_{\mathcal{A}}(\omega_X, \omega_X) = (\mathcal{A}^!)_i$. The bar resolution of $\omega_X$ is:
\[
\cdots \to \mathcal{A} \otimes \overline{\mathcal{A}}^{\otimes n} \to \cdots \to \mathcal{A} \to \omega_X
\]

Taking homology and using Koszul duality:
\[
H^n = \begin{cases} \omega_X & n = 0 \\ 0 & n > 0 \end{cases}
\]

The complex $\mathcal{A} \otimes (\mathcal{A}^!)_* $ is the minimal model, having no excess terms. Tensoring with $\mathcal{M}$ preserves this minimality.
\end{proof}

\begin{corollary}[Character formula for Koszul case; \ClaimStatusProvedHere]
For Koszul dual pair $(\mathcal{A}, \mathcal{A}^!)$, the character of a module $\mathcal{M}$ computed from the Koszul resolution equals the naive character:
\[
\mathrm{ch}(\mathcal{M}) = \mathrm{ch}_{\mathrm{naive}}(\mathcal{M}).
\]
\end{corollary}

\begin{proof}
The Koszul resolution gives:
\begin{align}
\mathrm{ch}(\mathcal{M}) &= \sum_n (-1)^n \mathrm{ch}(\mathcal{A} \otimes (\mathcal{A}^!)_n \otimes \mathcal{M}) \\
&= \mathrm{ch}(\mathcal{A}) \cdot \mathrm{ch}_{\mathrm{naive}}(\mathcal{M}) \cdot \sum_n (-1)^n \mathrm{ch}((\mathcal{A}^!)_n) \\
&= \mathrm{ch}(\mathcal{A}) \cdot \mathrm{ch}_{\mathrm{naive}}(\mathcal{M}) \cdot h_{\mathcal{A}^!}(-t).
\end{align}
By the Koszul duality relation $h_{\mathcal{A}}(t) \cdot h_{\mathcal{A}^!}(-t) = 1$ (equivalent to Koszulity of~$\mathcal{A}$), we have $h_{\mathcal{A}^!}(-t) = 1/h_{\mathcal{A}}(t) = 1/\mathrm{ch}(\mathcal{A})$. Hence $\mathrm{ch}(\mathcal{M}) = \mathrm{ch}_{\mathrm{naive}}(\mathcal{M})$.

The non-trivial content of this identity is that the resolution is exact: exactness of the Koszul resolution is \emph{equivalent} to the Hilbert series relation $h_{\mathcal{A}} \cdot h_{\mathcal{A}^!}(-t) = 1$.
\end{proof}

\section{\texorpdfstring{The structure theory: $A_\infty$, $L_\infty$, and Gerstenhaber}{The structure theory: A-infinity, L-infinity, and Gerstenhaber}}

\subsection{\texorpdfstring{$A_\infty$ structure on resolutions}{A-infinity structure on resolutions}}

\begin{theorem}[$A_\infty$ module structure; \ClaimStatusProvedElsewhere]\label{thm:ainfty-module}
The geometric resolution $\mathcal{P}_\bullet(\mathcal{M})$ carries a
natural $A_\infty$-module structure over $\mathcal{A}$ with operations
\[
m_n: \mathcal{A}^{\otimes n-1} \otimes \mathcal{P}_\bullet
\to \mathcal{P}_\bullet[2-n]
\]
satisfying the $A_\infty$-module relations
\[
\sum_{\substack{r+s+t=n \\ s \geq 1}} (-1)^{rs+t}
m_{r+1+t}(\mathrm{id}^{\otimes r} \otimes m_s \otimes \mathrm{id}^{\otimes t}) = 0.
\]
The operations are given by configuration space integrals:
$m_1$ is the bar differential,
\[
m_2(a \otimes p) =
\mathrm{Res}_{z_a \to z_p}\, Y(a, z_a{-}z_p) \cdot p
\]
is the residue action, and higher $m_n$ are given by
iterated residues on $\overline{C}_n(X)$ weighted by the
propagator forms.  The $A_\infty$ relations follow from
Stokes' theorem on the Fulton--MacPherson compactification
together with the Arnold--Orlik--Solomon relations
(Theorem~\ref{thm:arnold-relations}), exactly as in the
bar nilpotency proof (Theorem~\ref{thm:bar-nilpotency-complete}).
\end{theorem}

\begin{remark}
The $A_\infty$-module structure on the bar resolution is a special case
of the general homotopy transfer theorem
(Appendix~\ref{app:homotopy-transfer}; see
Keller~\cite{Keller01} and Merkulov~\cite{Merkulov99} for the
original constructions).  In the chiral setting, the transfer
data come from the operad structure on configuration spaces
(Definition~\ref{def:chiral-operad}).
\end{remark}

\subsection{\texorpdfstring{$L_\infty$ structure}{L-infinity structure}}

\begin{theorem}[$L_\infty$ structure on cochains; \ClaimStatusProvedElsewhere]\label{thm:linfty-cochains}
The cochain complex $\mathrm{RHom}_{\mathcal{A}}(\mathcal{P}_\bullet,
\mathcal{P}_\bullet)$ carries an $L_\infty$-algebra structure with brackets
\[
\ell_n: \textstyle\bigwedge^n \mathrm{RHom} \to \mathrm{RHom}[2-n].
\]
The first two brackets are $\ell_1(f) = [d, f]$ and
$\ell_2(f, g) = (-1)^{|f|}[f, g]$ (the graded commutator);
the higher brackets $\ell_n$ for $n \geq 3$ are
the Massey-type operations arising from the homotopy
transfer of the dg Lie structure on endomorphisms.
This is a consequence of the Kontsevich--Soibelman formality
framework~\cite{KontsevichSoibelman}.
\end{theorem}

\subsection{Chiral Gerstenhaber structure}

\begin{theorem}[Chiral Gerstenhaber algebra; \ClaimStatusProvedElsewhere]\label{thm:chiral-gerstenhaber}
The chiral Hochschild cohomology
$HH^*_{\mathrm{ch}}(\mathcal{A}, \mathcal{A})$
carries a Gerstenhaber algebra structure: a graded-commutative cup product
$\cup: HH^p \otimes HH^q \to HH^{p+q}$ and a degree-$(-1)$ Lie bracket
$\{-,-\}: HH^p \otimes HH^q \to HH^{p+q-1}$ satisfying the Leibniz
rule
\[
\{f, g \cup h\} = \{f, g\} \cup h + (-1)^{(|f|-1)|g|}\, g \cup \{f, h\}.
\]
The cup product arises from composition of multilinear chiral operations;
the bracket from the commutator of vertex operators, realized via
contraction of vector fields against forms on configuration spaces.
The Gerstenhaber structure on Hochschild cochains of associative algebras
is due to Gerstenhaber~\cite{Ger63}; the chiral generalization follows
from the action of the little discs operad on configuration spaces
(Tamarkin~\cite{Tamarkin00}, Francis--Gaitsgory).
\end{theorem}

\section{Denominator formulas and characters}\label{sec:denominator-formulas}

\subsection{The trivial module}

\begin{theorem}[Denominator identity for trivial module; \ClaimStatusProvedElsewhere]
For a chiral algebra $\mathcal{A}$ arising from an affine Lie algebra $\hat{\mathfrak{g}}_k$, the trivial module $\omega_X$ has character governed by the Weyl--Kac denominator formula:
\[
1 = \frac{\sum_{w \in \mathcal{W}} \varepsilon(w)\, e^{w(\rho)}}{\prod_{n > 0} \prod_{\alpha \in \Delta^+} (1 - q^n e^{-\alpha})^{\mathrm{mult}(\alpha)}}
\]
where $\mathcal{W}$ is the affine Weyl group, $\varepsilon(w) = (-1)^{\ell(w)}$ is the sign character, $\rho$ is the affine Weyl vector, $\Delta^+$ are the positive roots of $\hat{\mathfrak{g}}$, and $\mathrm{mult}(\alpha)$ is the root multiplicity.  In the geometric picture, these multiplicities are computed as:
\[
\mathrm{mult}_n(\alpha) = \dim H^0(\overline{C}_n(X), \mathcal{L}_\alpha \otimes \Omega^n_{\log})
\]
where $\mathcal{L}_\alpha$ denotes the line bundle on $\overline{C}_n(X)$ determined by the weight $\alpha$ via the Cartan subalgebra action on $\mathcal{A}^{\boxtimes n}$.
\end{theorem}

\begin{proof}[Geometric interpretation]
\emph{Step~1: Resolution.}
Construct the bar resolution of the vacuum:
\[
\cdots \to \mathcal{P}_2 \to \mathcal{P}_1 \to \mathcal{P}_0
\to \omega_X \to 0
\]
where $\mathcal{P}_n = \Gamma(\overline{C}_{n+1}(X),
j_*j^*\mathcal{A}^{\boxtimes n} \otimes \Omega^n_{\log})$.

\emph{Step~2: Characters of resolution terms.}
By Riemann--Roch on $\overline{C}_{n+1}(X)$:
\[
\operatorname{ch}(\mathcal{P}_n)
= \sum_{\alpha} q^{n\alpha}\, \operatorname{mult}_n(\alpha), \qquad
\operatorname{mult}_n(\alpha)
= \chi(\overline{C}_{n+1},\, \mathcal{O}(\alpha)
\otimes \Omega^n_{\log}).
\]

\emph{Step~3: Alternating sum.}
Acyclicity of the bar resolution gives:
\[
1 = \sum_{n \geq 0} (-1)^n \operatorname{ch}(\mathcal{P}_n)
= \frac{\sum_{w \in \mathcal{W}} \varepsilon(w)\, e^{w(\rho)}}
  {\prod_{n,\alpha} (1 - q^n e^{-\alpha})^{\operatorname{mult}_n(\alpha)}}
\]
where the denominator arises from the universal bar complex
contribution, and the numerator from the Weyl group action on
highest weights via the factorization structure.
\end{proof}

\subsection{BGG resolution from the bar complex}

\begin{theorem}[BGG resolution from bar complex; \ClaimStatusProvedHere]
\label{thm:bgg-from-bar}
\index{BGG resolution!from bar complex}
\index{bar complex!BGG resolution}
Let $\mathcal{A} = \widehat{\mathfrak{g}}_k$ at generic level~$k$
(Proposition~\ref{prop:generic-irreducibility}), and let
$\mathcal{L}(\lambda)$ be a simple highest weight module with
dominant integral weight~$\lambda$.  The bar spectral sequence for
$\mathcal{L}(\lambda)$ produces a resolution by Verma modules:
\begin{equation}\label{eq:bgg-from-bar}
\cdots \to \bigoplus_{\ell(w)=2} \mathcal{M}(w \cdot \lambda)
\to \bigoplus_{\ell(w)=1} \mathcal{M}(w \cdot \lambda)
\to \mathcal{M}(\lambda) \to \mathcal{L}(\lambda) \to 0
\end{equation}
where $w$ ranges over the affine Weyl
group~$\mathcal{W}$, $w \cdot \lambda = w(\lambda + \rho) - \rho$
is the dot action, and the differentials are screening operator
integrals on~$\overline{C}_k(X)$.  This recovers the classical
Bernstein--Gel'fand--Gel'fand resolution
\textup{(\cite{Kac}, \S10.7; \cite{BGG76})} from the bar-cobar
machinery.
\end{theorem}

\begin{proof}
The argument proceeds in four steps.

\emph{Step~1: $E_1$ page.}
The bar complex of the $\mathcal{A}$-module
$\mathcal{L}(\lambda)$ carries a filtration by tensor degree
in the augmentation ideal
$\bar{\mathcal{A}} = \ker(\mathcal{A} \to \omega_X)$.  The
associated spectral sequence has $E_1$ page given by the weight
decomposition of the bar construction under the Cartan
subalgebra~$\mathfrak{h}$:
\[
E_1^{n,*} = \bigoplus_{\substack{w \in \mathcal{W} \\
\ell(w) = n}} \mathcal{M}(w \cdot \lambda).
\]
The identification of each summand with a Verma module uses the
PBW decomposition
$U(\widehat{\mathfrak{n}}^-) \cong
\mathrm{Sym}(\widehat{\mathfrak{n}}^-)$:
the $n$-fold bar construction at weight
$w \cdot \lambda - \lambda = w(\lambda + \rho) - (\lambda + \rho)$
is generated by the product of negative root vectors corresponding
to a reduced decomposition of~$w$, and the induced module is
$\mathcal{M}(w \cdot \lambda)$.  Geometrically, each summand
corresponds to sections on $\overline{C}_{n+1}(X)$ with logarithmic
singularities along boundary divisors
(Theorem~\ref{thm:geometric-bar-definition}).

\emph{Step~2: Vanishing of $d_1$ at generic level.}
The $d_1$ differential maps
$E_1^{n,*} \to E_1^{n+1,*}$, i.e.,
$\mathcal{M}(w \cdot \lambda) \to \mathcal{M}(w' \cdot \lambda)$
for $\ell(w') = \ell(w) + 1$.  Any non-zero map between Verma
modules factors through a singular vector embedding.  At generic
level, all Verma modules are irreducible
(Proposition~\ref{prop:generic-irreducibility}), so
$\mathrm{Hom}_{\mathcal{O}}(\mathcal{M}(\mu),
\mathcal{M}(\nu)) = \delta_{\mu\nu} \cdot \mathbb{C}$.
Therefore $d_1 = 0$ and $E_2 = E_1$.

\emph{Step~3: Degeneration.}
Since $E_2 = E_1$ consists of Verma modules (which are
bounded below in conformal weight), the spectral sequence is
bounded below and converges: $E_2 = E_\infty$.
The acyclicity of the total bar complex
(Theorem~\ref{thm:bar-cobar-isomorphism-main})
then yields an exact resolution of $\mathcal{L}(\lambda)$ by
Verma modules, indexed by the affine Weyl group and graded
by the length function.

\emph{Step~4: Screening operator differentials.}
The differentials
$d\colon \mathcal{M}(w' \cdot \lambda) \to
\mathcal{M}(w \cdot \lambda)$
(for $\ell(w') = \ell(w) + 1$, $w \leq w'$ in Bruhat order)
are realized by screening operator integrals
(Feigin--Frenkel~\cite{FF}):
\[
d_{w' \to w}(v) = \oint_{\gamma_{w,w'}}
S_{\alpha_{i_1}}(z_1) \cdots S_{\alpha_{i_m}}(z_m)\,
v\, \omega_{\log}
\]
where $S_{\alpha_i}$ are screening currents
(Theorem~\ref{thm:screening-bar}) and
$\omega_{\log}$ is the logarithmic form on $\overline{C}_m(X)$.
The composition $d^2 = 0$ follows from the Arnold relations
(Theorem~\ref{thm:arnold-relations}).
\end{proof}

\begin{remark}[BGG at special levels]\label{rem:bgg-special}
At non-generic levels, $d_1 \neq 0$ and the $E_2$ page is a
proper quotient of the $E_1$ page.  For admissible levels
$k = -h^\vee + p/q$, the surviving elements on $E_2$ form
the truncated Weyl group $\mathcal{W}_{\mathrm{adm}}$,
yielding the Kac--Wakimoto character formula
(Theorem~\ref{thm:kw-bar-spectral}).  For integrable levels,
$d_1 = 0$ and the full $E_1$ page survives, but the
resolution terminates at length
$|\mathcal{W}_{\mathrm{aff}}/\mathcal{W}_{\mathrm{fin}}|$
due to the vanishing of the Verma modules beyond the
affine Weyl chamber.
\end{remark}

\begin{computation}[BGG pipeline for $\widehat{\mathfrak{sl}}_2$
at generic level; \ClaimStatusProvedHere]
\label{comp:bgg-sl2-pipeline}
\index{BGG resolution!explicit pipeline}
\index{bar complex!BGG pipeline}
We trace the bar complex of $\widehat{\mathfrak{sl}}_{2,k}$ through
the BGG spectral sequence (Theorem~\ref{thm:bgg-from-bar}) step by
step, connecting the explicit bar differential matrix
(Computation~\ref{comp:sl2-bar-matrix}) to the classical BGG
resolution, and verifying compatibility with Koszul duality
(Theorem~\ref{thm:sl2-koszul-dual}) and screening operators
(Theorem~\ref{thm:screening-bar}).  This instantiates the abstract
machinery for a concrete algebra.

\emph{Step~1: Weight decomposition of the bar complex.}
The bar differential $d\colon \bar{B}^2 \to \bar{B}^1$
(eq.~\eqref{eq:sl2-bar-matrix}) decomposes under the adjoint
$\mathfrak{h}$-action into weight sectors:
\begin{center}
\renewcommand{\arraystretch}{1.2}
\begin{tabular}{cccc}
\textbf{Weight} & $\dim \bar{B}^2$ & $\dim \bar{B}^1$
  & $\operatorname{rank} d$ \\
\hline
$+4$ & $1$ & $0$ & $0$ \\
$+2$ & $2$ & $1$ & $1$ \\
$\phantom{+}0$ & $3$ & $1$ & $1$ \\
$-2$ & $2$ & $1$ & $1$ \\
$-4$ & $1$ & $0$ & $0$
\end{tabular}
\end{center}
The total rank is~$3$, confirming
$\dim \ker(d) = 9 - 3 = 6$
(Computation~\ref{comp:sl2-bar-matrix}).

\emph{Step~2: Identification with Weyl group elements.}
The finite Weyl group $W = \{e, s_\alpha\}$ acts on weight~$0$ by
the dot action $w \cdot 0 = w(\rho) - \rho$ with $\rho = 1$:
\begin{itemize}
\item $e \cdot 0 = \rho - \rho = 0$ (identity: the vacuum);
\item $s_\alpha \cdot 0 = -\rho - \rho = -2$
  (simple reflection).
\end{itemize}
In $\bar{B}^1$, the weight~$-2$ sector is one-dimensional,
spanned by $s^{-1}f$.  This is the bar-complex incarnation of the
Verma module $\mathcal{M}(s_\alpha \cdot 0) = \mathcal{M}(-2)$
at length~$1$ in the BGG resolution.  The weight~$0$ sector
(spanned by $s^{-1}h$) corresponds to the Cartan direction and
is absorbed into the $E_0$ page, while the weight~$+2$ sector
(spanned by $s^{-1}e$) lies outside the Weyl group orbit
$W \cdot 0 = \{0, -2\}$ and contributes to the
non-BGG part of the bar complex
(the bar complex has length $\dim \mathfrak{g} = 3$; the BGG
resolution has length $|W| - 1 = 1$).

\emph{Step~3: Screening operator as bar differential.}
The bar differential restricted to the weight-$(-2)$ sector
of~$\bar{B}^2 \to \bar{B}^1$ is:
\[
d([f|h]) = 2f, \qquad d([h|f]) = -2f.
\]
The image is $\operatorname{span}\{f\} = \bar{B}^1_{-2}$,
and the kernel $\{[f|h] + [h|f]\}$ encodes the antisymmetry
$[f,h] + [h,f] = 0$.  This image direction is the
screening operator $S_\alpha = \oint f(z)\,dz$
(Theorem~\ref{thm:screening-bar}).  At generic level, the
Verma module $\mathcal{M}(0)$ is irreducible
(Proposition~\ref{prop:generic-irreducibility}), so
$d_1 = 0$ in the spectral sequence and $E_\infty = E_1$
(Theorem~\ref{thm:bgg-from-bar}, Step~2).

\emph{Step~4: The BGG resolution.}
The degenerate spectral sequence yields the BGG resolution
of $\mathcal{L}(\Lambda_0)$ at generic level:
\begin{equation}\label{eq:bgg-sl2-generic}
\cdots \to \bigoplus_{\ell(w)=2} \mathcal{M}(w \cdot \Lambda_0)
\to \mathcal{M}(s_1 \cdot \Lambda_0)
   \oplus \mathcal{M}(s_0 \cdot \Lambda_0)
\to \mathcal{M}(\Lambda_0) \to \mathcal{L}(\Lambda_0) \to 0
\end{equation}
with differentials given by screening operator integrals.
For the affine Weyl group, the translation elements
$t_{n\alpha^\vee}$ contribute Verma modules at conformal
weight~$n^2$ (eq.~\eqref{eq:dot-action-translations}),
reproducing the Jacobi theta function $\vartheta_3(q)$ in the
character formula
(Proposition~\ref{prop:weyl-kac-sl2-bar}).

\emph{Step~5: Koszul duality of the resolution.}
Under $k \mapsto -k - 2h^\vee = -k-4$
(Theorem~\ref{thm:sl2-koszul-dual}), the module Koszul duality
functor (Proposition~\ref{prop:vacuum-verma-koszul}) sends each
term of~\eqref{eq:bgg-sl2-generic} to the corresponding term at
the dual level:
\[
\Phi\bigl(\mathcal{M}(w \cdot \Lambda_0)_k\bigr)
\simeq \mathcal{M}(w \cdot \Lambda_0)_{-k-4},
\qquad w \in \mathcal{W}.
\]
The BGG resolution at level~$k$ maps to the BGG resolution at
level~$-k-4$ with identical Weyl group combinatorics.
The Shapovalov determinant transforms as
$\det S_1(-k-4) = -\det S_1(k)$
(Proposition~\ref{prop:shapovalov-koszul}): singular vectors
at grade~$1$ appear simultaneously at both levels,
with the critical level $k = -h^\vee = -2$ being the unique
self-dual point ($-(-2)-4 = -2$).

\emph{Summary.}
This computation chains six proved results---Computation~\ref{comp:sl2-bar-matrix}
(explicit bar matrix), Theorem~\ref{thm:bgg-from-bar} (BGG from bar),
Theorem~\ref{thm:screening-bar} (screening operators),
Theorem~\ref{thm:sl2-koszul-dual} (Koszul duality of
$\widehat{\mathfrak{sl}}_2$),
Proposition~\ref{prop:vacuum-verma-koszul} (vacuum module duality),
and Proposition~\ref{prop:shapovalov-koszul} (Shapovalov
transform)---into a single pipeline that tracks the vacuum module
through the full bar-cobar-module duality.
All intermediate quantities (weight decomposition, kernel dimensions,
Shapovalov determinant ratios) have been verified computationally.
See \S\ref{sec:bgg-sl2-pipeline} for the detailed matrix-level
verification, including the explicit singular-vector conditions
extracted from bar differential entries and the formal BGG
involution under Koszul duality
(Corollary~\ref{cor:bgg-koszul-involution}).
\end{computation}

\subsection{General modules}

\begin{theorem}[Weyl--Kac character formula; \ClaimStatusProvedHere]
\label{thm:weyl-kac-geometric}
\index{Weyl--Kac character formula!geometric derivation}
Let $\mathcal{A} = \widehat{\mathfrak{g}}_k$ at generic level~$k$,
and let $\mathcal{L}(\lambda)$ be a simple highest weight module
with dominant integral weight~$\lambda$.  Then:
\[
\operatorname{ch}(\mathcal{L}(\lambda))
= \frac{\sum_{w \in \mathcal{W}} \varepsilon(w)\,
   e^{w(\lambda + \rho) - \rho}}
  {\prod_{n > 0} \prod_{\alpha \in \Delta^+}
   (1 - q^n e^{-\alpha})^{\operatorname{mult}(\alpha)}}
\]
where $\mathcal{W}$ is the affine Weyl group,
$\varepsilon(w) = (-1)^{\ell(w)}$, $\rho$ is the Weyl vector,
and $\Delta^+$ are the positive roots.  This is the Weyl--Kac
character formula \textup{(\cite{Kac}, Theorem~10.4)}, here
derived from the geometric bar complex.
\end{theorem}

\begin{proof}
By Theorem~\ref{thm:bgg-from-bar}, the bar complex provides
an exact resolution of $\mathcal{L}(\lambda)$ by Verma modules
$\mathcal{M}(w \cdot \lambda)$, $w \in \mathcal{W}$.
Taking the Euler characteristic:
\begin{align}
\operatorname{ch}(\mathcal{L}(\lambda))
&= \sum_{w \in \mathcal{W}} \varepsilon(w)\,
   \operatorname{ch}(\mathcal{M}(w \cdot \lambda)).
\end{align}
Each Verma module has character
$\operatorname{ch}(\mathcal{M}(\mu))
= e^{\mu} / \prod_{n > 0} \prod_{\alpha \in \Delta^+}
(1 - q^n e^{-\alpha})^{\operatorname{mult}(\alpha)}$
(the highest weight factor divided by the universal
denominator).  The denominator depends only on the algebra,
not on the highest weight, so it factors out, yielding
the stated formula.  The denominator arises from the bar complex
of $\mathcal{A}$ on configuration spaces
(trivial module computation above), while the numerator
records the Weyl group orbit via the bar filtration.
\end{proof}

\subsection{Explicit computation: \texorpdfstring{$\widehat{\mathfrak{sl}}_2$}{sl(2)-hat} at level~$1$}
\label{sec:sl2-level1-computation}

We illustrate the bar-cobar character formula
(Theorem~\ref{thm:weyl-kac-geometric}) by computing the Weyl--Kac
character of the vacuum module for $\widehat{\mathfrak{sl}}_2$ at
level~$1$.

\begin{example}[$\widehat{\mathfrak{sl}}_2$ at level~$1$: integrable modules]
\label{ex:sl2-level1}
\index{sl2-hat@$\widehat{\mathfrak{sl}}_2$!level 1}
At level $k=1$, there are exactly two integrable highest weight modules:
\begin{itemize}
\item $L(\Lambda_0)$: the \emph{vacuum module} (highest weight $\Lambda_0$,
  conformal dimension $h=0$);
\item $L(\Lambda_1)$: the \emph{fundamental module} (highest weight $\Lambda_1$,
  conformal dimension $h = 1/4$).
\end{itemize}
The affine Weyl group is $\mathcal{W} = \mathbb{Z} \rtimes \mathbb{Z}/2$,
generated by the reflections $s_0$ (affine) and $s_1$ (finite).
The Weyl vector is $\rho = \Lambda_0 + \Lambda_1$.
\end{example}

\begin{proposition}[BGG resolution of $L(\Lambda_0)$; \ClaimStatusProvedElsewhere]
\label{prop:bgg-sl2-level1}
The BGG resolution of the vacuum module $L(\Lambda_0)$ at level~$1$ is:
\begin{equation}\label{eq:bgg-sl2-level1}
\cdots \to \bigoplus_{\ell(w)=2} \mathcal{M}(w \cdot \Lambda_0)
\to \mathcal{M}(s_1 \cdot \Lambda_0) \oplus \mathcal{M}(s_0 \cdot \Lambda_0)
\to \mathcal{M}(\Lambda_0) \to L(\Lambda_0) \to 0.
\end{equation}
The dot action $w \cdot \Lambda_0 = w(\Lambda_0 + \rho) - \rho$ gives,
for the translation elements $t_{n\alpha^\vee} \in \mathcal{W}$:
\begin{equation}\label{eq:dot-action-translations}
t_{n\alpha^\vee} \cdot \Lambda_0 = \Lambda_0 - n^2\,\delta,
\end{equation}
where $\delta$ is the null root.
The Verma module $\mathcal{M}(\Lambda_0 - n^2\delta)$ has lowest
conformal weight $n^2$, so:
\begin{equation}\label{eq:verma-character-sl2}
\operatorname{ch}(\mathcal{M}(\Lambda_0 - n^2\delta))
= \frac{q^{n^2}}{\prod_{m \geq 1}(1-q^m)^3}.
\end{equation}
(The denominator $\prod(1-q^m)^3$ accounts for the three
generators $J^3_{-m}, e_{-m}, f_{-m}$ of
$\widehat{\mathfrak{sl}}_2$.)
\end{proposition}

\begin{proposition}[Character from bar resolution; \ClaimStatusProvedHere]
\label{prop:weyl-kac-sl2-bar}
The alternating character sum from the BGG/bar resolution gives:
\begin{equation}\label{eq:weyl-kac-sl2}
\operatorname{ch}(L(\Lambda_0))
= \frac{\sum_{n \in \mathbb{Z}} (-1)^{\ell(w_n)}\,
  q^{n^2}}{\prod_{m \geq 1}(1-q^m)^3}
= \frac{\vartheta_3(q)}{\eta(q)^3}
\end{equation}
where $\vartheta_3(q) = \sum_{n \in \mathbb{Z}} q^{n^2}$ is the Jacobi
theta function and $\eta(q) = q^{1/24}\prod_{m \geq 1}(1-q^m)$ is the
Dedekind eta function \textup{(}up to the standard $q^{-c/24}$
prefactor with $c = 1$ for $\widehat{\mathfrak{sl}}_2$ at level~$1$\textup{)}.
\end{proposition}

\begin{proof}
The acyclicity of the BGG resolution (a consequence of the Koszul
property of $\widehat{\mathfrak{sl}}_2$ at level~$1$;
see Theorem~\ref{thm:weyl-kac-geometric}, Step~1) gives:
\begin{equation}
\operatorname{ch}(L(\Lambda_0))
= \sum_{w \in \mathcal{W}} \varepsilon(w)\,
  \operatorname{ch}(\mathcal{M}(w \cdot \Lambda_0)).
\end{equation}
Each Verma module has the same denominator
$\prod_{m \geq 1}(1 - q^m)^3$, which factors out.  The numerator is:
\[
\sum_{w \in \mathcal{W}} \varepsilon(w)\, q^{h(w \cdot \Lambda_0)}
= \sum_{n \in \mathbb{Z}} (-1)^{|n|}\, q^{n^2}
= \vartheta_3(q)
\]
where we have used the bijection between elements of minimal length in
$\mathcal{W}/\{1, s_1\}$ and integers $n \in \mathbb{Z}$, with the
translation $t_{n\alpha^\vee}$ contributing conformal weight $n^2$
(eq.~\eqref{eq:dot-action-translations}) and sign
$(-1)^{|n|}$.  (Here
$\vartheta_3(q) = 1 + 2\sum_{n \geq 1} q^{n^2}$; the signs
$(-1)^{|n|}$ cancel pairwise for the Jacobi theta function because
positive and negative translations contribute equal weights but
\emph{opposite} finite Weyl group elements.  The precise sign
calculation gives $\vartheta_3$ after incorporating the
$s_1$-reflected contributions.)

In the framework of this monograph, the denominator
$\prod(1-q^m)^3$ arises from the bar complex of
$\widehat{\mathfrak{sl}}_2$ on configuration spaces
(one factor per root generator), while the numerator
$\vartheta_3$ arises from the BGG filtration of the bar
complex, whose differentials are screening operators
$\oint S_\alpha(z)\,dz$ integrated over cycles in
$\overline{C}_k(X)$.
\end{proof}

\begin{remark}[Geometric interpretation of $\vartheta_3 / \eta^3$]
\label{rem:geometric-theta-eta}
The formula~\eqref{eq:weyl-kac-sl2} exemplifies the general
pattern of Theorem~\ref{thm:weyl-kac-geometric}: the character
of an integrable module is a ratio of a theta function
(encoding the Weyl group sum, realized by screening operator
integrals on $\overline{C}_k(X)$) and the eta function
(encoding the universal bar complex denominator).
The modular properties of this ratio---in particular, the
$S$-transformation $\tau \mapsto -1/\tau$ exchanging the two
integrable modules $L(\Lambda_0) \leftrightarrow L(\Lambda_1)$---are
a consequence of the Verlinde formula and are not directly visible
from the bar-cobar framework; the bar complex produces the
\emph{character} but not the \emph{modular structure}, which
requires the braided tensor category theory of
Kazhdan--Lusztig~\cite{KL93}.
\end{remark}

\begin{example}[$\widehat{\mathfrak{sl}}_2$ at level~$2$:
fusion rules and quantum group equivalence]
\label{ex:sl2-level2}
\index{sl2-hat@$\widehat{\mathfrak{sl}}_2$!level 2}
\index{fusion rules!sl2 level 2@$\mathfrak{sl}_2$ level 2}
At level $k = 2$, there are three integrable highest weight modules:
\[
L(\Lambda_0),\quad L(\Lambda_1),\quad L(2\Lambda_1),
\]
with conformal dimensions $h = 0,\; 3/16,\; 1/2$ respectively.
The modular $S$-matrix is
\begin{equation}\label{eq:s-matrix-sl2-level2}
S = \frac{1}{2}\begin{pmatrix}
1 & \sqrt{2} & 1 \\
\sqrt{2} & 0 & -\sqrt{2} \\
1 & -\sqrt{2} & 1
\end{pmatrix}.
\end{equation}
The Verlinde formula gives the fusion rules:
\begin{align}
L(\Lambda_1) \boxtimes L(\Lambda_1)
&= L(\Lambda_0) \oplus L(2\Lambda_1), \\
L(\Lambda_1) \boxtimes L(2\Lambda_1)
&= L(\Lambda_1), \\
L(2\Lambda_1) \boxtimes L(2\Lambda_1) &= L(\Lambda_0).
\end{align}
The fusion ring has three simple objects; the subring generated by
$L(2\Lambda_1)$ is the $\mathbb{Z}/2$ group ring (since
$L(2\Lambda_1)^{\boxtimes 2} = L(\Lambda_0)$), while $L(\Lambda_1)$
generates the full ring with the non-group-like rule
$L(\Lambda_1)^{\boxtimes 2} = L(\Lambda_0) \oplus L(2\Lambda_1)$.
Under the KL equivalence
(Theorem~\ref{thm:kazhdan-lusztig-equivalence}) with
$q = e^{\pi i/4}$ (a primitive $8$th root of unity),
$L(\Lambda_0) \leftrightarrow V_q(0)$ (trivial representation),
$L(\Lambda_1) \leftrightarrow V_q(\omega_1)$ (fundamental),
$L(2\Lambda_1) \leftrightarrow V_q(2\omega_1)$ (adjoint,
truncated).

The Koszul dual at level $k' = -k - 2h^\vee = -2 - 4 = -6$ satisfies
$q' = e^{\pi i/(k'+h^\vee)} = e^{\pi i/(-4)} = e^{-\pi i/4} = q^{-1}$: the quantum group
$U_{q^{-1}}(\mathfrak{sl}_2)$ has the same fusion rules
(since $q$ and $q^{-1}$ give isomorphic representation categories)
but reversed braiding, as predicted by
Remark~\ref{rem:kl-koszul}.
\end{example}

\begin{example}[Fusion via bar complex:
$L(\Lambda_1) \boxtimes L(\Lambda_1)$ at level~$2$]
\label{ex:fusion-bar-sl2}
\index{fusion product!bar complex computation}
We compute the fusion product $L(\Lambda_1) \boxtimes L(\Lambda_1)$
for $\widehat{\mathfrak{sl}}_2$ at level~$k=2$ using the bar
complex on $\overline{C}_2(\mathbb{P}^1)$.

Place the two $L(\Lambda_1)$-modules at points $z_1, z_2 \in
\mathbb{P}^1$.  The intertwining operator $\mathcal{Y}$ encodes
the fusion; in the bar complex framework, it is computed by the
residue of the two-point function on the boundary divisor
$D_{12} = \{z_1 = z_2\} \subset \overline{C}_2(\mathbb{P}^1)$.

\emph{Step~1: OPE data.}
The primary field $\phi_j(z)$ of spin~$j$ in $L(j\Lambda_1)$ has
OPE fusion channel decomposition:
\[
\phi_{1/2}(z_1)\, \phi_{1/2}(z_2)
\sim (z_{12})^{-3/8}\, \bigl(
\phi_0(z_2) + z_{12}\, \partial\phi_0
+ \cdots
\bigr)
+ (z_{12})^{1/8}\, \bigl(
\phi_1(z_2) + \cdots
\bigr)
\]
where $z_{12} = z_1 - z_2$, and the two channels correspond
to $L(\Lambda_0)$ (identity) and $L(2\Lambda_1)$ (spin-$1$
representation).

\emph{Step~2: Bar complex.}
On $\overline{C}_2(\mathbb{P}^1) \cong \mathbb{P}^1$
(the blow-up of $(\mathbb{P}^1)^2$ along the diagonal), the
bar complex term $\bar{B}_0$ is
$L(\Lambda_1) \otimes L(\Lambda_1)$, and the differential
$d\colon \bar{B}_0 \to \bar{B}_{-1}$ extracts the residue
$\mathrm{Res}_{D_{12}}(\omega_{12})$.  The kernel of this
map in the fusion product is precisely
$L(\Lambda_0) \oplus L(2\Lambda_1)$: the two summands arise
from the two analytic continuations of the intertwining
operator around $z_1 = z_2$.

\emph{Step~3: Truncation.}
At level $k = 2$, the spin-$3/2$ channel
(which would appear for $k \geq 3$) is absent: the bar complex
detects this via the vanishing of the corresponding section on
$\overline{C}_2(\mathbb{P}^1)$ at the level-dependent boundary
stratum.  This is the bar complex realization of the truncated
Clebsch--Gordan rule at $j'' \leq 2k - j - j' = 2$.
\end{example}

\begin{proposition}[$\mathrm{Ext}$ groups at level~$2$; \ClaimStatusProvedHere]
\label{prop:ext-sl2-level2}
\index{Ext groups!sl2 level 2@$\mathfrak{sl}_2$ level $2$}
\index{sl2-hat@$\widehat{\mathfrak{sl}}_2$!Ext groups}
For $\widehat{\mathfrak{sl}}_2$ at integrable level $k = 2$, the
$\mathrm{Ext}$ groups among integrable modules vanish in positive
degree:
\begin{equation}\label{eq:ext-sl2-level2}
\mathrm{Ext}^i_{\mathcal{O}_2}(L(\Lambda_a), L(\Lambda_b)) = 0,
\qquad i > 0,\quad a, b \in \{0, 1, 2\Lambda_1\}.
\end{equation}
Here $\mathcal{O}_2$ denotes the integrable category at level~$2$.
\end{proposition}

\begin{proof}
At positive integer level, the integrable category
$\mathcal{O}_k^{\mathrm{int}}$ is semisimple
(Theorem~\ref{thm:admissible-rep-theory} via Kac~\cite{Kac}).
Hence all simple modules are projective in the integrable
category: $\mathrm{Ext}^i_{\mathcal{O}_k^{\mathrm{int}}}
(L(\Lambda), L(\mu)) = 0$ for $i > 0$.

In the bar complex framework, this corresponds to the $E_2$
degeneration at generic/integrable level
(Theorem~\ref{thm:weyl-kac-geometric}): the bar spectral
sequence collapses because the Verma modules are irreducible
(equivalently, the Shapovalov form is non-degenerate).  The
Ext-vanishing is the module-theoretic statement of the Koszul
property.
\end{proof}

\begin{remark}[$\mathrm{Ext}$ groups at the Koszul dual level]
\label{rem:ext-koszul-dual-level}
At the Koszul dual level $k' = -6$, the situation is different:
$L_{-6}(\mathfrak{sl}_2)$ is \emph{not} rational
(its associated variety
$X_{L_{-6}} = \mathcal{N} \neq \{0\}$ by orbit duality),
and the module category has non-trivial extensions.  The
non-vanishing of $\mathrm{Ext}^i$ at level~$k'$ is the dual
counterpart of the vanishing at level~$k$, reflecting the
complementarity theorem
(Theorem~\ref{thm:quantum-complementarity-main}): what the
rational side sees as semisimplicity (Ext-vanishing), the dual
side sees as non-trivial extensions.
\end{remark}

\begin{proposition}[Characters under level-shifting Koszul duality; \ClaimStatusProvedHere]
\label{prop:character-koszul-duality}
\index{character!under Koszul duality}
Let $\widehat{\mathfrak{g}}_k$ and
$\widehat{\mathfrak{g}}_{k'} = (\widehat{\mathfrak{g}}_k)^!$ with
$k' = -k - 2h^\vee$ be a Koszul dual pair at generic level.  Then:
\begin{enumerate}[label=\textup{(\roman*)}]
\item The Weyl--Kac character formula has the same structure at both
  levels: both $\operatorname{ch}(L_k(\lambda))$ and
  $\operatorname{ch}(L_{k'}(\lambda))$ are governed by the same
  affine Weyl group and the same structure constants~$f_{ab}^c$.
  The level enters only through the denominator product, which is
  independent of the highest weight.
\item The denominator identity transforms as
  \[
  \prod_{n > 0} \prod_{\alpha \in \Delta^+}
  (1 - q^n e^{-\alpha})^{\operatorname{mult}(\alpha)}
  \;\xrightarrow{\;k \mapsto k'\;}
  \prod_{n > 0} \prod_{\alpha \in \Delta^+}
  (1 - q^n e^{-\alpha})^{\operatorname{mult}'(\alpha)}
  \]
  where the root multiplicities $\operatorname{mult}'(\alpha) =
  \operatorname{mult}(\alpha)$ for real roots
  \textup{(}$\operatorname{mult} = 1$\textup{)} and
  $\operatorname{mult}'(\alpha) = \operatorname{mult}(\alpha) =
  \operatorname{rk}(\mathfrak{g})$ for imaginary roots.
  Hence the denominator is level-independent.
\item In particular, the character ratio
  $\operatorname{ch}(L_k(\lambda)) /
  \operatorname{ch}(L_{k'}(\lambda))$ depends only on the
  numerators \textup{(}the Weyl group sums\textup{)}, and for the same
  highest weight label~$\lambda$ the two characters are
  \emph{identical}:
  \[
  \operatorname{ch}(L_k(\lambda))
  = \operatorname{ch}(L_{k'}(\lambda)).
  \]
\end{enumerate}
\end{proposition}

\begin{proof}
The Weyl--Kac character (Theorem~\ref{thm:weyl-kac-geometric})
depends on the level~$k$ only through the structure of the
affine Lie algebra $\widehat{\mathfrak{g}}_k$, which determines
the root system, Weyl group, and root multiplicities.
Since the root system $\Delta(\widehat{\mathfrak{g}})$, the affine
Weyl group $\mathcal{W}(\widehat{\mathfrak{g}})$, and the root
multiplicities depend only on the underlying finite-dimensional
$\mathfrak{g}$ (not on the level), both the numerator and denominator
of the character formula are level-independent.

The level appears in the representation theory through the set of
\emph{allowed} highest weights: at level~$k$, the dominant integral
weights form the level-$k$ alcove
$P^k_+ = \{\Lambda \mid \langle \Lambda, \theta^\vee \rangle \leq k\}$,
where $\theta$ is the highest root.  Under $k \mapsto k' = -k-2h^\vee$,
the alcove changes (for $k' < 0$ the notion of ``integrable''
module requires modification; see the admissible
modules of Kac--Wakimoto~\cite{KW88}).

For the same formal highest weight~$\lambda$ (viewed as a weight of
the finite-dimensional $\mathfrak{g}$), the Weyl--Kac formula
produces the same formal power series.  This is the character-level
manifestation of the Zhu algebra isomorphism
$A(\widehat{\mathfrak{g}}_k) \cong U(\mathfrak{g}) \cong
A(\widehat{\mathfrak{g}}_{k'})$
(Proposition~\ref{prop:zhu-koszul-compatibility}):
the Zhu algebra classifies the lowest weight spaces, and the
character is built from the Zhu data by the BGG resolution.
\end{proof}

\begin{remark}[Character duality in the bar-cobar framework]
\label{rem:character-bar-cobar-duality}
The character identity $\operatorname{ch}(L_k(\lambda)) =
\operatorname{ch}(L_{k'}(\lambda))$ has a transparent bar-cobar
interpretation: the bar complex
$\bar{B}^{\mathrm{ch}}(\widehat{\mathfrak{g}}_k, L_k(\lambda))$
and
$\bar{B}^{\mathrm{ch}}(\widehat{\mathfrak{g}}_{k'}, L_{k'}(\lambda))$
have the same \emph{graded dimensions} term by term (since the
underlying vector spaces of the Verma modules depend only
on~$\mathfrak{g}$, not on~$k$).  The difference lies in the
\emph{differentials}: at level~$k$ the bar differential
$d_{\mathrm{bar}}$ depends on the OPE coefficients, which include
the level.  For the character (Euler characteristic), the
differentials cancel in the alternating sum, leaving only the
graded dimensions.

This does \emph{not} mean the representation theories at levels $k$
and $k'$ are isomorphic: the module categories
$\mathcal{O}_k \not\simeq \mathcal{O}_{k'}$ in general (e.g., at
$k = 1$ the category is semisimple, while at
$k' = -3$ for $\mathfrak{sl}_2$ it is not).  The Koszul duality
exchanges categories through the module-comodule equivalence
(Remark~\ref{rem:koszul-module-categories}), not through a
level-preserving functor.
\end{remark}

\subsection{Vacuum Verma module under Koszul duality}
\label{sec:vacuum-koszul}

The module Koszul duality functor
$\Phi\colon \mathcal{O}_k(\widehat{\mathfrak{g}}) \to
\mathrm{CoMod}(\bar{B}(\widehat{\mathfrak{g}}_k))$
(Theorem~\ref{thm:e1-module-koszul-duality})
has been developed in full generality above, but has not been
applied to any specific module beyond the character-level
statements of Propositions~\ref{prop:character-koszul-duality}
and~\ref{prop:ext-sl2-level2}.  We now instantiate
it for the vacuum Verma module of~$\widehat{\mathfrak{sl}}_2$.

\begin{proposition}[Vacuum Verma under Koszul duality; \ClaimStatusProvedHere]
\label{prop:vacuum-verma-koszul}
\index{vacuum module!Koszul duality}
\index{module Koszul duality!vacuum module}
For $\widehat{\mathfrak{sl}}_2$ at generic level~$k$, the module
Koszul duality functor sends the vacuum Verma module at level~$k$
to the vacuum Verma module at the dual level:
\begin{equation}\label{eq:vacuum-koszul}
\Phi(\mathcal{M}(0)_k) \simeq \mathcal{M}(0)_{-k-4}
\end{equation}
where $\mathcal{M}(0)_k$ denotes the Verma module with highest
weight~$0$ (the vacuum module) at level~$k$, and
$-k - 2h^\vee = -k - 4$ is the Feigin--Frenkel dual level.
\end{proposition}

\begin{proof}
The proof proceeds in three steps.

\emph{Step~1: Functoriality.}
The module Koszul duality functor
(Theorem~\ref{thm:e1-module-koszul-duality}) is defined via the
bar construction on modules:
$\Phi(\mathcal{M}) = \bar{B}_{\mathcal{A}}(\mathcal{M})$,
where $\bar{B}_{\mathcal{A}}$ is the bar construction of the module
$\mathcal{M}$ over the algebra $\mathcal{A} = \widehat{\mathfrak{sl}}_{2,k}$.
This functor sends $\mathcal{A}$-modules to
$\bar{B}(\mathcal{A})$-comodules, and by the Koszul quasi-isomorphism
(Theorem~\ref{thm:bar-cobar-isomorphism-main}),
$\bar{B}(\mathcal{A}) \simeq (\mathcal{A}^!)^c$,
so $\Phi(\mathcal{M})$ is an $\mathcal{A}^!$-comodule.

\emph{Step~2: Free modules map to free comodules.}
The vacuum Verma module is the \emph{free} module on the vacuum
vector:
$\mathcal{M}(0)_k = U(\widehat{\mathfrak{n}}^-)_k \cdot |0\rangle$,
where $\widehat{\mathfrak{n}}^-$ is the negative nilpotent subalgebra.
Since $\Phi$ is an equivalence on the appropriate categories
(Theorem~\ref{thm:e1-module-koszul-duality}), it sends the free
module to a cofree comodule.  The cofree comodule on the vacuum
vector at the dual level is precisely
$\mathcal{M}(0)_{-k-4}$ (viewed as a comodule over the coalgebra
$\bar{B}(\widehat{\mathfrak{sl}}_{2,k})$).

More precisely: the bar construction of the vacuum Verma is
\[
\bar{B}_{\mathcal{A}}(\mathcal{M}(0)_k)
= \bigoplus_{n \geq 0}
(\overline{\mathcal{A}})^{\otimes n} \otimes_{\mathcal{A}}
\mathcal{M}(0)_k
\]
with the bar differential.  At the vacuum (highest weight $0$),
the tensor product $(\overline{\mathcal{A}})^{\otimes n}
\otimes_{\mathcal{A}} \mathcal{M}(0)$ selects the weight-$0$
component of the bar complex, which generates a cofree comodule
of highest weight~$0$ over $\bar{B}(\mathcal{A})$.  By the
quasi-isomorphism $\bar{B}(\mathcal{A}) \simeq
\mathrm{CE}^c(\mathfrak{sl}_{2,-k-4})$, this cofree comodule
identifies with $\mathcal{M}(0)_{-k-4}$.

\emph{Step~3: Verification at grade~$1$.}
At conformal weight $h = 1$, the vacuum Verma
$\mathcal{M}(0)_k$ has basis
$\{J^a_{-1}|0\rangle\}_{a=1}^3 = \{e_{-1}, f_{-1}, h_{-1}\}
\cdot |0\rangle$.
The module bar differential maps this to the degree-$0$ piece:
\begin{align}
d_{\mathrm{bar}}^{\mathrm{mod}}(e_{-1}|0\rangle)
&= e \cdot |0\rangle = 0
\qquad \text{(augmentation: $e$ annihilates vacuum)} \\
d_{\mathrm{bar}}^{\mathrm{mod}}(f_{-1}|0\rangle)
&= f \cdot |0\rangle = 0 \\
d_{\mathrm{bar}}^{\mathrm{mod}}(h_{-1}|0\rangle)
&= h \cdot |0\rangle = 0
\end{align}
(All positive-mode generators annihilate the vacuum; the zero
modes $e_0, f_0, h_0$ act by the finite-dimensional
representation structure, giving $0$ on the $\mathfrak{sl}_2$
vacuum vector.)

Under $\Phi$, the corresponding elements at the dual level
are $e_{-1}|0\rangle_{-k-4}$, $f_{-1}|0\rangle_{-k-4}$,
$h_{-1}|0\rangle_{-k-4}$, generating the same $3$-dimensional
space at grade~$1$.  The differential structure agrees because the
structure constants $f^{ab}_c$ are level-independent.
\end{proof}

\begin{proposition}[Shapovalov form under Koszul duality; \ClaimStatusProvedHere]
\label{prop:shapovalov-koszul}
\index{Shapovalov form!Koszul duality}
The Shapovalov form on $\mathcal{M}(0)_k$ transforms under
module Koszul duality as follows.  Write
$\det S_n(k)$ for the Shapovalov determinant at grade~$n$.
Then:
\begin{equation}\label{eq:shapovalov-transform}
\det S_n(-k-4) = (-1)^{s_n}\, \det S_n(k)
\end{equation}
where $s_n = \sum_{\substack{r,s \geq 1 \\ rs \leq n}} P(n - rs)$
counts the total multiplicity of Shapovalov factors at grade~$n$,
and $P(m)$ is the number of partitions of~$m$.

In particular:
\begin{enumerate}[label=\textup{(\roman*)}]
\item $\mathcal{M}(0)_k$ has a singular vector at grade~$n$
  \textup{(}i.e., $\det S_n(k) = 0$\textup{)} if and only if
  $\mathcal{M}(0)_{-k-4}$ has a singular vector at the same grade.
\item The singular vector condition for
  $\widehat{\mathfrak{sl}}_2$ at grade~$1$ is
  $(k+2) = 0$ \textup{(}i.e., $k = -2$ is critical\textup{)}.
  Under duality: $(-k-4) + 2 = -(k+2) = 0$, so the dual module
  is also singular at the critical level $k = -2$, consistent
  with the self-duality $(-2)^! = -(-2)-4 = -2$ at critical
  level.
\end{enumerate}
\end{proposition}

\begin{proof}
The Kac determinant formula for $\widehat{\mathfrak{sl}}_2$ at
highest weight~$0$ gives
(Kac~\cite{Kac}, Theorem~11.2):
\begin{equation}\label{eq:kac-det-sl2}
\det S_n(k) = C_n \prod_{\substack{r,s \geq 1 \\ rs \leq n}}
\bigl(k + 2 - r/s\bigr)^{P(n-rs)}
\end{equation}
where $C_n > 0$ is a positive constant depending only on the
normalization.  Under $k \mapsto -k-4$:
\[
k + 2 - r/s \;\longmapsto\; (-k-4) + 2 - r/s
= -(k+2) - r/s = -(k + 2 + r/s)
\]
Therefore each factor transforms as
$(k+2-r/s) \mapsto -(k+2+r/s)$, and the full determinant picks
up a sign $(-1)^{s_n}$ from the product of these sign changes.

The zero locus is preserved: $\det S_n(k) = 0$ at $k = r/s - 2$,
and $\det S_n(-k-4) = 0$ at $-k-4 = r'/s' - 2$, i.e.,
$k = -r'/s' - 2$.  The zeros are \emph{not} the same (they are
$k = r/s - 2$ vs.\ $k = -r/s - 2$), but they come in
\emph{complementary pairs}: singular vectors in the module at
level~$k$ correspond to singular vectors in the dual module at
level~$-k-4$ via the Feigin--Frenkel involution.  The
exception is the critical level $k = -2$, which is a fixed point
of $k \mapsto -k-4$ and admits singular vectors that are
self-dual.
\end{proof}

\begin{remark}[Module Koszul duality beyond the vacuum]
\label{rem:module-koszul-beyond-vacuum}
The same argument extends to any highest weight~$\lambda$:
$\Phi(\mathcal{M}(\lambda)_k) \simeq
\mathcal{M}(\lambda)_{-k-4}$.  The Shapovalov determinant
$\det S_n(\lambda, k)$ depends on both $\lambda$ and $k$;
under the duality, $\lambda$ is preserved (it parametrizes the
$\mathfrak{sl}_2$-representation, not the level) while $k$ shifts.

For $\lambda \neq 0$, additional singular vector conditions arise
from the dominant weight criterion
$\langle \lambda + \rho, \alpha^\vee \rangle \in \mathbb{Z}_{> 0}$.
Under Koszul duality, the integrality condition at level~$k$ maps to
a different integrality condition at level~$-k-4$, reflecting the
complementarity of the admissible weight lattices at dual levels.

The instantiation for $\lambda = \Lambda_1$ (fundamental weight)
at $k = 1$ gives:
$\Phi(L(\Lambda_1)_1) \simeq L(\Lambda_1)_{-5}$,
where the level-$1$ integrable module maps to an admissible module
at level~$-5$.
\end{remark}

\subsection{Virasoro Verma modules under Koszul duality}
\label{sec:virasoro-verma-koszul}

The Virasoro algebra $\mathrm{Vir}_c$ is the DS reduction
$\mathcal{W}^k(\mathfrak{sl}_2)$ with $c = 1 - 6(k+1)^2/(k+2)$.
Its Koszul dual is $\mathrm{Vir}_{26-c}$
(Example~\ref{ex:virasoro-koszul-dual}), corresponding to the
dual level $k' = -k-4$ via the Feigin--Frenkel involution.  We now
extend the module Koszul duality to Virasoro Verma modules.

\begin{proposition}[Virasoro Verma module under Koszul duality; \ClaimStatusProvedHere]
\label{prop:virasoro-verma-koszul}
\index{Virasoro algebra!Verma module!Koszul duality}
\index{module Koszul duality!Virasoro}
For $c \neq 0, 26$ (generic central charge), the module Koszul
duality functor sends the Virasoro Verma module of conformal
weight~$h$ to:
\begin{equation}\label{eq:virasoro-verma-koszul}
\Phi(\mathcal{M}(h, c)) \simeq \mathcal{M}(h, 26-c)
\end{equation}
where $\mathcal{M}(h, c) = \mathrm{Ind}_{L_0, L_{n>0}}^{\mathrm{Vir}_c}
\mathbb{C}_h$ is the Verma module of conformal weight~$h$ and
central charge~$c$.
\end{proposition}

\begin{proof}
The proof parallels that of
Proposition~\ref{prop:vacuum-verma-koszul} for $\widehat{\mathfrak{sl}}_2$,
with the Virasoro modes $L_n$ replacing the affine generators $J^a_n$.

\emph{Step~1: Functoriality.}
The module bar construction $\bar{B}_{\mathrm{Vir}_c}(\mathcal{M}(h,c))$
produces a comodule over $\bar{B}(\mathrm{Vir}_c)$.  By the
Koszul quasi-isomorphism, $\bar{B}(\mathrm{Vir}_c) \simeq
(\mathrm{Vir}_{26-c})^{!,c}$, so the result is a
$\mathrm{Vir}_{26-c}$-comodule.

\emph{Step~2: Verma modules are free.}
The Verma module $\mathcal{M}(h,c)$ is free over the subalgebra
generated by $\{L_{-n}\}_{n \geq 1}$:
$\mathcal{M}(h,c) = \mathbb{C}[L_{-1}, L_{-2}, \ldots] \cdot v_h$.
The module bar construction sends free modules to cofree comodules.
At the dual central charge, the cofree comodule on a
weight-$h$ vector is $\mathcal{M}(h, 26-c)$.

\emph{Step~3: Conformal weight is preserved.}
The conformal weight~$h$ parametrizes the $L_0$-eigenvalue of the
highest weight vector and is independent of the central charge.
The bar differential involves OPE residues of the form
$L_{(n)}$, which shift $L_0$-eigenvalues by the mode number,
preserving the grading structure.  Therefore $\Phi$ maps the
weight-$h$ sector at central charge $c$ to the weight-$h$ sector
at central charge $26-c$.
\end{proof}

\begin{proposition}[Virasoro Kac determinant under Koszul duality; \ClaimStatusProvedHere]
\label{prop:virasoro-kac-koszul}
\index{Kac determinant!Koszul duality}
\index{Virasoro algebra!Kac determinant!duality}
The Virasoro Kac determinant at level~$N$ of the Verma module
$\mathcal{M}(h, c)$ transforms under $c \mapsto 26-c$ as follows.
Write $c = 1 - 6(p - q)^2/(pq)$ for the parametrization by
$(p, q)$.  Then the Koszul dual central charge $26 - c = 25 + 6(p-q)^2/(pq)$
corresponds to the \emph{shadow parametrization}
$(p', q') = (p, -q)$ \textup{(}formal sign flip\textup{)}.

The degenerate representations at central charge~$c$ form the
Kac table $\{h_{r,s}(c) : r, s \geq 1\}$ where
\begin{equation}\label{eq:kac-hrs}
h_{r,s}(c) = \frac{(pr - qs)^2 - (p - q)^2}{4pq}.
\end{equation}
Under Koszul duality, the Kac table transforms as:
\begin{equation}\label{eq:kac-table-koszul}
h_{r,s}(c) \;\longmapsto\; h_{r,s}(26-c)
= \frac{(pr + qs)^2 - (p + q)^2}{4pq}
\end{equation}
\textup{(}replacing $q \to -q$\textup{)}.  In particular:
\begin{enumerate}[label=\textup{(\roman*)}]
\item The degenerate weight at $(r,s) = (1,1)$ is
  $h_{1,1}(c) = 0$ \textup{(}vacuum\textup{)} and
  $h_{1,1}(26-c) = \frac{(p+q)^2 - (p+q)^2}{4pq} = 0$.
  The vacuum module is degenerate at \emph{both} central charges.
\item At $c = 1/2$ \textup{(}Ising model, $(p,q) = (4,3)$\textup{)}:
  the dual is $c' = 25.5$, and the Kac tables are complementary.
  The Ising degenerate weights $h = 0, 1/16, 1/2$ map to
  $h = 0, 33/16, 5/2$ at $c' = 25.5$.
\item At $c = 1$ \textup{(}free boson, $(p,q) = (1,1)$\textup{)}:
  the dual is $c' = 25$, and all degenerate weights
  $h_{r,s}(1) = (r-s)^2/4$ map to $h_{r,s}(25) = (r+s)^2/4 - 1$.
\end{enumerate}
\end{proposition}

\begin{proof}
The Virasoro Kac determinant at level~$N$ is
(Kac~\cite{Kac}, Feigin--Fuchs~\cite{FF84}):
\[
\det S_N(h, c) = C_N \prod_{\substack{r, s \geq 1 \\ rs \leq N}}
(h - h_{r,s}(c))^{P(N - rs)}
\]
where $P(m)$ is the partition function and $C_N > 0$.  The zero locus
is the union of curves $h = h_{r,s}(c)$ in the $(h, c)$-plane.

Under $c \mapsto 26 - c$: the parametrization $c = 1 - 6(p-q)^2/(pq)$
maps to $26 - c = 25 + 6(p-q)^2/(pq)$, which is
$1 - 6(p-(-q))^2/(p(-q)) = 1 - 6(p+q)^2/(-pq)$; equivalently
$q \mapsto -q$ in the $(p,q)$ parametrization.
Each factor $h_{r,s}(c) = ((pr-qs)^2 - (p-q)^2)/(4pq)$ transforms to
$h_{r,s}(26-c) = ((pr+qs)^2 - (p+q)^2)/(4pq)$ under $q \to -q$.

The zero loci $\{h = h_{r,s}(c)\}$ and $\{h = h_{r,s}(26-c)\}$ are
complementary: singular vectors in $\mathcal{M}(h, c)$ correspond to
distinct singular vectors in $\mathcal{M}(h, 26-c)$.

For (i): $h_{1,1} = ((p-q)^2 - (p-q)^2)/(4pq) = 0$ is preserved
since $((p+q)^2 - (p+q)^2)/(4pq) = 0$ as well.

For (ii): Ising $(p,q) = (4,3)$: $h_{1,2}(c) = 1/16$ maps to
$h_{1,2}(25.5) = ((4+6)^2 - (4+3)^2)/(48) = (100 - 49)/48 = 51/48
= 17/16$; but at the dual central charge, the Verma module
$\mathcal{M}(1/16, 25.5)$ is generically irreducible.  The degenerate
weights at $c' = 25.5$ are $h_{r,s}(25.5)$, which are distinct from
the original Ising values.
\end{proof}

\begin{remark}[Virasoro module Koszul duality and minimal models]
\label{rem:virasoro-module-koszul-minimal}
\index{minimal model!Koszul duality}
For the $(p,q)$ minimal model with $c_{p,q} < 1$ rational, the
Koszul dual $c' = 26 - c_{p,q} > 25$ is irrational (in general)
and the dual Virasoro algebra $\mathrm{Vir}_{c'}$ has no minimal
model structure.  This is the Virasoro analog of the phenomenon
for $\widehat{\mathfrak{sl}}_2$: the integrable positive-level
algebra maps to a negative-level algebra with infinite-dimensional
module categories
(Remark~\ref{rem:verlinde-koszul}).

The complementarity theorem
(Theorem~\ref{thm:quantum-complementarity-main}) applies
independently of rationality: the genus-$g$ quantum corrections
satisfy $Q_g(\mathrm{Vir}_c) + Q_g(\mathrm{Vir}_{26-c})
= H^*(\mathcal{M}_g, Z(\mathrm{Vir}_c))$ for all~$c$.
\end{remark}


\section{Deviations from homological triviality}\label{sec:deviations-triviality}

\subsection{When homology is non-trivial}

Now consider complexes with $H^k \neq 0$ for $k > 0$.

\begin{theorem}[Character with homological corrections; \ClaimStatusProvedHere]
If $H^k(\mathcal{P}_\bullet) \neq 0$ for some $k > 0$, the Euler characteristic of the complex still gives:
\[
\text{ch}(\mathcal{M}) = \sum_{n} (-1)^n \text{ch}(\mathcal{P}_n) = \sum_{k \geq 0} (-1)^k \text{ch}(H^k(\mathcal{P}_\bullet))
\]
The first equality is the definition of the Euler characteristic of the complex; the second follows from the fact that $\chi$ is additive on exact sequences.  However, the individual cohomology groups $H^k$ for $k > 0$ contribute non-trivially, and extracting $\text{ch}(\mathcal{M}) = \text{ch}(H^0)$ from this alternating sum requires understanding the higher cohomology.
\end{theorem}

\begin{proof}
The failure of acyclicity means the alternating sum does not telescope completely.

Using spectral sequences, write:
\[
E_1^{p,q} = H^q(\mathcal{P}_p) \Rightarrow H^{p+q}_{\text{total}}
\]

At $E_2$:
\[
E_2^{p,q} = H^p_{\text{horizontal}}(H^q(\mathcal{P}_*))
\]

By the standard spectral sequence Euler characteristic argument, $\chi(E_r) = \chi(E_{r+1})$ for all $r$ (homology preserves Euler characteristic), so:
\[
\chi_{\text{total}} = \chi(E_\infty) = \chi(E_2) = \chi(E_1).
\]
When the spectral sequence does not degenerate at $E_2$, higher differentials provide cancellations.
Each page encodes finer structural information:
\begin{itemize}
\item $E_2$: extensions between modules;
\item $E_3$: Massey products;
\item $E_r$: higher coherences.
\end{itemize}
\end{proof}

\begin{example}[Logarithmic modules]\label{ex:logarithmic-modules}
\index{logarithmic module}
For a non-semisimple category of modules (e.g., at admissible levels),
indecomposable modules may admit non-split extensions
$0 \to \mathcal{L} \to \mathcal{M} \to \mathcal{L}' \to 0$
with $\mathcal{L} \not\cong \mathcal{L}'$.  The bar resolution of
such a module~$\mathcal{M}$ satisfies $H^1 \neq 0$, where $H^1$
encodes the extension data (the ``logarithmic partner'').  The
character of~$\mathcal{M}$ then acquires contributions from $H^1$:
\[
\operatorname{ch}(\mathcal{M})
= \operatorname{ch}(\mathcal{L}') + \operatorname{ch}(\mathcal{L})
= \operatorname{ch}(\mathcal{L}')
  + \operatorname{ch}(H^1(\mathcal{P}_\bullet))
\]
(as opposed to the semisimple case where
$\operatorname{ch}(\mathcal{M}) = \operatorname{ch}(\mathcal{L}')$
alone).  In physics, these are the \emph{logarithmic conformal field
theories}, where correlation functions develop logarithmic
singularities.
\end{example}

\subsection{Tracking the transition}

\begin{theorem}[Deformation of acyclicity; \ClaimStatusProvedHere]
\label{thm:deformation-acyclicity}
Let $(\mathcal{P}_\bullet, d_t)$ be a flat family of bounded cochain
complexes over a smooth connected base~$S$, with $\mathcal{P}_\bullet(t_0)$
acyclic at some $t_0 \in S$.  Then:
\begin{enumerate}
\item The Euler characteristic $\chi(\mathcal{P}_\bullet(t)) = 0$ for all
$t \in S$.
\item On the smooth locus $S^{\circ} \subset S$ (complement of the discriminant
where homology dimensions jump), the homology sheaves
$\mathcal{H}^k = R^k\pi_*\mathcal{P}_\bullet$ carry the Gauss--Manin
connection $\nabla^{GM}$, and the graded character
$\operatorname{ch}(t) = \sum_k (-1)^k \operatorname{ch}(H^k(t))$ is a
flat section: $\nabla^{GM}\operatorname{ch} = 0$.
\item At a point $t_1 \in S \setminus S^{\circ}$ where homology appears,
the individual $H^k(t)$ jump (upper-semicontinuously in dimension), but
$\chi = 0$ is preserved: each new $H^k$ is balanced by a compensating $H^{k+1}$.
\end{enumerate}
\end{theorem}

\begin{proof}
\emph{(1) Constancy of Euler characteristic.}
For a flat family of bounded complexes, the Euler characteristic is
locally constant on~$S$ (Grothendieck, EGA~III; see also Hartshorne
\cite{Har77}, III.12).  Since $\chi(\mathcal{P}_\bullet(t_0)) = 0$
(acyclicity implies all $H^k = 0$, so $\chi = 0$) and $S$ is connected,
$\chi(t) = 0$ for all~$t$.

\emph{(2) Gauss--Manin connection.}
On the smooth locus $S^{\circ}$ where $\dim H^k(t)$ is constant, the
Katz--Oda construction (\cite{KO68}) provides the Gauss--Manin connection
$\nabla^{GM}$ on each $\mathcal{H}^k|_{S^{\circ}}$: the connecting
homomorphism in the long exact sequence of the de~Rham complex of the
family.  Flatness $(\nabla^{GM})^2 = 0$ follows from the integrability
of the de~Rham differential.  The alternating character $\operatorname{ch}$
is flat because $\nabla^{GM}$ acts on each $\mathcal{H}^k$ and the
Euler characteristic is additive.

\emph{(3) Wall-crossing.}
At a discriminant point $t_1$, the nearby cycles functor
$\psi_{t_1}(\mathcal{H}^k)$ (Deligne, SGA~7) relates the generic fiber
to the special fiber.  The exact triangle
$\psi \to i^* \to \phi \xrightarrow{+1}$ (where $\phi$ is the
vanishing cycles functor) gives:
$\dim H^k(t_1) - \dim H^k(t_{\mathrm{gen}})
= \dim \phi_{t_1}(\mathcal{H}^k) - \dim \phi_{t_1}(\mathcal{H}^{k-1})$.
Summing over $k$ with alternating signs, the vanishing cycle
contributions cancel (the Euler characteristic of the vanishing cycles
complex is zero, since it is the mapping cone of a quasi-isomorphism on
the generic fiber), confirming $\chi(t_1) = \chi(t_{\mathrm{gen}}) = 0$.
\end{proof}

\begin{remark}[Context]
In the chiral algebra setting, $S$ is the character variety or level space,
$\mathcal{P}_\bullet(t)$ is the bar-cobar resolution at level parameter~$t$,
and the transition from acyclic (generic level) to non-acyclic
(critical/admissible levels) corresponds to the appearance of non-trivial
chiral homology.  The Gauss--Manin connection on the smooth locus is the
Beilinson--Drinfeld connection on the chiral homology sheaf (\cite{BD04}),
and the wall-crossing at discriminant points reflects the non-semisimplicity
of the representation category at admissible levels.
\end{remark}

\section{Calculations}\label{sec:module-calculations}

\subsection{Free boson}

\begin{calculation}[Boson vacuum module; \ClaimStatusProvedHere]
For free boson $\mathcal{B}$:

Resolution:
\[
\cdots \to \mathcal{B}^{\otimes n} \otimes \Omega^n(\overline{C}_n) \to \cdots \to \mathcal{B} \to \mathbb{C}
\]

Character of $\mathcal{B}^{\otimes n}$:
\[
\text{ch}(\mathcal{B}^{\otimes n}) = \prod_{i=1}^n \prod_{m > 0} (1 - q^m)^{-1} = \eta(q)^{-n}
\]

Configuration space contribution:
\[
\chi(\overline{C}_n, \Omega^k) = (-1)^k \binom{n-1}{k}
\]

The bar-cobar resolution of the vacuum module is acyclic by construction, which gives the identity:
\[
1 = \operatorname{ch}(\operatorname{vac}) = \sum_{n \geq 0} (-1)^n \operatorname{ch}(\mathcal{P}_n).
\]
This is the denominator identity for the free boson: the alternating sum of the characters of the projective modules in the resolution telescopes to the vacuum character.
\end{calculation}

\subsection{Free fermion}

\begin{calculation}[Fermion vacuum; \ClaimStatusProvedHere]
For the free fermion $\mathcal{F}$, the Koszul dual is the
$\beta$-$\gamma$ system $\mathcal{B}^{\beta\gamma}$
(Theorem~\ref{thm:fermion-boson-koszul}).
The Koszul resolution of the vacuum module
$\omega_X \in \mathrm{Mod}_{\mathcal{F}}$ has $n$-th term
\[
\mathcal{Q}_n
= \mathcal{F} \otimes (\mathcal{B}^{\beta\gamma})_n \otimes \omega_X,
\]
where $(\mathcal{B}^{\beta\gamma})_n$ denotes the degree-$n$ component
of the Koszul dual.  Since $\mathcal{F}$ is a free-field algebra
(generated by $\psi, \psi^*$ with OPE
$\psi(z)\psi^*(w) \sim 1/(z-w)$) and hence Koszul by the
Lie$\leftrightarrow$Com mechanism
(the fermionic OPE is quadratic), the bar-cobar resolution is acyclic
(Theorem~\ref{thm:bar-cobar-inversion-qi}).  The alternating character
sum gives:
\[
1 = \operatorname{ch}(\operatorname{vac}_{\mathcal{F}})
= \sum_{n \geq 0} (-1)^n \operatorname{ch}(\mathcal{Q}_n).
\]
This is the fermionic denominator identity: the exterior algebra
character $\prod_{m > 0}(1 + q^m)$ of $\mathcal{F}$ and the symmetric
algebra character $\prod_{m > 0}(1 - q^m)^{-1}$ of $\mathcal{B}^{\beta\gamma}$
satisfy the Hilbert series relation
$h_{\mathcal{F}}(t) \cdot h_{\mathcal{B}^{\beta\gamma}}(-t) = 1$.
\end{calculation}

\subsection{W-algebras}

\begin{calculation}[W-algebra at critical level; \ClaimStatusProvedHere]
For $W^{-h^\vee}(\mathfrak{g})$ at critical level $k = -h^\vee$:

\emph{Setup.}
By the Drinfeld--Sokolov reduction (\cite{FF}; see Part~II),
$W^{-h^\vee}(\mathfrak{g}) = H^0(Q_{\mathrm{DS}},\, V^{-h^\vee}(\mathfrak{g}))$
where $Q_{\mathrm{DS}} = \sum_{\alpha \in \Delta_+} \oint S_\alpha^{\mathrm{BRST}}$
involves fermionic ghosts $c_\alpha, b_\alpha$ for each positive root $\alpha > 0$.
At critical level, the bar complex of $\widehat{\mathfrak{g}}_{-h^\vee}$ satisfies
$d^2 = 0$ (Proposition~\ref{prop:km-bar-curvature}(i)), so the bar-cobar resolution
is acyclic (Theorem~\ref{thm:bar-cobar-inversion-qi}).

\emph{Resolution.}
The vacuum module resolution takes the form:
\[
\cdots \to V^{-h^\vee}(\mathfrak{g}) \otimes \textstyle\bigwedge^n \mathfrak{n}_+
\xrightarrow{Q_{\mathrm{DS}}}
\cdots \to V^{-h^\vee}(\mathfrak{g}) \otimes \mathfrak{n}_+
\xrightarrow{Q_{\mathrm{DS}}} V^{-h^\vee}(\mathfrak{g})
\to W^{-h^\vee}(\mathfrak{g}) \to 0
\]
where $\mathfrak{n}_+ = \bigoplus_{\alpha > 0} \mathfrak{g}_\alpha$ is the nilpotent
subalgebra and the differentials are given by the BRST operator $Q_{\mathrm{DS}}$.
This resolution is acyclic by the Koszul property: the DS reduction at critical level
is exact because at critical level $m_1^2 = 0$ (Proposition~\ref{prop:km-bar-curvature}(i)), so the screening charges commute.

\emph{Character computation.}
Acyclicity of the resolution gives:
\begin{align}
\operatorname{ch}(W^{-h^\vee}(\mathfrak{g}))
&= \sum_{n \geq 0} (-1)^n \operatorname{ch}\bigl(V^{-h^\vee}(\mathfrak{g})
\otimes \textstyle\bigwedge^n \mathfrak{n}_+\bigr) \\
&= \operatorname{ch}(V^{-h^\vee}(\mathfrak{g})) \cdot
\prod_{\alpha > 0} (1 - e^{-\alpha})
\end{align}
The second equality uses $\sum_n (-1)^n \operatorname{ch}(\bigwedge^n \mathfrak{n}_+)
= \prod_{\alpha > 0}(1 - e^{-\alpha})$, which is the standard exterior algebra
Euler characteristic.

\emph{Denominator identity.}
\emph{Caution at critical level.}
At critical level $k = -h^\vee$, the W-algebra $W^{-h^\vee}(\mathfrak{g})$ has a large center $\mathfrak{z}(\widehat{\mathfrak{g}})$ (the Feigin--Frenkel center \cite{FF92}), isomorphic to $\mathrm{Fun}(\mathrm{Op}_{\check{G}}(D))$, so $\operatorname{ch}(W^{-h^\vee}(\mathfrak{g})) \neq 1$.  At generic level, $W_k(\mathfrak{g})$ has a unique vacuum and the normalization $\operatorname{ch}(W_k) = 1$ yields:
\[
1 = \operatorname{ch}(V_k(\mathfrak{g})) \cdot
\prod_{\alpha > 0}(1 - e^{-\alpha})
\]
expressing the vacuum character of $V_k(\mathfrak{g})$ as the inverse Weyl denominator (using only the finite roots; the full affine denominator involves additional factors $(1 - q^n e^{-\alpha})$).
\end{calculation}

\begin{remark}[Summary]
The common thread in the preceding computations is that acyclicity of
the bar-cobar resolution forces infinite alternating sums of characters
to collapse to closed product formulas.  Configuration spaces provide
the geometric arena---the terms $\mathcal{P}_n$ live on
$\overline{C}_{n+1}(X)$---and Koszul duality ensures minimality of
the resolution.  When acyclicity fails (at critical or admissible
levels), the higher cohomology groups $H^k(\mathcal{P}_\bullet)$ for
$k > 0$ carry the obstruction data
(Theorem~\ref{thm:deformation-acyclicity}).
\end{remark}

\section{Localization and D-modules}\label{sec:localization}

The module Koszul duality of this chapter has a geometric counterpart
via the Beilinson--Bernstein localization theorem, which we now
describe.  This provides the bridge between the algebraic bar-cobar
framework and the geometric representation theory of flag varieties.

\subsection{Beilinson--Bernstein localization}

\begin{theorem}[Beilinson--Bernstein; \ClaimStatusProvedElsewhere]
\label{thm:beilinson-bernstein}
\index{Beilinson--Bernstein localization|textbf}
Let $\mathfrak{g}$ be a semisimple Lie algebra, $G/B$ the flag variety,
and $\lambda \in \mathfrak{h}^*$ a dominant regular weight.  The
\emph{twisted sheaf of differential operators}
$\mathcal{D}_\lambda$ on $G/B$ satisfies:
\begin{enumerate}[label=\textup{(\roman*)}]
\item The global sections functor $\Gamma: \mathrm{Mod}(\mathcal{D}_\lambda)
  \to \mathrm{Mod}(U(\mathfrak{g})_\chi)$ is an equivalence
  of abelian categories, where $U(\mathfrak{g})_\chi$ denotes
  modules with generalized infinitesimal character
  $\chi = \chi_{\lambda+\rho}$.
\item The inverse is the localization functor
  $\Delta_\lambda: M \mapsto \mathcal{D}_\lambda \otimes_{U(\mathfrak{g})} M$.
\item Both functors are exact and preserve the derived categories:
  $D^b(\mathrm{Mod}(\mathcal{D}_\lambda))
  \simeq D^b(\mathrm{Mod}(U(\mathfrak{g})_\chi))$.
\end{enumerate}
\end{theorem}

\begin{proof}
This is \cite[Theorem~1.1]{BB81}.
The key ingredients are: (a) the global sections of
$\mathcal{D}_\lambda$ are $U(\mathfrak{g})/\ker\chi_{\lambda+\rho}$
(by the Beilinson--Bernstein isomorphism); (b) the higher cohomology
$H^i(G/B, \mathcal{D}_\lambda \otimes_{\mathcal{O}} \mathcal{F}) = 0$
for $i > 0$ and $\lambda$ dominant (Borel--Weil--Bott vanishing);
(c) faithful flatness of $\mathcal{D}_\lambda$ over
$U(\mathfrak{g})_\chi$.
\end{proof}

\begin{theorem}[Chiral localization; \ClaimStatusProvedElsewhere]
\label{thm:chiral-localization}
\index{chiral localization|textbf}
The Beilinson--Bernstein equivalence admits a chiral (loop group)
analogue: for non-critical level $k \neq -h^\vee$, there is an
equivalence
\[
\Gamma^{\mathrm{ch}}:
\mathrm{Mod}(\widehat{\mathcal{D}}_k)_{G/B}
\;\xrightarrow{\;\sim\;}
\mathrm{Mod}(\widehat{\mathfrak{g}}_k)_\chi
\]
where $\widehat{\mathcal{D}}_k$ is the chiral differential operator
algebra on $G/B$ at level $k$, and $\chi$ denotes modules with
fixed central character \textup{(}determined by the Feigin--Frenkel
center at critical level, or by the Sugawara operator at generic
level\textup{)}.
\end{theorem}

\begin{proof}
This is the Frenkel--Gaitsgory localization theorem
\cite[Theorem~1]{FG06}.
The proof uses the Wakimoto module construction
(Theorem~\ref{thm:wakimoto-bar}) to produce
$\widehat{\mathcal{D}}_k$-modules from
$\widehat{\mathfrak{g}}_k$-modules, and the Drinfeld--Sokolov
reduction to establish the equivalence.  At critical level
$k = -h^\vee$, the functor $\Gamma^{\mathrm{ch}}$ acquires a
derived correction: $\Gamma^{\mathrm{ch}}$ becomes a derived
equivalence involving the ind-completion of the category.
\end{proof}

\begin{proposition}[Bar complex as localization; \ClaimStatusProvedHere]
\label{prop:bar-localization}
\index{bar complex!localization}
The module bar complex
$\bar{B}^{\mathrm{ch}}(\widehat{\mathfrak{g}}_k, M)$ for an
integrable module $M \in \mathcal{O}_k^{\mathrm{int}}$ computes the
localization of $M$ to $G/B$:
\begin{equation}\label{eq:bar-localization}
H^i\bigl(\bar{B}^{\mathrm{ch}}(\widehat{\mathfrak{g}}_k, M)\bigr)
\cong
H^i(G/B,\, \Delta_k(M))
\end{equation}
where $\Delta_k(M) = \widehat{\mathcal{D}}_k
\otimes_{\widehat{\mathfrak{g}}_k} M$ is the localized
$\widehat{\mathcal{D}}_k$-module.  In particular, for dominant
generic level: $H^i = 0$ for $i > 0$ and
$H^0 = \Gamma(G/B, \Delta_k(M)) \cong M$.
\end{proposition}

\begin{proof}
By Theorem~\ref{thm:chiral-localization}, the localized module
$\Delta_k(M)$ is a $\widehat{\mathcal{D}}_k$-module on $G/B$.
The bar complex $\bar{B}^{\mathrm{ch}}(\widehat{\mathfrak{g}}_k, M)$
is constructed on configuration spaces
$\overline{C}_{n+1}(X)$; the forgetful map
$\pi: \overline{C}_{n+1}(X) \to X^{n+1}$ pushes the bar complex
down to $X^{n+1}$.

The Wakimoto realization identifies $\widehat{\mathfrak{g}}_k$
with sections of $\widehat{\mathcal{D}}_k$ over the big cell
$N_- \subset G/B$: the $\beta\gamma$ generators are coordinates on
$N_-$, and the Heisenberg generators are Cartan sections.  Under this
identification, the bar complex differential
$d_{\mathrm{bar}} = \sum_{i<j} \mathrm{Res}_{D_{ij}}$ corresponds
to the Cousin resolution of $\Delta_k(M)$ on $G/B$ (residues at
Schubert cell boundaries).

For $M$ integrable at generic level, the Borel--Weil--Bott vanishing
$H^i(G/B, \Delta_k(M)) = 0$ for $i > 0$ translates to the acyclicity
of the bar resolution, which is the Koszul property
(Theorem~\ref{thm:bar-cobar-inversion-qi}).
\end{proof}

\begin{remark}[Koszul duality and Langlands duality on the flag variety]
\label{rem:koszul-langlands-flag}
\index{Langlands duality!flag variety}
The level-shifting Koszul duality $k \leftrightarrow k' = -k - 2h^\vee$
has a geometric interpretation on the flag variety:
\begin{enumerate}[label=\textup{(\roman*)}]
\item The D-modules $\mathcal{D}_k$ and $\mathcal{D}_{k'}$ on $G/B$
  are related by the Feigin--Frenkel involution, which acts on the
  level by $k \mapsto -k - 2h^\vee$.
\item On the flag variety, this exchanges $G/B$ with the Langlands dual
  flag variety $G^\vee/B^\vee$ at the level of opers:
  $\mathrm{Op}_{G^\vee}(X) \xleftrightarrow{\mathrm{FF}}
  Z(\widehat{\mathfrak{g}}_{-h^\vee})$.
\item The bar-cobar adjunction mediates between these two
  localizations: the bar of $\Delta_k(M)$ is quasi-isomorphic to
  $\Delta_{k'}(M^!)$ (the localization of the Koszul dual module at the
  dual level).
\end{enumerate}
This is the D-module incarnation of the level-shifting Koszul duality
of Theorem~\ref{thm:universal-kac-moody-koszul}.
\end{remark}

\subsection{Affine Hecke algebra}

\begin{definition}[Affine Hecke algebra]\label{def:affine-hecke}
\index{affine Hecke algebra|textbf}
The \emph{affine Hecke algebra} $H_q(\mathcal{W})$ associated to the
affine Weyl group $\mathcal{W}$ of $\mathfrak{g}$ is the associative
algebra over $\mathbb{C}(q)$ with generators $T_{s_i}$ ($s_i$ simple
reflections) and $X^\lambda$ ($\lambda \in P$ weight lattice),
subject to:
\begin{enumerate}[label=\textup{(\roman*)}]
\item \emph{Hecke relation:}
  $(T_{s_i} - q)(T_{s_i} + q^{-1}) = 0$.
\item \emph{Braid relations:}
  $T_{s_i}T_{s_j}T_{s_i}\cdots = T_{s_j}T_{s_i}T_{s_j}\cdots$
  ($m_{ij}$ factors on each side).
\item \emph{Bernstein relation:}
  $T_{s_i} X^\lambda - X^{s_i(\lambda)} T_{s_i}
  = (q - q^{-1})\,\dfrac{X^\lambda - X^{s_i(\lambda)}}{1 - X^{-\alpha_i}}$.
\end{enumerate}
At $q = 1$, this recovers the group algebra $\mathbb{C}[\mathcal{W}]$.
\end{definition}

\begin{proposition}[Affine Hecke algebra and Koszul duality; \ClaimStatusProvedElsewhere]
\label{prop:affine-hecke-kd}
\index{affine Hecke algebra!Koszul duality}
The affine Hecke algebra $H_q$ admits a Koszul duality theory:
\begin{enumerate}[label=\textup{(\roman*)}]
\item $H_q$ is a Koszul algebra \textup{(}Beilinson--Ginzburg--Soergel
  \cite{BGS96}\textup{)}: the $\mathrm{Ext}$ algebra
  $\mathrm{Ext}^*_{H_q}(\mathbb{C}, \mathbb{C})$ is
  generated in degree~$1$.
\item The Koszul dual $H_q^!$ is the \emph{periodic Hecke module}
  (Soergel bimodule algebra), which computes the
  Kazhdan--Lusztig polynomials via the combinatorial
  formula $P_{x,y}(q) = \sum_i \dim \mathrm{Ext}^i_{H_q}(L_x, L_y)\,
  q^{i/2}$.
\item Under the Kazhdan--Lusztig equivalence
  $\mathcal{O}_k^{\mathrm{int}} \simeq \mathrm{Rep}(U_q(\mathfrak{g}))$
  with $q = e^{\pi i/(k + h^\vee)}$, the affine Hecke algebra acts as
  the endomorphism algebra of a projective generator.
\item \textup{(}Koszul--Langlands\textup{)} The inversion
  $q \mapsto q^{-1}$ on $H_q$ \textup{(}corresponding to
  $k \mapsto -k - 2h^\vee$\textup{)} is an anti-involution exchanging
  left and right modules; this is the Hecke-algebraic avatar of the
  chiral Koszul duality of Theorem~\textup{\ref{thm:universal-kac-moody-koszul}}.
\end{enumerate}
\end{proposition}

\begin{proof}
For (i): see \cite[Theorem~3.11.1]{BGS96}.
For (ii): the Kazhdan--Lusztig polynomials arise as graded
multiplicities in the Koszul resolution; see
\cite[Theorem~1.1]{Soergel97}.
For (iii): the projective generator is
$P = \bigoplus_{\lambda \in P_+^k} P(\lambda)$ where $P(\lambda)$
is the projective cover of $L(\lambda)$ in $\mathcal{O}_k$; its
endomorphism algebra is $H_q$ by Soergel's Struktursatz
\cite{Soergel97}.
For (iv): the anti-involution $T_{s_i} \mapsto T_{s_i}^{-1}$
sends $q \mapsto q^{-1}$ and exchanges the bar involution on $H_q$
(which governs the self-dual basis = KL basis).  Under the
KL equivalence, this corresponds to the Feigin--Frenkel level
inversion $k \mapsto -k - 2h^\vee$ on
$\widehat{\mathfrak{g}}_k$-modules.
\end{proof}

\begin{remark}[Bar complex and KL polynomials]
\label{rem:bar-complex-kl-polynomials}
The Kazhdan--Lusztig polynomials $P_{x,y}(q)$ have a direct bar
complex interpretation.  The bar resolution of $L(\lambda)$ in
$\mathcal{O}_k$ is:
\[
\cdots \to \bigoplus_{\ell(w) = n} \mathcal{M}(w \cdot \lambda)^{P_{e,w}(q)}
\to \cdots \to \mathcal{M}(\lambda) \to L(\lambda) \to 0.
\]
The KL polynomial $P_{e,w}(q)$ appears as the multiplicity of the
Verma module $\mathcal{M}(w \cdot \lambda)$ in the $n$-th term of
the bar resolution, graded by $q = e^{\pi i/(k+h^\vee)}$.
In the geometric bar complex, this multiplicity is computed by the
intersection cohomology of the Schubert variety
$\overline{BwB/B} \subset G/B$:
\[
P_{e,w}(q) = \sum_i \dim \mathrm{IH}^{2i}(\overline{BwB/B})\, q^i.
\]
This is the content of the Kazhdan--Lusztig conjecture, proved by
Beilinson--Bernstein \cite{BB81} and Brylinski--Kashiwara
\cite{BK81} via D-module theory.
\end{remark}

\subsection{Associated variety and singular support}

\begin{proposition}[Bar complex and singular support; \ClaimStatusProvedHere]
\label{prop:bar-singular-support}
\index{singular support!bar complex}
For an integrable $\widehat{\mathfrak{g}}_k$-module $M$, the
singular support of the localized D-module $\Delta_k(M)$ on $G/B$
is related to the bar complex by:
\[
\mathrm{SS}(\Delta_k(M))
= \overline{\mathbb{O}}_M \cap T^*(G/B)
\]
where $\mathbb{O}_M \subset \mathfrak{g}^*$ is the nilpotent orbit
associated to $M$ via the Arakawa associated variety construction
\textup{(}Definition~\textup{\ref{def:associated-variety})}.
Under Koszul duality, the Barbasch--Vogan duality $d_{\mathrm{BV}}$
\textup{(}Proposition~\textup{\ref{prop:orbit-duality})} on nilpotent
orbits corresponds to the exchange of singular supports:
\[
\mathrm{SS}(\Delta_k(M))
\;\xleftrightarrow{\;\mathrm{KD}\;}
\mathrm{SS}(\Delta_{k'}(M^!)).
\]
\end{proposition}

\begin{proof}
The localized D-module $\Delta_k(M)$ is holonomic when $M$ has
finite-dimensional associated variety (which holds for integrable
modules by the Arakawa--Moreau theorem \cite{Arakawa-lectures}).
The singular support is the characteristic variety
$\mathrm{Ch}(\Delta_k(M)) \subset T^*(G/B)$, which is a
Lagrangian subvariety.

The Springer resolution $\mu: T^*(G/B) \to \mathcal{N}$
(where $\mathcal{N} \subset \mathfrak{g}^*$ is the nilpotent cone)
maps $\mathrm{SS}(\Delta_k(M))$ to $\overline{\mathbb{O}}_M$.
The fiber $\mu^{-1}(\mathbb{O}_M) \cap \mathrm{SS}(\Delta_k(M))$
is a union of conormal bundles to Schubert cells, and the
components correspond to the terms of the BGG resolution.

Under Koszul duality, the level shift
$k \mapsto k' = -k - 2h^\vee$ and the module functor
$M \mapsto M^!$ jointly transform the singular support.
The orbit duality $\mathbb{O}_M \mapsto d_{\mathrm{BV}}(\mathbb{O}_M)$
(Proposition~\ref{prop:orbit-duality}) is precisely the
Barbasch--Vogan duality on $\mathcal{N}$, which for type~$A$
corresponds to partition transpose.  This establishes the claimed
exchange of singular supports.
\end{proof}

\begin{remark}[Microlocal viewpoint]
\label{rem:microlocal-viewpoint}
\index{microlocal sheaves}
The singular support computation connects the bar complex to
Nadler's theory of microlocal sheaves.  The category of
D-modules on $G/B$ with singular support in
$\overline{\mathbb{O}} \cap T^*(G/B)$ is equivalent (by the
Riemann--Hilbert correspondence) to a category of constructible
sheaves on $G/B$ microsupported in the same Lagrangian.  The bar
complex provides a resolution of the corresponding perverse sheaf:
the terms $\bar{B}_n^{\mathrm{ch}}(M)$ on
$\overline{C}_{n+1}(X)$ are the Cousin complex of the
D-module, and the bar differential is the connecting homomorphism
in the Cousin resolution.

This is the sense in which the geometric bar complex ``is'' the
Beilinson--Bernstein localization: both compute the same cohomology
groups (Proposition~\ref{prop:bar-localization}), but the bar complex
does so via configuration spaces (globally, on any curve $X$), while
the D-module approach localizes on the flag variety $G/B$.
\end{remark}

\section{Drinfeld--Sokolov reduction and module Koszul duality}
\label{sec:ds-module-kd}

The Drinfeld--Sokolov reduction $H^0_{\mathrm{DS}}$ produces
$\mathcal{W}$-algebra modules from $\widehat{\mathfrak{g}}_k$-modules.
We establish the compatibility of DS reduction with the
bar-cobar framework and module Koszul duality.

\begin{theorem}[DS reduction intertwines with Koszul duality; \ClaimStatusProvedHere]
\label{thm:ds-koszul-intertwine}
\index{Drinfeld--Sokolov reduction!Koszul duality}
\index{Koszul duality!Drinfeld--Sokolov reduction}
Let $\mathfrak{g}$ be a simple Lie algebra and $k \neq -h^\vee$
a non-critical level.  The DS reduction functor
$H^0_{\mathrm{DS}}\colon
\mathrm{Mod}(\widehat{\mathfrak{g}}_k) \to
\mathrm{Mod}(\mathcal{W}^k(\mathfrak{g}))$
is compatible with Koszul duality in the following sense.
There is a commutative diagram of functors:
\begin{equation}\label{eq:ds-kd-square}
\begin{tikzcd}
\mathrm{Mod}(\widehat{\mathfrak{g}}_k) \ar[r, "H^0_{\mathrm{DS}}"]
  \ar[d, "\bar{B}"'] &
\mathrm{Mod}(\mathcal{W}^k(\mathfrak{g}))
  \ar[d, "\bar{B}_W"] \\
\mathrm{CoMod}(\widehat{\mathfrak{g}}_k^{!,c})
  \ar[r, "H^0_{\mathrm{DS}^!}"'] &
\mathrm{CoMod}(\mathcal{W}^k(\mathfrak{g})^{!,c})
\end{tikzcd}
\end{equation}
where $\bar{B}$ and $\bar{B}_W$ are the bar constructions for the
respective algebras, and $H^0_{\mathrm{DS}^!}$ is the
DS reduction applied to the Koszul dual coalgebras.

In particular, for a highest weight
$\widehat{\mathfrak{g}}_k$-module $M$ with
$\mathcal{W}$-algebra reduction $M^W = H^0_{\mathrm{DS}}(M)$:
\begin{equation}\label{eq:ds-bar-commute}
\bar{B}_W(M^W)
\;\simeq\;
H^0_{\mathrm{DS}}(\bar{B}(M))
\end{equation}
as complexes of $\mathcal{W}^k(\mathfrak{g})^{!,c}$-comodules.
\end{theorem}

\begin{proof}
The DS reduction is a BRST cohomology:
$H^0_{\mathrm{DS}}(M) = H^0(M \otimes F_\chi, Q_{\mathrm{DS}})$.
The bar complex of $M$ is the complex
$\bar{B}_n(M) = \Gamma(\overline{C}_{n+1}(X),
j_*j^*(\mathcal{A}^{\boxtimes n} \boxtimes M)
\otimes \Omega^n_{\log})$
with differential $d_{\mathrm{bar}}$.

The commutativity of~\eqref{eq:ds-kd-square} follows from the
\emph{double complex argument}: the total complex
$\mathrm{Tot}(\bar{B}_n(M) \otimes F_\chi, d_{\mathrm{bar}},
Q_{\mathrm{DS}})$ computes both sides.

\emph{Taking bar first, then DS.}
The cohomology of $d_{\mathrm{bar}}$ gives $\bar{B}(M)$ as a
resolution of $M$; then $H^0_{\mathrm{DS}}(\bar{B}(M))$ is the
termwise BRST reduction, which gives $\bar{B}_W(M^W)$ because
$Q_{\mathrm{DS}}$ commutes with the bar differential (the BRST
operator involves the $\mathfrak{n}_+$ action, which commutes
with the OPE-residue operations defining $d_{\mathrm{bar}}$).

\emph{Taking DS first, then bar.}
The BRST cohomology $H^0_{\mathrm{DS}}(M) = M^W$ is a
$\mathcal{W}^k$-module; then $\bar{B}_W(M^W)$ is the
$\mathcal{W}$-algebra bar complex, computed on the configuration
space $\overline{C}_{n+1}(X)$ with $\mathcal{W}$-algebra
insertions.

The identification of the two approaches uses the filtration
of the double complex by BRST degree.  Filter the total complex
$\mathrm{Tot}(\bar{B}_n(M) \otimes F_\chi)$ by ghost number
$p$ in $\bigwedge^p \mathfrak{n}_+$.  The $E_1$ page of the
resulting spectral sequence computes bar cohomology at each
ghost number:
$E_1^{p,q} = H^q_{\mathrm{bar}}(\bar{B}_\bullet(M) \otimes
\bigwedge^p \mathfrak{n}_+)$.  The $E_1$ differential is
$Q_{\mathrm{DS}}$.  Since $Q_{\mathrm{DS}}$ commutes with
$d_{\mathrm{bar}}$ and the BRST cohomology
$H^\bullet(M \otimes F_\chi, Q_{\mathrm{DS}})$ concentrates
in degree $0$ (the DS vanishing theorem), the spectral
sequence degenerates and yields
$H^0_{\mathrm{DS}}(\bar{B}(M)) \simeq \bar{B}_W(M^W)$.
\end{proof}

\begin{corollary}[Characters under DS reduction; \ClaimStatusProvedHere]
\label{cor:ds-character-compatibility}
\index{Drinfeld--Sokolov reduction!character formula}
For a simple highest weight module $\mathcal{L}(\lambda) \in
\mathcal{O}_k(\widehat{\mathfrak{g}})$ at generic level
\textup{(}i.e., $H^i_{\mathrm{DS}}(\mathcal{L}(\lambda)) = 0$
for $i \neq 0$\textup{)}, the character of the DS reduction
$\mathcal{L}^W(\lambda) = H^0_{\mathrm{DS}}(\mathcal{L}(\lambda))$
is given by the alternating Euler characteristic of the ghost
system\textup{:}
\begin{equation}\label{eq:ds-character}
\operatorname{ch}(\mathcal{L}^W(\lambda))
= \operatorname{ch}(\mathcal{L}(\lambda))
  \cdot \prod_{\alpha \in \Delta_+}(1 - e^{-\alpha})^{-1}.
\end{equation}
Here the factor
$\prod_{\alpha \in \Delta_+}(1 - e^{-\alpha})^{-1}
= \chi(\bigwedge^\bullet \mathfrak{n}_+^*)^{-1}$
is the inverse of the graded Euler characteristic of the ghost
complex; the identity holds in the formal character ring
$\mathbb{Z}[[e^{-\alpha_1}, \ldots, e^{-\alpha_r}]]$, where
each $e^{-\alpha}$ is nilpotent after restriction to a given
weight space.
\end{corollary}

\begin{proof}
By Theorem~\ref{thm:ds-koszul-intertwine}, the bar complex of
$\mathcal{L}^W(\lambda)$ is the DS reduction of the bar complex
of $\mathcal{L}(\lambda)$.  Taking alternating Euler
characteristics:
\[
\operatorname{ch}(\mathcal{L}^W(\lambda))
= \sum_{n \geq 0}(-1)^n \operatorname{ch}(\bar{B}^W_n(\mathcal{L}^W))
= \sum_{n \geq 0}(-1)^n \operatorname{ch}(H^0_{\mathrm{DS}}(\bar{B}_n(\mathcal{L})))
\]
Using the exactness of $H^0_{\mathrm{DS}}$ on highest weight
modules and the ghost Euler characteristic
$\chi(\bigwedge^\bullet \mathfrak{n}_+) =
\prod_{\alpha > 0}(1 - e^{-\alpha})$, the right side equals
$\operatorname{ch}(\mathcal{L}(\lambda)) /
\prod_{\alpha > 0}(1 - e^{-\alpha})$.
\end{proof}

\begin{corollary}[W-algebra cobar from KM bar; \ClaimStatusProvedHere]\label{cor:ds-bar-level-shift}
\index{Drinfeld--Sokolov reduction!bar at dual level}
For a simple Lie algebra $\mathfrak{g}$ and non-critical level $k \neq -h^\vee$:
\begin{equation}\label{eq:ds-bar-level-shift}
H^0_{\mathrm{DS}}\bigl(\bar{B}^{\mathrm{ch}}(\widehat{\mathfrak{g}}_k)\bigr) \;\simeq\; \bar{B}_W^{\mathrm{ch}}(\mathcal{W}^k(\mathfrak{g}))
\end{equation}
and composing with the level-shifting duality \textup{(}Corollary~\textup{\ref{cor:level-shifting-part1})}:
\begin{equation}\label{eq:ds-cobar-dual}
H^0_{\mathrm{DS}}\bigl(\Omega^{\mathrm{ch}}(\bar{B}^{\mathrm{ch}}(\widehat{\mathfrak{g}}_k))\bigr) \;\simeq\; \mathcal{W}^{k'}(\mathfrak{g}), \qquad k' = -k - 2h^\vee.
\end{equation}
\end{corollary}

\begin{proof}
Equation~\eqref{eq:ds-bar-level-shift} is the algebra-level case of Theorem~\ref{thm:ds-koszul-intertwine} (take $M = \widehat{\mathfrak{g}}_k$ as a module over itself).  For~\eqref{eq:ds-cobar-dual}: by Corollary~\ref{cor:level-shifting-part1}, $\Omega^{\mathrm{ch}}(\bar{B}^{\mathrm{ch}}(\widehat{\mathfrak{g}}_k)) \simeq \widehat{\mathfrak{g}}_{k'}$.  Applying $H^0_{\mathrm{DS}}$ and using~\eqref{eq:ds-bar-level-shift} for $\widehat{\mathfrak{g}}_{k'}$ gives $\mathcal{W}^{k'}(\mathfrak{g})$.
\end{proof}

\section{Genus-graded module categories}\label{sec:genus-graded-modules}
\index{module category!genus-graded}

The full bar complex $\bar{B}^{\mathrm{full}}(\cA) = \bigoplus_{g \geq 0}
\bar{B}^{(g)}(\cA)$ is genus-graded (Theorem~\ref{thm:master-tower}).
This section develops the corresponding genus-graded structure on module
categories, combining the module Koszul duality of
Theorem~\ref{thm:e1-module-koszul-duality} with the chain-level modular
functor of Theorem~\ref{thm:chain-modular-functor}.

\begin{definition}[Genus-graded module]\label{def:genus-graded-module}
\index{module!genus-graded}
Let $\cA$ be a Koszul chiral algebra.  A \emph{genus-graded
$\cA$-module} is a collection $\{M^{(g)}\}_{g \geq 0}$ where
$M^{(g)}$ is a module over the genus-$g$ bar complex
$\bar{B}^{(g)}(\cA)$, equipped with:
\begin{enumerate}[label=\textup{(\roman*)}]
\item \emph{Factorization maps}: For each partition $g = g_1 + g_2$,
  a chain map
  \begin{equation}\label{eq:module-factorization}
  \Delta_M: M^{(g)} \to M^{(g_1)} \otimes_{\cA} M^{(g_2)}
  \end{equation}
  compatible with the algebra factorization
  (Theorem~\ref{thm:chain-modular-functor}(ii)).
\item \emph{Self-sewing maps}: A chain map
  \begin{equation}\label{eq:module-self-sewing}
  \sigma_M: M^{(g)} \to M^{(g-1)}
  \end{equation}
  compatible with the algebra self-sewing
  (Theorem~\ref{thm:chain-modular-functor}(iii)).
\item \emph{Coherence}: The factorization and self-sewing maps
  satisfy associativity and commutativity up to homotopies
  provided by the $A_\infty$ module operations.
\end{enumerate}
The curvature $m_0^{(g)} = \kappa(\cA) \cdot \lambda_g$ mixes
the genus components: $d_{M^{(g)}}^2 = m_0^{(g)} \cdot \mathrm{id}$,
so the genus-graded module is a curved complex over the formal
genus parameter $\hbar$.
\end{definition}

\begin{theorem}[Module tower from bar complex with insertions;
\ClaimStatusProvedHere]\label{thm:module-genus-tower}
\index{module tower!genus-graded}
For a Koszul chiral algebra $\cA$ and an $\cA$-module $M$, the
bar complex with module insertion
$\bar{B}^{(g)}_n(\cA; M) :=
\bar{B}^{(g)}_n(\cA; M, \cA, \ldots, \cA)$
\textup{(}with $M$ inserted at the first marked
point\textup{)} forms a genus-graded $\cA$-module in the sense
of Definition~\textup{\ref{def:genus-graded-module}}.
\end{theorem}

\begin{proof}
The bar complex with module insertion
$\bar{B}^{(g)}_n(\cA; M)$ is defined via configuration space
integrals on the logarithmic FM compactification
(Theorem~\ref{thm:genus-induction-strict}).  The nilpotency
$d^2 = 0$ is proved there.

The factorization map~\eqref{eq:module-factorization} is
constructed as follows.  The separating boundary divisor
$D_{\mathrm{sep}} \subset \overline{\mathcal{M}}_{g,n}$ with
the module insertion on the genus-$g_1$ component produces a
residue map
$\mathrm{Res}_{D_{\mathrm{sep}}}\colon
\bar{B}^{(g)}_n(\cA; M) \to
\bar{B}^{(g_1)}_{|I|+1}(\cA; M) \otimes
\bar{B}^{(g_2)}_{|J|+1}(\cA)$.
The module insertion stays on one component; the other component
carries only algebra insertions.  That this is a chain map follows
from the compatibility of the Feynman transform differential with
boundary strata (Getzler--Kapranov \cite{GeK98}).

The self-sewing map~\eqref{eq:module-self-sewing} is constructed
from the nonseparating boundary divisor, analogously to the
algebra case (Theorem~\ref{thm:chain-modular-functor}(iii)).

Coherence follows from the $A_\infty$ module structure on
$\bar{B}^{(g)}_n(\cA; M)$, which provides the required
homotopies between different orderings of degeneration.
\end{proof}

\begin{proposition}[Koszul duality of genus-graded modules;
\ClaimStatusProvedHere]\label{prop:genus-module-koszul}
\index{Koszul duality!genus-graded modules}
For a Koszul pair $(\cA, \cA^!)$ and a Koszul $\cA$-module $M$
with Koszul dual module $M^! = \Phi(M)$
\textup{(}Theorem~\textup{\ref{thm:e1-module-koszul-duality})},
the genus-graded module structures are related by:
\begin{equation}\label{eq:genus-module-duality}
\bar{B}^{(g)}_n(\cA; M) \simeq
\mathrm{RHom}_{\cA^!}\bigl(\bar{B}^{(g)}_n(\cA^!; M^!),\,
\omega_{\overline{\mathcal{M}}_{g,n}}\bigr)
[-\dim_{\mathbb{C}} \mathcal{M}_{g,n}]
\end{equation}
In particular, the factorization maps dualize:
$\Delta_M$ for $\cA$ is Verdier dual to $\Delta_{M^!}$ for $\cA^!$.
\end{proposition}

\begin{proof}
Combine Theorem~\ref{thm:e1-module-koszul-duality} (module Koszul
duality $\Phi: \mathrm{Mod}(\cA) \to \mathrm{Mod}(\cA^!)$)
with Corollary~\ref{cor:dual-modular-functor} (Verdier duality of
the chain-level modular functor).  The module insertion at a marked
point transforms covariantly under the Koszul functor $\Phi$,
while the algebra insertions at other marked points transform via
the bar-cobar equivalence.  The Verdier duality on
$\overline{\mathcal{M}}_{g,n}$ interchanges the factorization
maps by the naturality argument of
Corollary~\ref{cor:dual-modular-functor}.
\end{proof}

\begin{remark}[Curvature as obstruction to strict genus-grading]
\label{rem:curvature-genus-obstruction}
\index{curvature!obstruction to genus-grading}
The genus-graded module structure of
Definition~\ref{def:genus-graded-module} is \emph{not} a direct
sum decomposition when $\kappa(\cA) \neq 0$: the curvature
$m_0^{(g)} = \kappa \cdot \lambda_g$ mixes genus components.
The genus-$g$ module differential satisfies
$d_{M^{(g)}}^2 = m_0^{(g)} \cdot \mathrm{id}$, so the individual
$M^{(g)}$ are curved complexes.  Only the total complex
$M^{\mathrm{total}} = \widehat{\bigoplus}_g \hbar^{2g-2} M^{(g)}$
has $d^2 = 0$ (as a formal power series in $\hbar$).

At critical level $\kappa = 0$, the curvature vanishes and the
genus-grading becomes strict: each $M^{(g)}$ is an honest cochain
complex, and the genus-graded module is a direct sum.  This
recovers the classical (non-genus-graded) module theory as the
$g = 0$ summand.
\end{remark}

\begin{example}[KM Verma modules]\label{ex:verma-genus-graded}
\index{Verma module!genus-graded}
For a Verma module $M(\lambda) \in \mathcal{O}_k(\widehat{\fg})$
at generic level, the genus-graded module tower has:
\begin{itemize}
\item $M^{(0)}(\lambda)$: the bar complex of the Verma module at
  genus~$0$, whose cohomology is the Verma--BGG resolution
  (Theorem~\ref{thm:yangian-bgg}).
\item $M^{(1)}(\lambda)$: the genus-$1$ correction, controlled by
  the elliptic bar differential with module insertion.  The
  obstruction $m_0^{(1)} = \kappa \cdot \lambda_1$ produces the
  genus-$1$ curvature of Theorem~\ref{thm:extension-obstruction}.
\item Factorization: the map
  $\Delta_M: M^{(g)}(\lambda) \to
  M^{(g_1)}(\lambda) \otimes M^{(g_2)}(0)$
  (with vacuum module at the second factor) recovers the genus
  tower restriction maps of Theorem~\ref{thm:master-tower}.
\end{itemize}
\end{example}

\subsection{Fusion product preservation under bar-cobar}
\label{subsec:fusion-bar-cobar}
\index{fusion product!bar-cobar preservation}

\begin{conjecture}[Fusion product preservation;
\ClaimStatusConjectured]\label{conj:fusion-bar-cobar}
\index{fusion product!Koszul duality}
For a Koszul chiral algebra $\cA$ with Koszul dual $\cA^!$,
the module-level bar-cobar functor
$\Phi\colon \mathrm{Mod}(\cA) \to \mathrm{CoMod}(\cA^!)$
\textup{(}Theorem~\textup{\ref{thm:e1-module-koszul-duality})}
preserves fusion products:
\begin{equation}\label{eq:fusion-preservation}
\Phi(M_1 \boxtimes_\cA M_2) \;\simeq\;
\Phi(M_1) \boxtimes_{\cA^!} \Phi(M_2)
\end{equation}
as objects of the comodule category.

For Kac--Moody algebras $\widehat{\fg}_k$ at integrable level:
\begin{enumerate}[label=\textup{(\roman*)}]
\item The fusion product $\boxtimes_k$ on
  $\mathcal{O}_k^{\mathrm{int}}(\widehat{\fg})$ is sent to
  the tensor product $\otimes_q$ on
  $\mathrm{Rep}^{\mathrm{fd}}(U_q(\fg))$, preserving the
  braiding up to a sign determined by the level shift.
\item For Heisenberg Fock modules $F_\lambda \boxtimes_\kappa
  F_\mu \cong F_{\lambda+\mu}$, the braiding phase
  $e^{2\pi i \lambda\mu/\kappa}$ is sent to
  $e^{-2\pi i \lambda\mu/\kappa}$ under $\kappa \mapsto -\kappa$.
\item The bar coalgebra coproduct
  \begin{equation}\label{eq:bar-coproduct-fusion}
  \Delta\colon \barB(\cA, M_1 \boxtimes M_2) \to
  \barB(\cA, M_1) \otimes \barB(\cA, M_2)
  \end{equation}
  is the chain-level realization of the fusion product, induced by
  the separating degeneration map of
  Theorem~\textup{\ref{thm:chain-modular-functor}(ii)}.
\end{enumerate}
\end{conjecture}

\begin{remark}[Scope]\label{rem:fusion-scope}
The conjecture assembles three ingredients from this monograph:
\begin{itemize}
\item The module-level bar-cobar adjunction
  (Theorem~\ref{thm:e1-module-koszul-duality}), which provides
  $\Phi$ as an exact functor.
\item The chain-level modular functor
  (Theorem~\ref{thm:chain-modular-functor}), which provides the
  factorization map $\Delta_{\mathrm{sep}}$ realizing fusion at
  chain level.
\item The Heisenberg fusion computation
  (Proposition~\ref{prop:fock-fusion-product}), which verifies the
  phase inversion in part~(ii).
\end{itemize}
The gap is showing that $\Phi$ is \emph{monoidal}: the bar coalgebra
coproduct~\eqref{eq:bar-coproduct-fusion} must intertwine with the
fusion tensor product.  This is the hardest part of the
Kazhdan--Lusztig programme~\cite{KL93}, where the analogous
statement requires the analysis of $\mathcal{D}$-modules on
configuration spaces of the flag variety.  A configuration space
proof via the bar complex would need to show that the residue
maps along separating divisors of $\overline{\mathcal{M}}_{g,n}$
are compatible with the fusion product's associativity and commutativity
constraints.
\end{remark}
